\chapter{Ancient solutions and singularity models}
\label{ch:ancient-models}

Ancient solutions describe geometry on every backward time scale.
Noncollapsing prevents a bounded-curvature region from losing dimension
under rescaling.  Together with nonnegative curvature, it allows the
extraction of shrinking models at backward infinity and compactness at
points of normalized scalar curvature.  These are different limits:
a backward rescaling has a soliton equation, whereas a limit of
normalized ancient solutions need not be a soliton.

Throughout, $N$ is a connected smooth $n$-manifold without boundary,
Hausdorff, regular, and second countable, with its Borel structure.
It need not be the original closed three-manifold.  Write
$A=|\operatorname{Rm}|$ for the full curvature norm.  We use
\[
 {\cal R}(X,Y)Z=\nabla_X\nabla_YZ-\nabla_Y\nabla_XZ-\nabla_{[X,Y]}Z,
 \qquad \operatorname{Rm}(u,v,w,z)=g({\cal R}(u,v)z,w),
\]
and $\operatorname{Ric}(u,v)=\sum_i\operatorname{Rm}(u,e_i,v,e_i)$.
The measure $\mu_g=c_n{\cal H}^n_{d_g}$ is calibrated to equal
Riemannian volume.  The Euclidean unit ball has volume $\omega_n$.
All balls below are open and have positive radius.

\section{Ancient noncollapsed flows}

\begin{definition}[Ancient $\kappa$-solution]
\label{def:ancient-kappa}
An ancient $\kappa$-solution is a smooth Ricci flow $g(t)$ on $N$,
$-\infty<t\leq0$, with $\kappa>0$, such that every slice is complete,
has nonnegative curvature operator, has bounded $A(t,\cdot)$, and is
nonflat.  The bound may initially depend on the slice.  For every
$t\leq0$, $p\in N$, and $r>0$,
\[
 \bigl(A(s,q)\leq r^{-2}\text{ for }t-r^2<s\leq t,\
          q\in B_t(p,r)\bigr)
 \ \Longrightarrow\ \mu_{g(t)}(B_t(p,r))\geq\kappa r^n .
\]
The spatial ball is measured at the terminal time throughout the cylinder.
The flow equation at zero uses derivatives within the closed time domain.
\lean{PoincareMT.AncientKappaSolution}
\source{MT2007}
\dcref{Definitions 9.1--9.2, pp. 179--180}
\end{definition}

Nonnegative curvature operator means positivity of the curvature
quadratic form on skew two-tensors; it implies nonnegative sectional
and Ricci curvature.  Nonflatness forces $n\geq2$, since the
curvature tensor vanishes in dimensions zero and one.
A static metric is \emph{metrically}
$\kappa$-noncollapsed when $A\leq r^{-2}$ on $B(p,r)$ implies
$\mu(B(p,r))\geq\kappa r^n$.  This is distinct from the parabolic condition.

\begin{theorem}[Structural estimates and normalization]
\label{thm:ancient-structure}
For an ancient $\kappa$-solution, at every point and $t\leq0$,
\[
 0<R(t,x),\quad A(t,x)\leq R(t,x)\leq nA(t,x),\quad
 0\leq\operatorname{Ric}(v,v)\leq Rg(v,v).
\]
For $s\leq t\leq0$,
\[
 R(s,x)\leq R(t,x),\quad g(t)\leq g(s),\quad d_t\leq d_s,\quad
 B_s(p,r)\subseteq B_t(p,r).
\]
For each $b\leq0$, $A$ is uniformly bounded on $N\times(-\infty,b]$;
indeed $A(t,x)\leq R(b,x)$ for $t\leq b$.
For $p\in N$ and $Q=R(b,p)>0$, the same-carrier flow
$\widehat g(s)=Qg(b+s/Q)$, $s\leq0$, is an ancient solution with the
same $\kappa$ and $\widehat R(0,p)=1$.
\lean{PoincareMT.ancientKappaStructuralConsequences}
\uses{def:ancient-kappa}
\source{MT2007}
\dcref{Theorem 4.40, p. 82; Definition 3.40, p. 61; Section 9.1}
\end{theorem}
\begin{proof}
The curvature inequalities are algebraic consequences of positivity of
the curvature operator, with the stated norm convention.
The ancient trace Harnack estimate is
\[
 \partial_tR+2\,dR(v)+2\operatorname{Ric}(v,v)\geq0 .
\]
Taking $v=0$ gives scalar monotonicity.  Its integrated path form
propagates positivity from any point at an earlier time to every point
at a later time.  Nonflatness and $A\leq R$ provide a positive point
on every earlier slice.  Thus $R>0$ everywhere.
Integrating $\partial_tg=-2\operatorname{Ric}\leq0$ against a fixed
tangent vector gives metric monotonicity.  Integrating along paths and
taking infima of lengths gives the distance and ball assertions.
If $A(b,\cdot)\leq C_b$, then
$A(t,x)\leq R(t,x)\leq R(b,x)\leq nC_b$ for all $t\leq b$.

Multiplication by $Q>0$ multiplies lengths, volumes, and curvature
norms by $\sqrt Q$, $Q^{n/2}$, and $Q^{-1}$, respectively; covariant
Ricci is unchanged.  The time change preserves the flow equation.
A radius-$r$ cylinder corresponds to radius $r/\sqrt Q$ in the
original flow, so its volume bound becomes exactly $\kappa r^n$.
Completeness, curvature sign, and nonflatness are preserved as well.
\end{proof}

\section{Prescribed rescalings and a single limit carrier}

Fix $o\in N$.  For $\tau>0$ define the reduced length and volume by
\[
 \ell_o(q,\tau)=\frac1{2\sqrt\tau}
 \inf_{\gamma(0)=o,\ \gamma(\tau)=q}
 \int_0^\tau\sqrt s\,
       \bigl(R(-s,\gamma(s))+|\dot\gamma(s)|_{g(-s)}^2\bigr)\,ds ,
\]
\[
 V_{\rm red}(\tau)=
 \int_N\tau^{-n/2}e^{-\ell_o(q,\tau)}\,d\mu_{g(-\tau)}(q).
\]
The admissible paths and endpoint action are as defined in reduced
geometry.  No factor $(4\pi)^{-n/2}$ is inserted; the Euclidean mass
is $(4\pi)^{n/2}$.

\begin{proposition}[A prescribed sequence]
\label{prop:ancient-sequence}
For every reference point $o$ and every prescribed sequence $\tau_i>0$
with $\tau_i\to\infty$, there are global minimizers $p_i$ of
$\ell_o(\cdot,\tau_i)$ with $\ell_o(p_i,\tau_i)\leq n/2$.
On the original carrier set
\[
 g_i(t)=\tau_i^{-1}g(\tau_i t),\quad
 R_i(t,x)=\tau_i R(\tau_i t,x),\quad
 A_i(t,x)=\tau_i A(\tau_i t,x),\qquad t<0.
\]
The construction retains the exact reference and scale sequence.
Monotonicity of the sequence is not required.
\lean{PoincareMT.ancientBlowupSequenceSetup}
\source{MT2007}
\dcref{Corollary 6.74, p. 142; Theorem 7.10 and Claim 7.11,
pp. 154--156; Theorem 9.11, pp. 183--185}
\end{proposition}
\begin{proof}
Finite backward intervals have uniform curvature bounds by
Theorem~\ref{thm:ancient-structure}.  Reduced-length coercivity gives
a quadratic distance lower bound with constants depending on the
interval.  Thus the continuous reduced length is proper on the complete
slice and attains its minimum.  The reduced-geometry minimum estimate
bounds that minimum by $n/2$.  Apply it on the interval of length
$\tau_i+1$, so that $\tau_i$ is an interior time, and choose a minimizer.
Parabolic rescaling gives the displayed formulas on the same carrier.
\end{proof}

\begin{lemma}[Local regularity of a weak parabolic solution]
\label{lem:ancient-weak-regularity}
Let $\Omega\subset\mathbb R^n\times\mathbb R$ be open, $n\geq1$.
Suppose that $a^{ij},b^i,c$ are smooth, that $a$ is symmetric,
and that $a(z)$ is positive definite at every $z\in\Omega$.
If a continuous function $w$ solves
\[
 Pw=\partial_t w-a^{ij}\partial_i\partial_jw
                  +b^i\partial_iw+cw=0
\]
in distributions, then $w$ is smooth.  Repeated spatial indices
are summed.  The distributional assertion means that
\[
 \int_\Omega w\bigl(-\partial_t\phi
       -\partial_i\partial_j(a^{ij}\phi)
       -\partial_i(b^i\phi)+c\phi\bigr)=0
 \quad(\phi\in C_c^\infty(\Omega)).
\]
No uniform ellipticity constant on all of $\Omega$ is required.
\end{lemma}
\begin{proof}
All norms and integrals in this proof use Euclidean spacetime
measure.  Around each point choose several nested balls with
compact closure in $\Omega$.  On the largest ball there is
$\kappa>0$ with $a\geq\kappa I$, and every needed coefficient
derivative is bounded.  Multiply coefficients and $w$ by a
compactly supported cutoff equal to one on this ball when
forming convolutions.  The extended coefficients need not be
elliptic outside the ball.  We use the equation only where
every translated convolution kernel stays inside the region
on which these extensions agree with the original functions.

\emph{The commutator and the first energy estimate.}
Let $\eta$ be a nonnegative smooth function of integral one
supported in the spacetime unit ball, and put
$\eta_\epsilon(z)=\epsilon^{-(n+1)}\eta(z/\epsilon)$.
For a smooth scalar coefficient $q$ and a locally integrable
function $v$, write
\begin{align*}
 C_{\epsilon,j}(q,v)
   &=q\,\partial_j(\eta_\epsilon*v)
        -(\partial_j\eta_\epsilon)*(qv)
        +\eta_\epsilon*((\partial_jq)v)\\
   &=\int\bigl[(q(z)-q(y))\partial_j\eta_\epsilon(z-y)
          +\eta_\epsilon(z-y)\partial_jq(y)\bigr]v(y)\,dy.
\end{align*}
On the compact region, the absolute value is bounded by the
convolution of $|v|$ with
\[
 L\bigl(|s|\,|\partial_j\eta_\epsilon(s)|
                                      +|\eta_\epsilon(s)|\bigr),
\]
where $L$ bounds the full spacetime derivative of $q$.
The $L^1$ norm of this kernel is independent of $\epsilon$:
the factor $|s|$ cancels the extra inverse scale from the
derivative.  Thus the commutator is uniformly bounded in
$L^\infty$ for bounded $v$, and in $L^2$ for $v\in L^2$
by the convolution inequality $\|k*v\|_2\leq\|k\|_1\|v\|_2$.

Set $d^i=b^i+\partial_ja^{ji}$ and $H=\partial_id^i-c$.
The original weak equation can be written
\[
 \partial_t w-\operatorname{div}(aDw)
                           =Hw-\operatorname{div}(dw).
\]
Here $Dw$ is initially a distribution.  Its mollification
$w_\epsilon=\eta_\epsilon*w$ satisfies
\[
 \partial_t w_\epsilon-\operatorname{div}(aDw_\epsilon)
        =g_\epsilon+\operatorname{div}F_\epsilon,\qquad
 g_\epsilon=\eta_\epsilon*(Hw),\qquad
 F_\epsilon^i=-C_{\epsilon,j}(a^{ij},w)
                              -\eta_\epsilon*(d^iw).
\]
Since $w$ is continuous, $w_\epsilon,g_\epsilon,F_\epsilon$
are bounded uniformly on the fixed inner ball.  Test this
smooth equation against $\chi^2w_\epsilon$, with a compact
spacetime cutoff $\chi$.  Integration by parts gives
\begin{align*}
 \int\chi^2\langle aDw_\epsilon,Dw_\epsilon\rangle
   ={}&\int\chi^2w_\epsilon g_\epsilon
          +\int\chi\chi_t w_\epsilon^2
          -2\int\chi w_\epsilon\langle aDw_\epsilon,D\chi\rangle\\
      &-\int\chi^2\langle Dw_\epsilon,F_\epsilon\rangle
          -2\int\chi w_\epsilon\langle D\chi,F_\epsilon\rangle.
\end{align*}
Young's inequality absorbs the two terms containing
$\chi Dw_\epsilon$ into half the elliptic left side.
Everything else is uniformly bounded on the compact support.
Consequently $\int\chi^2|Dw_\epsilon|^2$ is bounded.
Weak compactness in $L^2$, uniform convergence of
$w_\epsilon$ to $w$, and integration by parts identify the
weak limits as the spatial derivatives of $w$.
Thus $Dw\in L^2_{\rm loc}$.

\emph{The second energy estimate.}
Freeze $A=a(z_0)$ and shrink the ball until
$|a-A|_{\rm F}\leq\kappa/2$, where the subscript denotes
the matrix Frobenius norm.  For a smooth compactly supported
function $v$, integration by parts gives
\[
 \int v_t(A:D^2v)=0,\qquad
 \int(A:D^2v)\Delta v
       =\sum_k\int\langle A D\partial_kv,D\partial_kv\rangle
       \geq\kappa\int|D^2v|^2,
 \qquad \|\Delta v\|_2^2=\|D^2v\|_2^2.
\]
Cauchy--Schwarz in the middle identity implies
$\|A:D^2v\|_2^2\geq\kappa^2\|D^2v\|_2^2$.
Expanding the first square, then absorbing the variation
of $a$, yields
\begin{align*}
 \|v_t\|_2^2+\kappa^2\|D^2v\|_2^2
      &\leq\|v_t-A:D^2v\|_2^2,\\
 \|v_t\|_2^2+\frac{\kappa^2}{2}\|D^2v\|_2^2
      &\leq2\|v_t-a:D^2v\|_2^2.
\end{align*}
Apply the latter estimate to $\chi v$, with $\chi=1$ on
a smaller ball.  Its principal residual is
\[
 (\partial_t-a:D^2)(\chi v)
   =\chi(v_t-a:D^2v)+(\chi_t-a:D^2\chi)v
                              -2a^{ij}\partial_i\chi\,\partial_jv.
\]
It follows that the $L^2$ norms of $v_t$ and $D^2v$ on the
smaller ball are bounded by a constant times the sum of
the $L^2$ norms of $v_t-a:D^2v$, $v$, and $Dv$ on the
larger ball.  The cutoffs and the constant are fixed before
the mollifier scale tends to zero.

We need the forced version of this assertion.  Suppose
$Pw=f$ weakly, with $w,Dw,f\in L^2_{\rm loc}$, and write
$g_j=\partial_jw$ for its weak derivatives.  Testing the
weak derivative identities against translated kernels gives
$\partial_jw_\epsilon=\eta_\epsilon*g_j$.  The mollified
principal residual is exactly
\[
 \partial_t w_\epsilon-a^{ij}\partial_i\partial_jw_\epsilon
   =\eta_\epsilon*f-\eta_\epsilon*(b^ig_i)
       -\eta_\epsilon*(cw)-C_{\epsilon,i}(a^{ij},g_j).
\]
The commutator estimate now bounds this expression in $L^2$
uniformly in $\epsilon$.  Convolution also bounds
$w_\epsilon$ and $Dw_\epsilon$ in $L^2$.  The second
energy estimate therefore bounds $\partial_t w_\epsilon$
and $D^2w_\epsilon$ on the smaller ball.  Weak subsequential
limits, paired against compact tests, give the actual
distributional derivatives $w_t,D^2w\in L^2_{\rm loc}$.
This argument applies to weak solutions with these $L^2$
assumptions; it does not require their continuity.

\emph{All derivatives and the continuous representative.}
Apply the forced estimate first to the original equation
with $f=0$.  Differentiate it in a spatial direction:
\[
 P(\partial_jw)
    =\partial_jf+(\partial_ja^{ik})\partial_i\partial_kw
       -(\partial_jb^i)\partial_iw-(\partial_jc)w.
\]
This identity holds in distributions by the weak product
rule and commutation of constant derivatives on test
functions.  For the homogeneous original equation,
the right side is now in $L^2$, and $\partial_jw$ already
has its first spatial derivatives in $L^2$.  The forced
estimate gives one more spatial order.  More generally,
the equation for $D^\alpha w$ has forcing consisting of
smooth coefficient derivatives times derivatives of $w$
of order at most $|\alpha|+1$.  Induction, shrinking the
ball finitely many times for each prescribed order, gives
all spatial weak derivatives locally in $L^2$.

The equation itself says
$w_t=a^{ij}\partial_i\partial_jw-b^i\partial_iw-cw$.
For each $m$, repeated distributional differentiation and
this identity express $\partial_t^mw$ as a spatial
differential operator of order at most $2m$ applied to
$w$, with smooth coefficients.  To pass from $m$ to $m+1$,
differentiate those coefficients in time and substitute
the displayed expression for $w_t$; constant spatial
derivatives commute with the time derivative on tests.
Thus every mixed spacetime weak derivative is locally
in $L^2$ as well.

For clarity, the final Sobolev step is local in ordinary
spacetime dimension $d=n+1$.  To obtain a $C^r$
representative, take a cutoff of $w$ on a ball where
weak derivatives through order $r+1+d$ are in $L^2$.
Finite support volume puts them in $L^{p_0}$ for
$1<p_0<2$.  Choose $p_0$ avoiding the finitely many values
that make the following iteration hit $d$ exactly.
The Sobolev inequality $W^{1,p}\subset L^{dp/(d-p)}$
for $1<p<d$, applied to each remaining weak derivative,
trades one derivative for the reciprocal-exponent change
$1/p\mapsto1/p-1/d$.  In at most $d$ steps it gives
$W^{r+1,q}$ for some $q>d$.  Morrey's inequality gives
a $C^r$ representative on a smaller ball.  The continuous
original function agrees with this representative:
their difference vanishes almost everywhere, and a
nonzero continuous value would persist on an open set
of positive measure.  For every finite $r$ this proves
local $C^r$ regularity of that same function.  The
neighborhood may depend on $r$; smoothness follows
because being $C^r$ is local for each $r$.
\end{proof}

Only local ellipticity is needed here: on each compact ball,
pointwise positive definiteness of the smooth principal matrix
gives a positive minimum eigenvalue.  The preceding energy and
regularity argument applies on that ball before it is shrunk.
This supplies the regularity conclusion used in
\cite[Proposition 9.20, p.~198]{MT2007}.

\begin{theorem}[Asymptotic soliton equation]
\label{thm:ancient-asymptotic}
Every prescribed sequence above has a strictly increasing subsequence
map $\sigma$ and a pointed complete ancient limit
$(N_\infty,p_\infty,h(t))$, $t<0$, with nonnegative curvature operator.
There are nested connected relatively compact open sets $U_j$
containing $p_\infty$, exhausting $N_\infty$, and windows
$I_j=[-(j+1),-(j+1)^{-1}]$ exhausting negative time.
On compact subregions the rescaled flows converge by smooth embeddings
preserving time and spatial worldlines, with the basepoint mapped to
$p_{\sigma(i)}$ at time $-1$.  All spacetime metric jets converge
uniformly on compact coordinate sets.
The limit retains $\kappa$, is nonflat on every slice, has
$\partial_tR_h\geq0$, and admits a smooth function $f$ satisfying
\[
 \operatorname{Ric}_{h(t)}+\operatorname{Hess}_{h(t)}f(\cdot,t)
                          +\frac1{2t}h(t)=0.              \tag{A}
\]
At this stage neither global slice curvature bounds nor global
self-similarity have been asserted.
\lean{PoincareMT.ancientAsymptoticSolitonLimits}
\uses{prop:ancient-sequence}
\source{MT2007}
\dcref{Theorem 9.11 and Remark 9.12, pp. 183--185;
Claims 9.15--9.16 and Corollary 9.17, pp. 186--187;
Claim 9.22, pp. 190--191; Corollary 9.26 through Claim 9.31,
pp. 195--198; Proposition 9.32 and Lemma 9.33, pp. 198--202}
\end{theorem}
\begin{proof}[Proof: geometric extraction]
\emph{Uniform curvature with two independent times.}
Write $\ell_i(x,\tau)=\ell_o(x,\tau_i\tau)$.
The ancient differential Harnack inequality, applied along a
minimizing reduced geodesic, bounds its reduced Harnack density
below by $-R/s$.  After multiplication by $s^{3/2}$ and
integration, its Harnack integral is at least minus the
reduced action $\mathcal L$.  Since
$\mathcal L=2\sqrt\tau\,\ell$, the regular-point variation
formulas of Chapter~\ref{ch:reduced-geometry} give
\[
 |\nabla\ell|^2+R\leq\frac{3\ell}{\tau},\qquad
 \partial_\tau\ell\geq R-\frac{2\ell}{\tau}
                                  \geq-\frac{2\ell}{\tau}.
\]
The first scalar consequence extends across the cut locus
by continuity and density of regular points.  For the
second inequality, local spacetime Lipschitz regularity
gives absolute continuity on compact worldline intervals.
Fubini gives the regular-point inequality at almost every
time for almost every spatial point.  Integrating
$\partial_\tau(\tau^2\ell)\geq0$ on those worldlines
and using spatial continuity extends the result to every
point, including worldlines lying in the cut locus.
The basepoint choice therefore gives, for $0<\sigma\leq1$,
\[
 \ell_i(p_i,\sigma)\leq\frac n{2\sigma^2},\qquad
 R_i(p_i,-\sigma)\leq\frac{3n}{2\sigma^3}.
\]

For $a<b<0$, join $x$ to $y$ by a minimizing geodesic
of $g_i(a)$, parametrized with constant speed on $[a,b]$.
Nonnegative Ricci curvature makes forward metrics decrease,
so this path has spacetime energy at most
$d_{g_i(a)}(x,y)^2/(b-a)$.  Integrated ancient Harnack yields
\[
 R_i(x,a)\leq R_i(y,b)
              \exp\!\left(\frac{d_{g_i(a)}(x,y)^2}{2(b-a)}\right).
\]
Now fix $-2\leq a<0$, $r>0$, and let $s,t\leq a$.
For $x\in B_{g_i(s)}(p_i,r)$ set $q=\max\{s,t\}$.
Metric monotonicity gives $d_{g_i(q)}(p_i,x)<r$, while
the structural curvature and scalar monotonicity estimates give
$|\operatorname{Rm}_{g_i(t)}|(x)\leq R_i(x,q)$.
Take $\sigma=-a/2$, so $a<-\sigma<0$ and $0<\sigma\leq1$.
The two preceding bounds imply
\[
 |\operatorname{Rm}_{g_i(t)}|(x)
 \leq \frac{3n}{2\sigma^3}
          \exp\!\left(\frac{r^2}{2(-\sigma-a)}\right).
\]
This constant is independent of $i$, of the ball's reference
time $s$, and of the curvature-evaluation time $t$.

\emph{Complete limits on buffered windows.}
Use the open windows
$J_j=(-(j+2),-(j+2)^{-1})$, which contain $I_j$ in their
interiors and contain the based time $-1$.
The two-time estimate supplies the compactness curvature
hypotheses on every fixed-radius ball throughout each
closed buffered window.  Completeness makes the closed
based balls compact.  If $C\geq0$ bounds curvature on
the based unit ball, choose $r_0=(C+1)^{-1}\leq1$.
Then $C\leq r_0^{-2}$ and the parabolic interval
$(-1-r_0^2,-1]$ lies in $[-(j+2),-(j+2)^{-1}]$.
The original noncollapse hypothesis gives
\[
 \mu_{g_i(-1)}(B_{g_i(-1)}(p_i,r_0))\geq\kappa r_0^n.
\]
In uncalibrated Hausdorff measure this is the same bound
with constant $\kappa/c_n$.  This checks the positive
based-volume input to Theorem~\ref{thm:rf-compactness},
in addition to its two-time curvature inputs.  After
shifting the based time to zero, that theorem produces
a pointed smooth subsequential flow on each $J_j$,
complete on every interior slice.  The singleton
$I_0=\{-1\}$ is thus always inside a genuine open interval.

\emph{Identifying two complete finite-window limits.}
Let $G$ and $H$ be limits on successive windows, with
common based time $-1$, and let $e_i,\widetilde e_i$ be
their spatial embeddings into the same members of the
source sequence.  Their domains exhaust the two complete
manifolds.  On each fixed compact set, the pullback tangent
norms have distortion between $D^{-1}$ and $D$ for any
fixed $D>1$, once $i$ is large enough.

A minimizing radial path gives
$e_i(B_G(p,r))\subset B_i(p_i,Dr)$.
Conversely, a source path of length less than $r$ can
be lifted by the local inverse starting at $p$.
As long as it is lifted, its length is at most $D r$.
Choose a larger limit ball with compact closure inside
the embedding domain.  The lift cannot first leave
that ball, and compactness and a local inverse continue
it to its endpoint.  Hence
$B_i(p_i,r)\subset e_i(B_G(p,Dr))$ eventually.
The same statements hold for $\widetilde e_i$.
With $D=2$ they show that
\[
 F_i=\widetilde e_i^{-1}\circ e_i
       :B_G(p,r)\longrightarrow B_H(\widetilde p,4r)
\]
is defined for all large $i$.  Completeness makes the
closed target ball compact.

On a larger ball the exact pullback identity and the
two norm comparisons bound $|DF_i|$ by $D^2$.
A minimizing path between two points of $B_G(p,r)$
stays in $B_G(p,3r)$, so this differential bound gives
$d_H(F_i x,F_i y)\leq D^2d_G(x,y)$ there.
The reverse comparisons give the other inequality,
using their exact cancellation on the smaller domains.
Arzela--Ascoli on a countable exhaustion and a diagonal
choice give a continuous compact-uniform limit $F$.
Sending $i\to\infty$ and then $D\downarrow1$ in the
two inequalities proves $d_H(Fx,Fy)=d_G(x,y)$.
For any $y\in H$, its reverse images
$x_i=e_i^{-1}\widetilde e_i(y)$ remain in one compact
ball of $G$.  A subsequence converges to $x$;
compact-uniform convergence and $F_i(x_i)=y$ give
$F(x)=y$.  Thus $F$ is a bijective distance isometry.
Its inverse comparisons converge uniformly on compact
sets: for $x_i=F_i^{-1}y$ in a fixed compact ball,
\[
 d_G(F^{-1}y,x_i)=d_H(y,Fx_i)
                   \leq\sup_x d_H(F_i x,Fx)\longrightarrow0.
\]

We also need smooth convergence of these same comparisons
and preservation of every metric in the common time
interval.  In fixed source and target charts let $A_i$
and $B_i$ denote the pulled-back metrics at the based
time.  On a compact chart neighborhood and a slightly
larger compact target neighborhood they converge with
all derivatives and have a common positive ellipticity
bound.  The identity $F_i^*B_i=A_i$ bounds the first
derivatives.  Preservation of their Levi-Civita connections
gives the coordinate recurrence
\[
 \partial_\alpha\partial_\beta F_i^k
   =\Gamma(A_i)^\gamma_{\alpha\beta}\partial_\gamma F_i^k
    -\Gamma(B_i)^k_{ab}(F_i)
                      \partial_\alpha F_i^a\partial_\beta F_i^b.
\]
Inverse matrices and their derivatives are bounded by
ellipticity and differentiation of $A_iA_i^{-1}=I$,
and similarly for $B_i$.  Differentiating the recurrence
bounds each next derivative of $F_i$ in terms of already
bounded derivatives.  Compact target confinement bounds
order zero.  Arzela--Ascoli for these derivatives on
smaller charts, with integration along coordinate
segments to identify consecutive derivatives, makes
every subsequence admit a smooth further limit.
Its order-zero limit must be $F$.  If a derivative of
the original sequence failed to converge locally
uniformly to the corresponding derivative of $F$,
a subsequence witnessing that failure would have just
such a smooth further limit, a contradiction.
Thus the same comparisons converge smoothly to $F$;
the reverse argument proves smoothness of $F^{-1}$.

The embeddings preserve spatial worldlines, so their
spatial maps are independent of time.  At every common
time their comparison has the exact metric pullback
identity.  Its smooth limit preserves that identity:
coefficient convergence is uniform on the compact
target neighborhood, including the moving points
$F_i(x)$, and first derivatives converge as just proved.
Consequently $F^*h_H(t)=h_G(t)$ for every common $t$.
The based points agree by the pointed embeddings
(equivalently, by confinement to arbitrarily small
based balls).

\emph{One manifold and one source subsequence.}
Inductively extract the next finite-window limit from
the preceding source subsequence.  Let $F_j:N_j\to N_{j+1}$
be the preceding pointed diffeomorphism, preserving
metrics on $J_j$, and put
$\Phi_0=\operatorname{id}$, $\Phi_{j+1}=F_j\circ\Phi_j$.
Every $\Phi_j$ is a global diffeomorphism from $N_0$
to $N_j$.  The metrics $\Phi_j^*h_j(t)$ agree exactly
on overlapping windows.  Define $h(t)$ on $N_0$ using
any window containing $t$.  Near each time it equals
one smooth pulled-back flow, so it is smooth and solves
the Ricci equation.  Each slice is complete because
it is isometric to the complete corresponding slice
of that finite-window limit.  No completeness assertion
is inferred merely from an exhaustion.

Choose a connected relatively compact open exhaustion
$U_j$ of $N_0$ from the first limit.  In window $j$
choose an index $k_j$ so large that $\Phi_j(\overline U_j)$
lies in its embedding domain.  Increase $k_j$ further
to make the first $j$ tests in a countable list of compact
chart sets and derivative orders have metric error less
than $1/(j+1)$, whenever their time sets lie in $J_j$.
This is possible because the metric in every test is
the same glued metric on its window.  Choose $k_j$
strictly increasing.  If $s_j$ is the source-subsequence
map at window $j$, then
$s_{j+1}=s_j\circ r_j$ with $r_j$ strictly increasing,
so $s_{j+1}(k)\geq s_j(k)$.  Hence
$\sigma(j)=s_j(k_j)$ is strictly increasing.
Composing that window's spatial embedding with $\Phi_j$
gives the required map on $U_j\times I_j$.
Every fixed compact chart test eventually occurs in
the finite list and inside $J_j$.  This proves convergence
of all spacetime metric derivatives on compact sets.

\emph{Preserving curvature and noncollapse.}
Curvature sign and the scalar time-derivative inequality
pass through smooth metric convergence.  To retain
the exact $\kappa$, fix $t<0$, $p$, and $r>0$ with
\[
 |\operatorname{Rm}_{h(s)}|(x)\leq r^{-2}
 \quad\text{for }s\in(t-r^2,t],\quad x\in B_{h(t)}(p,r).
\]
For $0<\rho<a<r$, the compact cylinder
$[t-\rho^2,t]\times\overline{B_{h(t)}(p,a)}$
lies strictly inside this region.  Its curvature bound
has the strict margin $r^{-2}<\rho^{-2}$.
Smooth convergence transfers that latter bound to its
source images.  The ball-containment argument above,
now centered at $p$, puts the source $\rho$-ball inside
the image of $B_{h(t)}(p,a)$, using distortion $a/\rho>1$.
Time independence of the spatial embedding makes the
same preimage valid at every $s\in[t-\rho^2,t]$.
Source noncollapse gives volume at least
$\kappa\rho^n$ for that source ball.

These source ball volumes converge to the limit ball
volume.  Indeed, the two ball containments and Jacobian
bounds give, for every $D>1$, the eventual inequalities
\[
 D^{-n}\mu_{h(t)}(B(p,\rho/D))
 \leq \mu_{g_i(t)}(B(e_i p,\rho))
 \leq D^n\mu_{h(t)}(B(p,D\rho)).
\]
The radius-volume function is continuous at positive
radii: a metric sphere has zero volume, since minimizing
geodesics cover it by the exponential image of the
tangent sphere, whose smooth image has zero
$n$-dimensional volume.  Compactness of finite balls
and continuity of measure then give the assertion.
Sending $D\downarrow1$ proves convergence.  It follows
that $\mu_{h(t)}(B(p,r))\geq\kappa\rho^n$ for every
$0<\rho<r$.  Letting $\rho\uparrow r$ proves the
original noncollapse constant, with no loss.
\end{proof}
\begin{proof}[Proof: the potential and equality]
Pull back $\ell_i(q,\tau)=\ell_o(q,\tau_{\sigma(i)}\tau)$ with these
same embeddings.  We first extract their limit while retaining the
geometric convergence just constructed.

\emph{Bounds through the cut locus.}
At a regular point the first-variation identities give
$2\ell_\tau=R-|\nabla\ell|^2-\ell/\tau$.
Together with the bounds already proved, this implies
$|\ell_\tau|\leq2\ell/\tau$.
The absolute-continuity and full-measure worldline argument
above therefore gives both monotonicity of $\tau^2\ell$
and reverse monotonicity of $\ell/\tau^2$, on all worldlines.
In particular, at every normalized time $\tau>0$,
\[
 \ell_i(p_i,\tau)\leq\frac n2\max\{\tau^2,\tau^{-2}\}.
\]
These bounds are invariant under the chosen parabolic rescaling.

Fix a coordinate ball with compact closure on the limit,
and compose its chart parametrization with the retained
spatial embeddings.  On $\alpha\leq\tau\leq\beta$,
$0<\alpha<\beta$, their tangent norms are bounded by
$B|v|$, with $B$ independent of large $i$.
Indeed, smooth metric convergence gives the bound at
time $-\beta$, and forward metric monotonicity preserves
it through time $-\alpha$.
Since $R>0$ and $R\leq3\ell/\tau$, the source lengths
are strictly positive.  Thus their square roots are
locally Lipschitz.  On the full-measure regular locus,
the gradient estimate gives the coordinate derivative bound
\[
 \bigl|D\sqrt{\ell_i(\cdot,\tau)}\bigr|
                        \leq\tfrac12\sqrt{3/\tau}\,B.
\]
This bounds differences at every pair of points of a
smaller convex coordinate ball.  To justify the assertion
when a segment lies in the exceptional set, translate
the whole segment by a small vector.  Fubini and
translation invariance give the derivative bound almost
everywhere along almost every translated segment.
Local Lipschitz regularity gives absolute continuity
there, so integration proves the stated difference
bound.  The original segment has positive distance
from the boundary of the larger coordinate domain.
Translations tending to zero stay in that domain,
and continuity gives the same bound on the original
segment.  Parabolic scaling cancels in the constant:
the source tangent norm contributes $\sqrt{\tau_i}$,
while the square-root gradient bound contributes
$\tau_i^{-1/2}$.

\emph{From one base value to compact bounds.}
At fixed $\tau$, let $E$ consist of points at which
the sequence $\sqrt{\ell_i}$ is eventually bounded.
The preceding coordinate estimate gives an eventual
uniform oscillation bound on a neighborhood of every
point.  A bound at any one point in such a neighborhood
therefore gives bounds at all its points.  Consequently
both $E$ and its complement are open: a neighborhood
meeting $E$ must be contained in $E$.
The based value estimate makes $E$ nonempty; connectedness
of the limit makes it the whole manifold.  Each point
now has a neighborhood with an eventual uniform bound.
A finite subcover of any compact set, and the maximum
of the finitely many thresholds and constants, give
one eventual bound on that compact set.  Squaring gives
a bound for $\ell_i$ itself.  This argument needs no
single chart containing both the basepoint and that set.

Let $A\geq0$ bound the lengths at time $\beta$ on the
larger coordinate ball, and put $M=A\beta^2/\alpha^2$.
The weighted time inequality gives $0\leq\ell_i\leq M$
throughout its cylinder.  Integration of the absolute
time-derivative bound, first on almost every worldline
and then on all by spatial continuity, gives
\[
 |\ell_i(y,s)-\ell_i(y,t)|
             \leq\frac{2A\beta^2}{\alpha^3}|s-t|.
\]
Factoring the difference of two squares in the spatial
square-root estimate gives, for points in the smaller ball,
\[
 |\ell_i(x,s)-\ell_i(y,t)|
 \leq 2\sqrt M\bigl(\tfrac12\sqrt{3/\alpha}\,B\bigr)|x-y|
              +\frac{2A\beta^2}{\alpha^3}|s-t|.
\]
Thus values and spacetime Lipschitz constants are bounded
on every compact coordinate cylinder, independently of
large $i$.  These are bounds on the actual pulled-back
functions and on the same fixed spatial embeddings.

\emph{A common further extraction.}
Choose countably many compact coordinate cylinders whose
interiors cover $N_\infty\times(0,\infty)$, each with a
slightly larger coordinate ball still in its chart and
positive lower time.  Second countability supplies such
a cover.  On each cylinder discard its finite exceptional
initial segment and apply Arzela--Ascoli to the bounded
Lipschitz functions.  Nested extractions and a strictly
increasing diagonal give uniform convergence on every
cylinder.  Limits agree on overlaps, since both are
pointwise limits of the same source functions there.
They define a continuous nonnegative function
$\ell_\infty$ on all positive backward times.
A compact spacetime set is covered by finitely many
cylinder interiors, so convergence is uniform on it.
Passing the Lipschitz inequality to the limit preserves
local spacetime Lipschitz regularity.

Reselect the existing geometric subsequence by this
further strictly increasing map.  Composition preserves
strict increase and every compact metric-jet limit.
For its $j$th member one may restrict the old embedding
to $U_j\times I_j$, since the new index is at least $j$;
the nested domains and windows contain that product.
Basepoint preservation and time independence persist.
We keep the notation $\sigma$ for the resulting source
subsequence.  Nothing identifies $\ell_\infty$ as a
reduced length based at a point of the limiting ancient
flow; it is the limit of the specified source lengths.

\emph{Strong convergence of the spatial gradients.}
Fix a coordinate ball $B_{2r}(a)$ with compact closure in a limit
chart, and $0<\alpha<\beta<\infty$.  Use the retained embeddings
to express the pulled-back reduced lengths as $w_i(x,\tau)$ and
their limit as $w(x,\tau)$ on this chart cylinder.  Write
\[
 \rho_i=\sqrt{\det g_i},\qquad
 B_i^{jk}=\rho_i(g_i^{-1})^{jk},\qquad
 F_i=-B_iDw_i,\qquad
 b_i=\rho_i\frac{w_i+n/2}{\tau},
\]
where $g_i$ is the pulled-back metric at forward time $-\tau$.
All derivatives here are coordinate derivatives, defined almost
everywhere for Lipschitz functions.  On a smaller cylinder the
$w_i$ have a common Lipschitz bound $L$ and a common value bound.
Smooth metric convergence gives uniform bounds and uniform
ellipticity for $B_i$, as well as uniform convergence to
$B=\rho h^{-1}$.  Thus $|b_i|$ and every component of $F_i$
are bounded by one constant $C$ there.

At each fixed $\tau$, the weak reduced-length inequality takes
the coordinate form
\[
              \int F_i\cdot D\psi\leq\int b_i\psi
              \qquad(\psi\geq0,\ \psi\in C_c^\infty).
\]
These are the coefficients of the actual rescaled metric.
Indeed a rescaling by a backward-time scale $s_i$ multiplies
the volume density by $s_i^{-n/2}$, the inverse metric by $s_i$,
and the source time by $s_i$.  Both sides of this inequality
therefore acquire the same positive factor $s_i^{1-n/2}$.
The coefficient bounds are taken after this cancellation.

Put $O=B_r(a)$, $U=B_{r/2}(a)$ and choose a smooth cutoff
$0\leq\phi\leq1$ supported in $B_{3r/4}(a)$, equal to one
on $U$, with $|\partial_j\phi|\leq D$.
For two sufficiently late indices write $u=w_i(\cdot,\tau)$,
$v=w_j(\cdot,\tau)$, and suppose
\[
 \|u-v\|_{L^\infty(O)}\leq\varepsilon,\qquad
 \max_{k,l}\|B_i^{kl}-B_j^{kl}\|_{L^\infty(O)}\leq\delta.
\]
The functions $\phi(v-u+\varepsilon)$ and
$\phi(u-v+\varepsilon)$ are nonnegative compactly supported
Lipschitz tests.  They are admissible in the weak inequalities:
extend the difference as a Lipschitz function, truncate it to
$[0,2\varepsilon]$, and multiply by $\phi$.
This agrees with the displayed test on its support.  Nonnegative
mollification with sufficiently small support approximates it in
$W^{1,2}$ by nonnegative smooth tests supported in $O$.
The fluxes and sources are in $L^2(O)$, so both integral pairings
pass to the limit by Cauchy--Schwarz.

Use the first test for $F_i,b_i$ and the second for $F_j,b_j$.
Adding their inequalities and expanding the derivatives gives
\[
 \int_O\phi(F_j-F_i)\cdot(Du-Dv)
       \leq4\varepsilon C(1+nD)|O|.
\]
In fact each source term is bounded by $2\varepsilon C$, and
each of the $n$ cutoff terms by $4\varepsilon CD$.
If $B_i\geq cI$ for a fixed $c>0$, then
\[
 (F_j-F_i)\cdot(Du-Dv)
  =B_i(Du-Dv)\cdot(Du-Dv)
       +(B_i-B_j)Dv\cdot(Du-Dv).
\]
Each coordinate derivative has absolute value at most $L$;
the last term has absolute value at most $2n^2\delta L^2$.
Consequently
\[
 c\int_U|Du-Dv|^2
   \leq\bigl(4\varepsilon C(1+nD)+2n^2\delta L^2\bigr)|O|.
\]
Uniform convergence of the functions and coefficient matrices
makes both errors tend to zero.  The coordinate gradients are
therefore Cauchy in $L^2(U)$.

Their limit is $Dw$, not an unidentified vector field.
For every compactly supported smooth $\psi$ and coordinate $k$,
Lipschitz integration by parts gives
$\int w_i\partial_k\psi=-\int\partial_kw_i\psi$.
Uniform convergence on a finite-measure set gives $L^2$ convergence
of $w_i$, while completeness gives the $L^2$ gradient limit.
Passing in this identity identifies that limit with the weak
derivative of $w$.  Rademacher's theorem and uniqueness of weak
derivatives identify it with the almost-everywhere classical
derivative.  Thus $Dw_i\to Dw$ strongly in $L^2(U)$ at every
fixed positive backward time.

The density-weighted energies
\[
 Q_i=\rho_i|\nabla w_i|_{g_i}^2
        =\sum_{j,k} B_i^{jk}\partial_jw_i\partial_kw_i
\]
converge to $Q=\rho|\nabla w|_h^2$ in $L^1(U)$.
To verify this, expand each product difference as
$\xi_i\eta_i-\xi\eta=(\xi_i-\xi)\eta_i+\xi(\eta_i-\eta)$
and use the uniform derivative
bounds and Cauchy--Schwarz.  The coefficient errors vanish
uniformly.  This argument also retains the volume density.
The spacetime Lipschitz bound makes these energies jointly
measurable, and the bounds $|B_i^{jk}|,|B^{jk}|\leq C_1$ give
\[
             \int_U|Q_i-Q|\leq2n^2C_1L^2|U|
\]
uniformly in $\tau\in[\alpha,\beta]$.
Dominated convergence in time and Fubini now give
$Q_i\to Q$ in $L^1(U\times[\alpha,\beta])$.
This is the local energy passage needed below.

\emph{The time derivative.}
On an open chart cylinder the rescaled Hamilton--Jacobi equation is
\[
  2\partial_\tau w_i+Q_i/\rho_i-R_i+w_i/\tau=0
                         \quad\hbox{almost everywhere}.
\]
For each $i$, its slice-wise almost-everywhere version gives the
spacetime version by Fubini: spacetime Lipschitz regularity
identifies the time-slice derivative with the corresponding
directional derivative almost everywhere, and the energy is
jointly measurable.  Intersect the resulting full-measure sets
for the countably many indices.
Local $L^1$ energy convergence admits a further subsequence with
$Q_i\to Q$ almost everywhere.  Positive metric density, smooth
convergence of $\rho_i$ and $R_i$, and uniform convergence of
$w_i$ then imply
\[
       \partial_\tau w_i\longrightarrow
       H:=\tfrac12(R-Q/\rho-w/\tau)
                             \quad\hbox{almost everywhere}.
\]
The common spacetime Lipschitz constant bounds these derivatives.
For every compactly supported smooth spacetime test $\psi$,
dominated convergence passes the integration-by-parts identity
\[
 \int w_i\partial_\tau\psi=-\int\partial_\tau w_i\psi
       \quad\hbox{to}\quad
 \int w\partial_\tau\psi=-\int H\psi.
\]
Uniqueness of distributional derivatives gives
$\partial_\tau w=H$ almost everywhere.  This proves the limiting
Hamilton--Jacobi equation before any smoothness assertion.
Each point is contained in such a chart cylinder; a countable
subcover gives the almost-everywhere equation globally.
No common pointwise-convergent subsequence over all charts is required.

\emph{The weak heat inequality.}
The second weak reduced-length inequality also passes to the limit.
Its flux is linear in $Dw_i$, its energy is $Q_i$, and its remaining
terms converge uniformly on compact sets.  The preceding $L^2$ and
$L^1$ limits therefore give, at each fixed $\tau$, for nonnegative
compact smooth $\psi$,
\[
 \int\left[\psi\left(-|\nabla\ell_\infty|^2+R+
                   \frac{\ell_\infty-n}{\tau}\right)
           -2\langle\nabla\psi,\nabla\ell_\infty\rangle\right]
                 d\mu_{h(-\tau)}\leq0.
\]
As with the comparison tests, positive Sobolev approximation permits
the compact Lipschitz test $\psi=e^{-\ell_\infty}\phi$.
The chain rule, followed by the almost-everywhere time equation,
changes the integrand, after division by two, to
\[
 e^{-\ell_\infty}\left[
    \left(-\ell_{\infty,\tau}+R-\frac n{2\tau}\right)\phi
      -\langle\nabla\phi,\nabla\ell_\infty\rangle\right].
\]
Set $u=\tau^{-n/2}e^{-\ell_\infty}>0$.
In fixed coordinates the backward volume density satisfies
$\partial_\tau\rho=R\rho$.  Thus the first term, multiplied by
$\tau^{-n/2}\rho$, is $\partial_\tau(\rho u)\phi$;
the second is $\rho\langle\nabla u,\nabla\phi\rangle$.
Both $u$ and $\rho u$ are Lipschitz on compact chart cylinders.
Fubini and Lipschitz integration by parts in space and time now give
\[
 T(\phi):=\int_0^\infty\int_{N_\infty}
          u(\partial_\tau\phi+\Delta_{h(-\tau)}\phi)
                         \,d\mu_{h(-\tau)}\,d\tau\geq0
                         \qquad(\phi\geq0).
\]
A finite nonnegative partition of unity on the compact support
passes from chart tests to global tests.  This is precisely the
weak inequality $(\partial_\tau-\Delta+R)u\leq0$.
The $R$ term is absorbed by the time derivative of volume in $T$;
it has not been discarded.

\emph{Gaussian tails and retention of mass.}
Let $p_i$ and $p$ be the retained basepoints.  Put
$C_n^{\rm q}=16(2n+604)^2$ and
$B_{\alpha,\beta}=\frac n2\max\{\beta^2,\alpha^{-2}\}+1$.
The source quadratic estimate gives
\[
 \ell_i(x,\tau)\geq
 \frac{d_{g_i(-\tau)}(p_i,x)^2}{C_n^{\rm q}\tau}
                   -B_{\alpha,\beta},
                \qquad\alpha\leq\tau\leq\beta.
\]
To transfer a sufficient lower bound to the limit, fix $x$ and
write $d=d_{h(-\tau)}(p,x)$.  If $d>0$, source-ball containment
eventually puts the source ball of radius $d/2$ inside the image
of the limit ball of radius $d$.  If the image of $x$ belonged to
that source ball, injectivity of the retained embedding on a common
larger compact region would put $x$ in the open ball of radius $d$,
a contradiction.  Thus the source distance of its image is at
least $d/2$ eventually.  The case $d=0$ is immediate.  It follows that
\[
 \ell_\infty(x,\tau)\geq
       \frac{d_{h(-\tau)}(p,x)^2}{4C_n^{\rm q}\beta}
                         -B_{\alpha,\beta}.
\]
In particular, for $a=(8C_n^{\rm q}\beta)^{-1}$ and
$A=e^{B_{\alpha,\beta}/2}$,
$e^{-\ell_\infty/2}\leq A e^{-a d^2}$ throughout the slab.
Bishop's bound is $\mu(B(p,r))\leq\omega_n r^n$, for source
and limit metrics.  On $d\geq R$ split the Gaussian into two
factors and sum over the annuli $m\leq d<m+1$:
\[
 \int_{d\geq R}e^{-a d^2}\,d\mu
   \leq e^{-aR^2/2}\omega_n
        \sum_{m=0}^\infty (m+1)^n e^{-am^2/2}.
\]
The series converges, since a Gaussian dominates every polynomial.
This proves uniform tails on the slab, also for the source
densities $u_i=\tau^{-n/2}e^{-\ell_i}$.

Here is the complete mass comparison at a fixed $\tau$.
For a relatively compact open set $V$ in the limit and any $C>1$,
the retained metric convergence eventually bounds tangent norms
in both directions by $C$.  The Jacobian bounds are consequently
$C^{-n}$ and $C^n$.  Change of variables compares the integral of
$u_i$ over the image of $V$ with its pulled-back integral over $V$
within these factors.  The latter converges to $\int_Vu$ by
uniform convergence and finite volume.  Each image integral is
bounded by the full source mass.  Letting $i\to\infty$, then
$C\downarrow1$, and exhausting the limit shows
$\int u\leq V_\infty$.

For the reverse inequality, choose $R$ so that every source tail
outside $B(p_i,R)$ is at most $\varepsilon$.  Completeness makes
the closure of the limit ball $V=B(p,2R)$ compact, and source-ball
containment puts $B(p_i,R)$ inside its image for all sufficiently
late $i$.  Thus the full source mass is at most
$C^n\int_V(u_i\hbox{ pulled back})+\varepsilon$.
First pass to the limit, then let $C\downarrow1$ and
$\varepsilon\downarrow0$.  We obtain
\[
 \int_{N_\infty}u(\cdot,\tau)\,d\mu_{h(-\tau)}
     =V_\infty:=\lim_{s\to\infty}V_{\rm red}(s),
                    \qquad\tau>0.
\]
Every dilated time $s_i\tau$ tends to infinity, so the right side
is independent of $\tau$.  Scalar positivity of the original
nonflat ancient flow, reduced-volume monotonicity and its strict
Gaussian rigidity give $V_\infty<(4\pi)^{n/2}$.

\emph{Vanishing cutoff error and equality.}
The source gradient estimate passes by local energy convergence:
an almost-everywhere subsequence of the weighted energies, followed
by division by the positive limiting density, gives
\[
       |\nabla\ell_\infty|^2\leq3\ell_\infty/\tau
                  \quad\hbox{almost everywhere on each slice}.
\]
Countably many coordinate balls globalize this assertion.  Since
$\ell_\infty\geq0$ and $\sqrt y\,e^{-y}\leq e^{-y/2}$ for $y\geq0$,
the Gaussian estimate above supplies one bound $B$ for
$\int\sqrt{\ell_\infty}e^{-\ell_\infty}\,d\mu_{h(-\tau)}$
over $\tau\in[\alpha,\beta]$.

Choose the cutoffs in the single complete metric $h(-\alpha)$.
To smooth its distance $r=d(p,\cdot)$, choose small convex chart
balls whose tangent norms and the norm at their centers differ
by factors at most $6/5$.  The triangle inequality and lengths
of chart segments make $r$ Lipschitz with constant $6/5$ in
each fixed center norm.  Extend this scalar function with the
same Lipschitz constant and convolve with a positive mollifier.
Pullback gives local smooth approximations $r_j$ with arbitrarily
small value error and $|dr_j|\leq(6/5)^2<3/2$.

Take a countable locally finite smooth partition $\theta_j$ with
compact supports subordinate to these balls, and let
$B_j\geq\sup|d\theta_j|$.  Choose the local value errors to be
at most $\varepsilon_j=\min\{1,2^{-j}/(4(B_j+1))\}$.
Then $v=\sum_j\theta_jr_j$ is smooth, $|v-r|\leq1$, and
\[
 dv=\sum_j\theta_j\,dr_j+\sum_j(r_j-r)\,d\theta_j,
 \qquad |dv|\leq\tfrac32+\sum_{j=0}^\infty2^{-j}/4\leq2.
\]
The cancellation uses $\sum_jd\theta_j=0$, and does not require
differentiating the nonsmooth $r$.
This smooth approximation $v$, with
$|v-d(p,\cdot)|\leq1$ and $|\nabla v|\leq2$, composed with a
fixed bump $b(v/(2R))$, gives $0\leq\chi_R\leq1$,
$\chi_R=1$ on $B(p,R)$, support in $\overline B(p,5R)$ and
$|\nabla\chi_R|\leq C_0/R$ for $R\geq1$.
Here $b=1$ on $[-1,1]$, vanishes outside $(-2,2)$, and
$|b'|\leq C_0$; the chain rule gives the claimed bound.
Completeness makes the support compact.
Backward nonnegative Ricci curvature makes $h(-\tau)$ increase
with $\tau$, so dual norms decrease: the same gradient bound
holds for all $\tau\geq\alpha$ with this fixed cutoff.

Cauchy--Schwarz in each tangent space now gives
\[
 \int u\,|\langle\nabla\ell_\infty,\nabla\chi_R\rangle|
                      \,d\mu_{h(-\tau)}
 \leq \frac{\alpha^{-n/2}\sqrt{3/\alpha}\,C_0B}{R}.
\]
Lipschitz Green's formula and the exponential chain rule give
$\int u\Delta\chi_R=\int u\langle\nabla\ell_\infty,
\nabla\chi_R\rangle$.  Thus the integrated Laplacian error
tends uniformly to zero on the slab, without a pointwise
Laplacian bound on $\chi_R$.

For a smooth time test $\eta$ supported in $(\alpha,\beta)$,
\[
 T(\chi_R\eta)=\int_\alpha^\beta
    \left(\eta'(\tau)\int u\chi_R\,d\mu
             +\eta(\tau)\int u\Delta\chi_R\,d\mu\right)d\tau
                       \longrightarrow V_\infty\int\eta'=0.
\]
Indeed $0\leq\int u\chi_R\leq V_\infty$ and these integrals
tend to $V_\infty$ at every time, so dominated convergence
handles the first term; the uniform flux bound handles the second.

Finally let $\phi$ be any compact smooth spacetime test in the
slab.  Choose $0\leq\eta\leq1$ equal to one on its time support,
and $K\geq\sup|\phi|$.  Eventually $\chi_R=1$ on its spatial
support, so $K\chi_R\eta\pm\phi\geq0$ everywhere.
Positivity and linearity imply
$|T(\phi)|\leq K T(\chi_R\eta)\to0$.
Every compact positive-time support lies in such a slab.
Consequently $T(\phi)=0$ for all compact tests, which is the
weak conjugate-heat equation for this same $u$.

\emph{Smoothness and the scalar identities.}
In a fixed chart write
$\Delta\phi=a^{jk}\partial_j\partial_k\phi+b^j\partial_j\phi$,
where $a=h^{-1}$ and the coefficients are smooth.  The coordinate
density $v=\rho u$, which is continuous and positive, satisfies
\[
 \partial_\tau v=\partial_j\partial_k(a^{jk}v)
                         -\partial_j(b^jv)
                         \quad\hbox{in distributions}.
\]
The principal matrix is positive definite; on every smaller compact
chart cylinder its smallest eigenvalue has a positive lower bound.
Expanding its distributional divergences puts it in the form of
Lemma~\ref{lem:ancient-weak-regularity}, with principal matrix $a$
and lower coefficients
\[
 \widehat b^{\,i}=b^i-\partial_j a^{ij}-\partial_j a^{ji},\qquad
 \widehat c=\partial_i\widehat b^{\,i}
                              +\partial_i\partial_j a^{ij}.
\]
These coefficients are smooth, and the lemma proves smoothness of
the original continuous $v$.  Thus
$\ell_\infty=-\log(v/(\rho\tau^{-n/2}))$ is smooth on the
original carrier.

The weak heat equality now holds pointwise, and taking its logarithm
gives the first identity below.  The almost-everywhere
Hamilton--Jacobi identity also holds everywhere: its now continuous
left side cannot be nonzero on a set of measure zero.  We have proved
\[
 \ell_{\infty,\tau}-\Delta\ell_\infty+|\nabla\ell_\infty|^2
                  -R+\frac n{2\tau}=0,\qquad
 2\ell_{\infty,\tau}+|\nabla\ell_\infty|^2-R
                  +\frac{\ell_\infty}{\tau}=0.
\]
Subtracting twice the first equation from the second yields
\[
 \tau(2\Delta\ell_\infty-|\nabla\ell_\infty|^2+R)+\ell_\infty-n=0.
\]
\emph{The tensor equation.}
Write $f=\ell_\infty$ temporarily in backward time, $q=|\nabla f|^2$,
$W=2\Delta f-q+R$, and
$\mathcal L=\partial_\tau-\Delta+2\langle\nabla f,\nabla\cdot\rangle$.
The heat identity is $f_\tau=\Delta f-q+R-n/(2\tau)$.
The backward metric variation, Bochner formula and scalar evolution give
\begin{align*}
 \mathcal L(\Delta f)
   &=-2|\operatorname{Hess}f|^2-2\operatorname{Ric}(\nabla f,\nabla f)
       +\Delta R-2\langle\operatorname{Ric},\operatorname{Hess}f\rangle,\\
 \mathcal Lq
   &=-4\operatorname{Ric}(\nabla f,\nabla f)
       +2\langle\nabla f,\nabla R\rangle-2|\operatorname{Hess}f|^2,\\
 \mathcal LR
   &=-2\Delta R-2|\operatorname{Ric}|^2
       +2\langle\nabla f,\nabla R\rangle.
\end{align*}
For example, $\partial_\tau\Delta f=\Delta f_\tau
-2\langle\operatorname{Ric},\operatorname{Hess}f\rangle$;
the contracted variation of the connection cancels by contracted
Bianchi.  Applying Bochner to $\Delta q$ then gives the first
two displayed identities.  Their combination is
$\mathcal LW=-2|\operatorname{Ric}+\operatorname{Hess}f|^2$.
Since $\mathcal Lf=q+R-n/(2\tau)$, the scalar expression
$P=\tau W+f-n$, already proved zero everywhere, satisfies
\[
 \mathcal LP=-2\tau
       \left|\operatorname{Ric}+\operatorname{Hess}f
                                   -\frac{h}{2\tau}\right|^2.
\]
Here completing the square uses
$\operatorname{tr}(\operatorname{Ric}+\operatorname{Hess}f)
=R+\Delta f$ and $|h|^2=n$.
The left side vanishes; since $\tau>0$, the tensor square is zero.
An orthonormal basis expresses it as a sum of squares of all
components, so the tensor itself vanishes.  Reversing time proves
(A) for $f(x,t)=\ell_\infty(x,-t)$.

\emph{Excluding a flat slice.}
If a slice were flat, (A) and the vanishing scalar expression would
give $\operatorname{Hess}f=h/(2\tau)$ and $f=\tau|\nabla f|^2$.
We verify the resulting global Gaussian normalization.
Along any complete geodesic $\gamma$ with $\gamma(0)=q$ and
$\gamma'(0)=v$, its constant speed and the Hessian equation give
\[
 f(\gamma(s))=f(q)+s\langle\nabla f(q),v\rangle
                                      +\frac{s^2|v|^2}{4\tau}.
\]
Take $v=-2\tau\nabla f(q)$ and $s=1$.  The gradient identity
makes this value zero.  Thus there is a point $p$ with $f(p)=0$;
the same identity gives $\nabla f(p)=0$.
No prior properness or existence of a minimum has been assumed.

For a geodesic issuing from $p$ with initial velocity $v$, put
$x=\gamma(1)$ and $a=|v|^2$.  Differentiating the displayed
quadratic and using constant speed gives
\[
 |\gamma'(1)|^2=a,\qquad
 \langle\nabla f(x),\gamma'(1)\rangle=\frac a{2\tau},\qquad
 |\nabla f(x)|^2=\frac a{4\tau^2}.
\]
Hence $|\gamma'(1)-2\tau\nabla f(x)|^2=0$.
If two such geodesics have the same endpoint, their terminal
velocities agree.  Uniqueness for the geodesic equation, run
backward, gives the same initial velocity.  Therefore
$\exp_p:T_pN_\infty\longrightarrow N_\infty$ is injective.
Hopf--Rinow gives a minimizing geodesic to every point, so it is
also surjective.

Choose an orthonormal identification $T_pN_\infty\cong\mathbb R^n$.
The complete exponential is smooth on every finite ball by smooth
dependence of the geodesic equation.  Its derivative in direction
$w$ is the endpoint of the radial Jacobi field with $J(0)=0$ and
$\nabla_sJ(0)=w$.  Flatness gives $\nabla_s^2J=0$, so
$J(s)=s\mathsf P_s w$, where $\mathsf P_s$ is parallel transport.
It follows that $(\exp_p)^*h$ is the Euclidean metric everywhere.
In particular its derivative is invertible.  The inverse function
theorem and the established bijectivity give a global smooth
inverse.  These maps preserve curve lengths in both directions,
hence distance and the corresponding Riemannian volume.
The quadratic formula becomes
$f(\exp_p v)=|v|^2/(4\tau)$, with no additive constant.
Consequently the mass is
\[
 \tau^{-n/2}\int_{\mathbb R^n}e^{-|v|^2/(4\tau)}\,dv
   =\tau^{-n/2}(4\pi\tau)^{n/2}=(4\pi)^{n/2},
\]
contradicting the strict mass gap.  The same argument applies to
any chosen negative time, so every limit slice is nonflat.
\end{proof}

The construction has a fixed order.  The
reduced-length Harnack estimates first give curvature bounds on each
buffered window, with the spatial ball and the curvature-evaluation time
chosen independently.  The complete pointed compactness argument then
produces one carrier with compatible finite-window embeddings; distance
distortion, ball coverage, and the strict-radius volume squeeze identify the
same limit at overlapping times and retain completeness, curvature, and
noncollapsing.  A countable space-time diagonal gives all metric jets on the
exhausting coordinate sets, rather than only a sequence of separate slice
limits.

The reduced-distance functions are extracted after the geometric diagonal by
connected local oscillation bounds and overlap agreement.  Gaussian shell
estimates give a constant actual mass, while the two weighted comparison
inequalities pass to a weak heat equality and its Hamilton--Jacobi identity.
The latter passage uses the compactly supported gradient-energy and cutoff
estimates proved above; local uniform convergence by itself would not
justify the quadratic term.  Lemma~\ref{lem:ancient-weak-regularity}
upgrades the weak equality to a smooth equation.

Finally, the smooth heat and Hamilton--Jacobi identities make the
scalar defect identically zero.  Its pointwise evolution then forces
the squared norm of
$\operatorname{Ric}+\operatorname{Hess}f+(2t)^{-1}h$ to vanish,
giving the soliton equation (A).  The complete smooth
negative-gradient flow and its pullback identity will identify
the retained solution as self-similar on this same carrier.
Compare \cite[Proposition 9.32 and Lemma 9.33,
pp.~200--202]{MT2007}.

\begin{remark}[Analytic inputs]
\label{rem:ancient-energy-interface}
The energy and time-derivative passages above use the reduced-length
weak inequalities and uniform estimates from
Chapter~\ref{ch:reduced-geometry}.  The Gaussian exhaustion and weak
heat equality and local parabolic regularity are developed above.
The finite-window identifications and ancient diagonal are developed
above, as is the extraction of the locally Lipschitz potential.
The complete pointed-compactness input is
Theorem~\ref{thm:rf-compactness}.  The regularity lemma uses
Sobolev and Morrey embeddings; the flat Gaussian argument uses
complete minimizing geodesics, radial Jacobi fields and the
Riemannian change-of-variables formula.
\end{remark}

\section{Surface rigidity}

A static gradient shrinking soliton is a complete metric $h$ and a
$C^2$ function $f$ satisfying
\[
       \operatorname{Ric}_h+\operatorname{Hess}_h f=\lambda h,
                         \qquad\lambda>0.                 \tag{B}
\]
The admissible static solitons used later additionally have nonnegative
curvature operator, bounded curvature, nonflatness, and metric
$\kappa$-noncollapsing for some $\kappa>0$.  Their normalization is
$\lambda=1/2$.

\begin{lemma}[Complete surface shrinkers]
\label{lem:surface-shrinker}
A complete connected nonflat surface with nonnegative curvature
satisfying (B) is compact and has constant Gaussian curvature $\lambda$.
Compactness here needs no curvature bound or noncollapsing hypothesis.
\source{MT2007,ChowKnopf2004}
\dcref{Morgan--Tian, Theorem 9.42 and its compact case, Claim 9.43,
pp. 206--208; Chow--Knopf, Proposition 5.21, pp. 118--119;
the surface argument below also supplies compactness under the stated weaker hypotheses}
\end{lemma}
\begin{proof}
\emph{Conserved quantities and compactness.}
Since $\operatorname{Ric}=(R/2)h$, the Hessian equation is
$\operatorname{Hess}f=(\lambda-R/2)h$.  This first makes $f$ smooth.
Indeed, in smooth coordinates,
\[
 \partial_i\partial_jf
   =\Gamma_{ij}^k\partial_kf+(\lambda-R/2)h_{ij}.
\]
If $f$ is $C^k$, with $k\geq2$, the right side is $C^{k-1}$,
so all second derivatives are $C^{k-1}$ and $f$ is $C^{k+1}$.
Iteration starts from the assumed $C^2$ regularity.  Taking a divergence
of the Hessian equation and using
$\Delta f=2\lambda-R$ gives
$-\tfrac12dR=-dR+\tfrac R2df$, hence $dR=R\,df$.
Differentiating the gradient norm then gives
\[
             dR=R\,df,\qquad R+|\nabla f|^2-2\lambda f=C .
\]
Connectedness and nonflatness imply $Re^{-f}=a>0$.
The second identity gives $f\geq-C/(2\lambda)$, hence the uniform
bound $R\geq ae^{-C/(2\lambda)}>0$.  Myers' theorem gives compactness.

\emph{The critical levels.}
Suppose $f$ is not constant, and let $m<M$ be its minimum and maximum.
At any critical point $p$, the conserved quantities imply
\[
 |\nabla f|^2(x)
   =2\lambda(f(x)-f(p))-R(p)\bigl(e^{f(x)-f(p)}-1\bigr).
\]
If $R(p)=2\lambda$, the inequality $e^u>1+u$ for $u\ne0$
forces $f(x)=f(p)$ everywhere, a contradiction.
The same inequality shows that a critical point with $R(p)<2\lambda$
is a global minimum: a negative value of $u=f(x)-f(p)$ would make
the displayed right side negative.  A critical point with
$R(p)>2\lambda$ is similarly a global maximum.
Thus every value in $(m,M)$ is regular, and both endpoint Hessians
are definite.

We also need uniqueness at each endpoint.  Let $p$ be a critical
point and suppose $q\ne p$ satisfies $f(q)=f(p)$.  A minimizing
geodesic $\gamma:[0,1]\longrightarrow N$ from $p$ to $q$ has constant
speed $v>0$.  Set $\phi=(f\circ\gamma)'$.  The Hessian and scalar
identities give
\[
 \phi'=(\lambda-R(\gamma)/2)v^2,
 \qquad \phi''=-\tfrac12R(\gamma)v^2\phi,
 \qquad \phi(0)=0,\quad \phi'(0)\ne0.
\]
For a parallel unit normal $E$ along $\gamma$, the field $J=\phi E$
is a nontrivial Jacobi field.  It cannot vanish at an interior time
$c$.  To see the obstruction, extend $J|_{[0,c]}$ by zero to $[c,1]$
and call the resulting continuous piecewise smooth field $W$.
For the index form
\[
 I(U,V)=\int_0^1\left(
   \langle D_tU,D_tV\rangle
   -\langle\operatorname{Rm}(U,\dot\gamma)\dot\gamma,V\rangle
                  \right)\,dt,
\]
integration by parts gives $I(W,W)=0$ and
$I(W,V)=\langle D_tJ(c),V(c)\rangle$ for every smooth field $V$
vanishing at the endpoints.  Uniqueness for the Jacobi equation gives
$D_tJ(c)\ne0$.  Choose $V(c)=-D_tJ(c)$.  For sufficiently small
$\epsilon>0$, $I(W+\epsilon V,W+\epsilon V)<0$, contradicting the
second variation for a minimizing segment.  The piecewise smooth
fields are admissible by using a fixed-endpoint variation on each
side of $c$ with the same value there.
But Rolle's theorem applied to $f\circ\gamma$ supplies just such
an interior zero of $\phi$.  This proves that each critical level is
a singleton.  Write its two points as $p$ and $z$, with $f(p)=m$
and $f(z)=M$, and put $a_0=\lambda-R(p)/2>0$.

\emph{The coarea density.}
For $m<s<M$ define
\[
 r(s)=R(p)e^{s-m},\qquad
 q(s)=2\lambda(s-m)-R(p)(e^{s-m}-1).
\]
On $f^{-1}(s)$ these are $R$ and $|\nabla f|^2$, respectively;
$q(s)>0$ and $q'(s)=2\lambda-r(s)$.
Let $L(s)$ denote the total length of this compact regular level.
On each closed interval of regular values, the flow of
$\nabla f/|\nabla f|^2$ identifies the levels.  It exists throughout
that interval because the slab is compact and its gradient is bounded
away from zero.  For a unit tangent $T$ to a level, its length element
has logarithmic derivative
\[
 \left\langle\nabla_T\left(\frac{\nabla f}{|\nabla f|^2}\right),T\right\rangle
   =\frac{\lambda-r(s)/2}{q(s)}.
\]
The same expression is the logarithmic derivative of $\sqrt{q(s)}$.
Consequently $L(s)/\sqrt{q(s)}=c$ is a single constant on $(m,M)$.
This uses the total level length and requires no prior identification
of the components as circles.  Coarea on a compact regular slab gives,
for every continuous function $H$ and $m<s<t<M$,
\[
 \int_{\{s\leq f\leq t\}}H(f)\,d\mu
       =c\int_s^t H(u)\,du.
\]
Taking $H=1$ and exhausting an open slab by closed regular slabs yields
$\mu\{m<f<t\}=c(t-m)$.

The value of $c$ is determined at the unique minimum.  In coordinates
centered at $p$ with $h_{ij}(p)=\delta_{ij}$,
\[
 f(x)-m=\tfrac12a_0|x|^2+o(|x|^2),
       \qquad d\mu=(1+o(1))\,dx.
\]
For each $K>1$, on a sufficiently small coordinate ball the quadratic
bounds have coefficients $a_0/(2K)$ and $a_0K/2$, and the volume
density lies between $K^{-1}$ and $K$.  Compactness and uniqueness
of the minimum ensure that every sufficiently small sublevel is
contained in this ball.  Comparing its coordinate image with the
disks of radii $\sqrt{2\eta/(a_0K)}$ and
$\sqrt{2K\eta/a_0}$ gives
\[
 \frac{2\pi}{a_0K^2}
    \leq\frac{\mu\{f<m+\eta\}}{\eta}
    \leq\frac{2\pi K^2}{a_0}.
\]
The singleton $\{p\}$ has zero area.  The open-slab formula therefore
identifies the middle expression with $c$ for small $\eta>0$.
Letting $K\downarrow1$ proves $c=2\pi/a_0$.

\emph{The curvature contradiction.}
Put $d=M-m>0$.  At the extrema,
\[
             R(z)=R(p)e^d,\qquad R(z)-R(p)=2\lambda d .
\]
Strict convexity of the exponential gives
$\int_0^d e^u\,du<d(1+e^d)/2$.
Substituting the two endpoint identities yields
$R(p)+R(z)>4\lambda$, or $R(z)-R(p)>4a_0$.
Use $H=r$ in the slab formula.  Since $R\geq0$, its integral over
each slab is at most the total scalar curvature.  Letting the regular
endpoints approach $m$ and $M$ gives
\[
 8\pi<\frac{2\pi}{a_0}(R(z)-R(p))
                  \leq\int_N R\,d\mu.
\]
This contradicts Gauss--Bonnet,
$\int R=4\pi\chi(N)\leq8\pi$, also valid for a nonorientable surface.
Therefore $f$ is constant, $R=2\lambda$, and the Gaussian curvature
is $\lambda$.
\end{proof}

\begin{theorem}[Ancient surface solutions]
\label{thm:surface-ancient}
Every two-dimensional ancient $\kappa$-solution is compact and has
constant positive sectional curvature on each slice, including time
zero.  Every two-dimensional asymptotic limit of
Theorem~\ref{thm:ancient-asymptotic} is compact and round at every
negative time.
\lean{PoincareMT.twoDimensionalAncientAndShrinkingSolitonClassification}
\uses{lem:surface-shrinker,thm:ancient-asymptotic}
\source{MT2007,ChowKnopf2004}
\dcref{Morgan--Tian, Corollaries 9.49--9.50, pp. 213--214;
Chow--Knopf, Lemma 5.38 and Proposition 5.39, pp. 133--134}
\end{theorem}
\begin{proof}
Apply the lemma to the limit equation on each slice with
$\lambda=-1/(2t)$.  At time $-1$ its scalar curvature is therefore
$1$.  The increasing open convergence domains cover this compact
limit, so one domain, and every later domain, contains the whole
surface.  The retained spatial embeddings are then injective local
diffeomorphisms on the whole limit.  Each image is open by the inverse
function theorem and closed by compactness, hence is the entire
original connected surface.  Thus the original surface is compact,
and the late convergence maps are global diffeomorphisms.

Scalar curvature is positive on every ancient slice.  Indeed,
$R_t=\Delta R+R^2$ makes $R\geq0$ a heat supersolution.  Nonflatness
at time $t-1$ gives a point of positive scalar curvature there, and
the strong maximum principle on $[t-1,t]$ propagates positivity to
every point at time $t\leq0$.  This is the scalar argument of
Proposition~\ref{prop:rf-positivity}; it supplies the strict positivity
needed for logarithms below.  Only the scalar maximum-principle step
of that three-dimensional proposition is used here.

\emph{Entropy and its derivative.}
To transfer exact roundness, define on a compact positive-scalar surface
\[
 \bar R=\frac{\int R\,d\mu}{\mu(N)},\qquad
 {\cal E}(g)=\int_N\bigl(R\log(R/\bar R)-R+\bar R\bigr)\,d\mu .
\]
For $r,x>0$, the function $x\log(x/r)-x+r$ is nonnegative and vanishes
only at $x=r$: its derivative is $\log(x/r)$ and its second derivative
is $1/x$.  If its integral vanishes, continuity and the positive area
of every nonempty open set make the density vanish everywhere.
Thus ${\cal E}\geq0$, with equality exactly for constant scalar
curvature.  Under $g\mapsto bg$, both $R$ and $\bar R$ are divided
by $b$, while $d\mu$ is multiplied by $b$ in dimension two.
The density is divided by $b$, so ${\cal E}$ is invariant under
homothety.

Write $V=\mu(N)$ and $S=\int R\,d\mu$.  The scalar equation and
$\partial_td\mu=-R\,d\mu$ give $V'=-S$ and
$S'=\int\Delta R\,d\mu=0$.  Consequently
$\bar R'=\bar R^2$.  Since the last two terms of the entropy
density have zero integral,
\[
 {\cal E}=\int R\log R\,d\mu-S\log\bar R.
\]
On a compact time interval inside the flow interval all integrands
are smooth and $R$ has a positive lower bound, so differentiation
under the integral is legitimate.  The evolving volume term cancels
the $R^2\log R$ term, giving
\[
 \frac d{dt}\int R\log R\,d\mu
    =\int\bigl((1+\log R)\Delta R+R^2\bigr)\,d\mu
    =\int R^2\,d\mu-\int\frac{|\nabla R|^2}{R}\,d\mu.
\]
Subtracting $(S\log\bar R)'=S\bar R$ therefore yields
\[
 \frac{d{\cal E}}{dt}
       =\int(R-\bar R)^2\,d\mu-\int\frac{|\nabla R|^2}{R}\,d\mu
       \leq0.                                           \tag{C}
\]
\emph{The Fisher information inequality.}
For any smooth real $v$ on the fixed closed surface, integrate the
Bochner identity and use $\operatorname{Ric}=(R/2)g$ to obtain
\[
 2\int(\Delta v)^2\,d\mu
    =2\int|\operatorname{Hess}v|^2\,d\mu
                         +\int R|\nabla v|^2\,d\mu.
\]
The pointwise inequality
$2|\operatorname{Hess}v|^2\geq(\Delta v)^2$ then gives
$\int R|\nabla v|^2\leq\int(\Delta v)^2$.
Expand the nonnegative integral of $|\nabla R+R\nabla v|^2/R$,
integrate its mixed term by parts, and use $\int\Delta v=0$:
\[
 2\int(R-\bar R)\Delta v
       \leq\int|\nabla R|^2/R+\int(\Delta v)^2 .
\]
Here and in the next paragraph all integrals use the fixed metric.
To choose the test functions, let $e_j$ be a smooth orthonormal
eigenbasis for the nonnegative operator $-\Delta$, with eigenvalues
$\lambda_j\geq0$.  One obtains this basis from the compact positive
self-adjoint resolvent $(1-\Delta)^{-1}$ on $L^2$: the weak energy
form gives its bounded inverse into $H^1$, Rellich compactness gives
compactness in $L^2$, and local elliptic regularity makes each weak
eigenfunction smooth.  A zero eigenfunction is constant, since
$\int|\nabla e_j|^2=-\int e_j\Delta e_j=0$ and the surface is
connected.
For $F=R-\bar R$, therefore, every zero-mode coefficient vanishes.
If $F=\sum_j a_je_j$ in $L^2$, put
\[
 v_A=-\sum_{j\in A,\,\lambda_j>0}\frac{a_j}{\lambda_j}e_j
\]
for finite sets $A$ exhausting the eigenbasis.  These are smooth and
$\Delta v_A\longrightarrow F$ in $L^2$.
The preceding inequality passes to the limit because its mixed
integral and squared norm are continuous in $\Delta v_A$ for this
topology.  It becomes
\[
 2\|F\|_2^2\leq\int\frac{|\nabla R|^2}{R}\,d\mu+\|F\|_2^2,
\]
which is the last inequality in (C).  This argument only requires
the finite spectral approximants, without assuming convergence of
the potentials $v_A$ themselves.

\emph{Passage from the compact limit.}
Let $g_i=\tau_{\sigma(i)}^{-1}g(-\tau_{\sigma(i)})$.
The retained global diffeomorphisms pull these metrics back to metrics
converging smoothly to the compact round limit at time $-1$.
In any fixed coordinate chart, connection coefficients are smooth
functions of the positive definite metric matrix and its first
derivatives.  Curvature additionally uses its second derivatives.
Thus convergence through order two gives scalar convergence, also
at moving points of the chart.  A finite cover of the compact limit
makes this convergence uniform.  Surjectivity of the same
diffeomorphisms yields
\[
       \sup_N|R_{g_i}-1|\longrightarrow0.
\]
In particular, $R_{g_i}\geq1/2$ eventually.  Gauss--Bonnet bounds
$\int R_{g_i}\,d\mu_{g_i}\leq8\pi$, so
$\mu_{g_i}(N)\leq16\pi$.
Moreover,
\[
 |\bar R_{g_i}-1|\leq\sup_N|R_{g_i}-1|\longrightarrow0.
\]
The continuous relative density
$(r,x)\mapsto x\log(x/r)-x+r$ vanishes at $(1,1)$; its values
at $(\bar R_{g_i},R_{g_i}(x))$ therefore tend uniformly to zero.
The uniform area bound now gives ${\cal E}(g_i)\longrightarrow0$.
Homothety invariance identifies this exactly with the entropy of the
original slices at times $-\tau_{\sigma(i)}$, which tend to $-\infty$.
For any fixed $t<0$ and all large $i$,
\[
 0\leq{\cal E}(g(t))\leq{\cal E}(g(-\tau_{\sigma(i)})).
\]
Thus $R(t,\cdot)$ is constant.  At time zero, continuity in time
gives $R(0,x)=R(0,p)$ for every fixed $x,p$.  This continuity can also
be read directly from
$R(t,x)=2\operatorname{Ric}_{g(t)}(v,v)/g(t)(v,v)$ for any fixed
nonzero tangent vector $v$ at $x$, using the metric evolution and
its regularity up to the terminal time.  Nonflatness at zero and
nonnegative curvature make this constant positive.  This permits
both the sphere and the projective plane; no orientability was assumed.
\end{proof}

\section{Three-dimensional shrinking models}

\begin{lemma}[A line in a retained pointed limit]
\label{lem:ancient-geometric-selected-line}
Let $(M,g)$ be a complete connected $n$-manifold, $n\geq2$,
with nonnegative sectional curvature.  Fix $p\in M$, points
$q_i$, and positive numbers $Q_i,r_i$.  Put
\[
 d_i=d_g(p,q_i),\qquad D_i=d_i\sqrt{Q_i},
 \qquad L_i=r_i\sqrt{Q_i}.
\]
Suppose $d_i,D_i,L_i\to\infty$ and $r_i/d_i\to0$.
Fix $\delta\geq0$.  For $\delta>0$, let $H_i(u)$ be complete
Ricci flows on open intervals containing $[-\delta,0]$, with
$H_i(0)=Q_i g$ and nonnegative sectional curvature.  Suppose
there is a fixed $\Lambda\geq0$ such that
\[
 0\leq\operatorname{Ric}_{H_i(u)}
       \leq\Lambda H_i(u)
 \quad\hbox{on }B_{H_i(0)}(q_i,L_i),\qquad -\delta\leq u\leq0.
\]
For $\delta=0$ only the metrics $H_i(0)=Q_i g$ are required.
If the pointed metrics $(M,H_i(-\delta),q_i)$ converge smoothly
along supplied exhaustion embeddings to a complete pointed
manifold $(M_\infty,g_\infty,o)$, that same limit contains an
isometric line $\gamma:\mathbb R\to M_\infty$ through $o$.
\end{lemma}
\begin{proof}
Restrict first to the subsequence defining the supplied pointed limit.
All divergence statements survive further subsequences, which retain
that limit and its embeddings.
All unmarked distances and tangent inner products in the selection
below are for $g$.  Discard finitely many terms so that $d_i>0$.
If $\delta=0$, set $\Lambda=0$.
Unit initial directions of minimizing geodesics from $p$ to
$q_i$ have a convergent subsequence, say $u_i\to u$.
Choose farther selected vertices $y_i=q_{j(i)}$, with $j(i)\geq i$,
such that $d(p,y_i)\geq(i+2)d_i+1$.
The hinge inequality of
Lemma~\ref{lem:rf-squared-distance-comparison} gives
\[
 d(q_i,y_i)^2\leq d_i^2+d(p,y_i)^2
                   -2d_i d(p,y_i)\langle u_i,u_{j(i)}\rangle.
\]
Let $c_i$ be the Euclidean comparison cosine at $q_i$ in the
triangle $p,q_i,y_i$.  The triangle inequality gives $c_i\geq-1$.
Writing $\epsilon_i=1/(i+1)$, the farther-vertex inequality
and the displayed hinge bound give, eventually,
\[
 c_i\leq\epsilon_i-(1-\epsilon_i)
                         \langle u_i,u_{j(i)}\rangle\longrightarrow-1.
\]
This cosine is unchanged by scalar rescaling.
With $K_0=4e^{\Lambda\delta}$ choose
\[
 s_i=\min\{\sqrt{L_i},\ L_i/(K_0+1),\ D_i/2\}.
\]
Then $s_i\to\infty$, $s_i/L_i\to0$, and $K_0s_i\leq L_i$.
In $H_i(0)$, trim the minimizing segments from $q_i$ toward
$p$ and $y_i$ to length $s_i/2$; call their endpoints
$x_i^-$ and $x_i^+$.  Both endpoints lie in $B(q_i,s_i)$.
Corresponding-side comparison makes their new comparison
cosine at most $c_i$ and at least $-1$, so it too tends to $-1$.

If $\delta=0$, retain these segments and set $C_\delta=0$ below.
For $\delta>0$, the earlier metric $H_i(-\delta)$ requires an
additive distance estimate.  On the controlled terminal ball, $0\leq\operatorname{Ric}
\leq\Lambda H_i$, and the margin $4e^{\Lambda\delta}s_i\leq L_i$ supplies
the current balls needed for that estimate.  Indeed, metric
comparison places the endpoints at current distance less than
$2e^{\Lambda\delta}s_i$.  The current ball of radius
$3e^{\Lambda\delta}s_i$ about either endpoint lies in the controlled
terminal ball, by forward distance monotonicity and the triangle
inequality.  In particular the minimizing segment between the
endpoints stays in this region.  On a segment of length $\ell$,
use parallel normal fields with scalar cutoff
\[
 \phi(v)=(1-e^{-\lambda v})(1-e^{-\lambda(\ell-v)}),
 \qquad\lambda>0.
\]
Writing the two exponentials as $a$ and $b$, one has
$\phi'=\lambda(a-b)$ and $1-\phi^2\leq2(a+b)$.
Consequently $\int(\phi')^2\leq2\lambda$ and
$\int(1-\phi^2)\leq4/\lambda$.  The index inequality and
$0\leq\operatorname{Ric}(\dot\gamma,\dot\gamma)\leq\Lambda$
give
\[
 \int\operatorname{Ric}(\dot\gamma,\dot\gamma)
 \leq n\int(\phi')^2+\Lambda\int(1-\phi^2)
 \leq2n\lambda+4\Lambda/\lambda.
\]
The time-distance support is obtained by splitting a minimizing
segment at an interior point and retaining the two paths while
varying time.  Their total length touches distance from above,
and its derivative is the negative sum of their Ricci integrals.
Applying the displayed estimate to both segments gives the
coarser bound $4n\lambda+8\Lambda/\lambda$.  At coincident endpoints
distance is identically zero.  Thus the support estimate gives
\[
 d_0(x,y)\leq d_{-\delta}(x,y)
 \leq d_0(x,y)+C_\delta,\qquad
 C_\delta=(4n\lambda+8\Lambda/\lambda)\delta,
 \quad x,y\in B_{H_i(0)}(q_i,s_i).
\]
To integrate the one-sided estimate, reverse time: an upper
distance support with forward derivative at least
$-(4n\lambda+8\Lambda/\lambda)$ becomes an upper support with
backward derivative at most that positive constant.
Continuity of distance and the one-sided mean-value inequality
then give the displayed bound, including cut-locus times.
Fix $\lambda=1$; the error is independent of $i$ and $s_i$.

Put $A_i=d_{-\delta}(q_i,x_i^-)$ and
$B_i=d_{-\delta}(q_i,x_i^+)$.  They lie between $s_i/2$ and
$s_i/2+C_\delta$, while the distance between the endpoints
does not decrease backward.  If $\widehat c_i$ is their earlier
comparison cosine and $c_i'$ the trimmed terminal cosine,
expanding the squares gives
\[
 0\leq\widehat c_i+1\leq c_i'+1
       +\frac{8C_\delta}{s_i}+\frac{8C_\delta^2}{s_i^2}.
\]
Hence $\widehat c_i\to-1$.  Take actual minimizing unit-speed
segments $\alpha_i^-,\alpha_i^+$ in the earlier metric to
these same endpoints.  For fixed $a,b\geq0$, corresponding-side
comparison and the triangle inequality yield, eventually,
\[
 a^2+b^2-2ab\widehat c_i
 \leq d_{-\delta}(\alpha_i^-(a),\alpha_i^+(b))^2
 \leq(a+b)^2.
\]
Distances on either individual segment are exact.  Gluing the
two segments at $q_i$ therefore produces maps whose pairwise
distances tend to the distances on $\mathbb R$.
We spell out why these source points have preimages under the
actual convergence maps $e_i$.  Fix $A>0$.  Completeness of the
limit makes $K=\overline B_{g_\infty}(o,4A)$ compact, so $K$ lies
inside one exhaustion stage.  Metric convergence gives
$|v|_{g_\infty}\leq2|de_i(v)|_{H_i(-\delta)}$ on $K$ for large $i$.
The image of $B(o,4A)$ is open and the image of $K$ is closed.
If a source path from $q_i$ of length less than $A$ first left
that open image, its inverse up to the first exit would have
length less than $2A$.  The inverse endpoint would still belong
to $B(o,4A)$, a contradiction.  Thus the whole source $A$-ball
is eventually covered, with inverse points in $K$.

Distances also converge for these moving preimages.  On any
fixed larger compact ball the norm comparison improves to
factors tending to one.  A source path between two points of a
fixed source ball, with length within one of their distance,
stays in another fixed source ball by the triangle inequality.
The preceding first-exit argument confines its inverse to a
fixed compact limit ball.  Its inverse length gives the lower
distance comparison.  Conversely, completeness supplies a
minimizing limit segment between the preimages; their compact
confinement bounds the radius of that segment, and its image
length gives the upper comparison.  Both errors tend to zero.

Define the glued source arc at $t\geq0$ by $\alpha_i^+(t)$ and
at $t<0$ by $\alpha_i^-(-t)$ whenever the parameter lies within
the segment; use $q_i$ elsewhere.  Pull it back where the image
contains it, and use $o$ elsewhere.  At zero the pullback is
always $o$.  For every fixed pair of parameters these arbitrary
extensions disappear eventually, and the preceding distance
comparison applies.
Choose one free ultrafilter containing the cofinite sets.
Compactness gives limits for all parameter values along this
same filter; continuity of distance gives
$d_{g_\infty}(\gamma(a),\gamma(b))=|a-b|$ and $\gamma(0)=o$.
This is the required minimizing line.

The initial-direction comparison is the argument of
\cite[Theorem 5.35, pp. 100--102]{MT2007}; compare
\cite[Proposition 41.13, p. 2678, and Appendix G, p. 2852]{KleinerLott2013}.
The time buffer and compact-image passage above retain the supplied
smooth limit throughout.  No volume assumption enters this lemma.
\end{proof}

\begin{lemma}[Strong positivity from lower heat contacts]
\label{lem:ancient-heat-contacts}
Let $U$ be a connected open manifold domain and let $g(t)$
be a smooth family of Riemannian metrics for $a\leq t\leq b$,
where $a<b$.  Suppose $v:U\times[a,b]\to[0,\infty)$ is
continuous and has the following contact property at every
$(x,t)$ with $a<t<b$.  Whenever a smooth spatial function
$\psi$ touches $v(\cdot,t)$ from below at $x$, there is a
differentiable time function $\tau$, defined near $t$, with
\[
 \tau(t)=v(x,t),\qquad v(x,s)\leq\tau(s)\ \hbox{near }t,
 \qquad \Delta_{g(t)}\psi(x)\leq\tau'(t).
\]
If $v(p,s)>0$ and $a\leq s<t\leq b$, then $v(q,t)>0$
for every $q\in U$.  Consequently a terminal zero forces
$v=0$ on the whole closed slab, and at every time $t>a$
the spatial zero set is either empty or all of $U$.
\end{lemma}
\begin{proof}
First compare $v$ with a smooth heat subsolution $w$ on a
compact spatial domain $D\Subset U$, assuming $w\leq v$
initially and on its lateral boundary.  If the comparison
failed before the terminal time, then for sufficiently small
$\varepsilon>0$ the continuous function
$w-v-\varepsilon(t-a)$ would have a positive maximum on a
compact truncated cylinder.  The maximum is spatially interior
and occurs after $a$.  Subtracting a constant from the spatial
function $w(\cdot,t)$ gives a lower contact to $v$ there.
For its upper time support $\tau$, the function
$w(x,s)-\tau(s)-\varepsilon(s-a)$ has a maximum relative
to the past, so its derivative is nonnegative.  But
\[
 \partial_tw-\tau'-\varepsilon
 \leq\Delta w-\tau'-\varepsilon\leq-\varepsilon,
\]
a contradiction.  Continuity includes the terminal time.
Only the contact property of $v$ has been used.

To propagate a positive value within one chart, move a small
coordinate ball along a smooth center path $\gamma(t)$ lying
in the interior of a compact chart domain.  Continuity gives
$v\geq\varepsilon_0>0$ on an initial ball.  Compactness of
the path gives positive clearance from the lateral boundary.
Choose $\rho>0$ below both radii and set
\[
 s=\rho^2-|z-\gamma(t)|^2,\qquad
 w(z,t)=\varepsilon_0e^{-C(t-a)}
 \begin{cases}e^{-1/s},&s>0,\\0,&s\leq0.\end{cases}
\]
The zero extension is smooth through the moving boundary,
with all profile derivatives zero there.  Write the actual
coordinate Laplacian as
$\Delta=a^{ij}\partial_i\partial_j+\beta^i\partial_i$,
where $a=(g_{ij})^{-1}$ and
$\beta^k=-g^{ij}\Gamma^k_{ij}$.  For $y=z-\gamma(t)$ put
\[
 Q=a^{ij}y_i y_j,\qquad
 H=\operatorname{tr}a+\beta\cdot y+y\cdot\gamma'(t).
\]
On the compact spacetime domain, positivity and continuity
give constants $0<\ell\leq\Lambda$ and $H_0\geq0$ with
$\ell|y|^2\leq Q\leq\Lambda|y|^2$ and $H\leq H_0$.
Inside its support, direct differentiation yields
\[
 \frac{(\Delta-\partial_t)w}{w}
       =C+\frac{4Q}{s^4}-\frac{8Q}{s^3}-\frac{2H}{s^2}.
\]
One uniform damping constant suffices even as $s\downarrow0$.
Choose
\[
 d=\min\left\{\frac{\rho^2}{2},\frac14,
                 \frac{\ell\rho^2}{2(H_0+1)}\right\},\qquad
 C=\frac{8\Lambda\rho^2}{d^3}+\frac{2H_0}{d^2}.
\]
For $s\leq d$, we have $Q\geq\ell\rho^2/2$,
$4Q-8Qs\geq2Q$ and $2H_0s^2\leq\ell\rho^2$.
Thus the last three terms have nonnegative sum.
For $s\geq d$, the two negative terms are bounded by $C$,
while $4Q/s^4\geq0$.  The residual is therefore nonnegative
everywhere, including the zero-extension boundary.
Initially $0\leq w\leq\varepsilon_0$ on its support and
it vanishes elsewhere; it vanishes on the lateral boundary
throughout.  Compact comparison gives $v\geq w$.
The final center value
$\varepsilon_0e^{-C(b-a)}e^{-1/\rho^2}$ is strictly positive.

Any two points of $U$ can be joined by a path in $U$.
Cover its compact image by finitely many coordinate balls
whose closed inverse images lie in $U$, and subdivide so
the endpoints of each piece lie in one such ball.  Replace
each piece by its straight coordinate segment.  Convexity
keeps the segment in that ball, and its compact image has
positive lateral clearance.  Assign each piece an equal
positive fraction of the prescribed elapsed time, including
constant pieces.  Repeated application of the moving bump
propagates positivity to the chosen endpoint at that time.
No completeness or global coefficient bound is required.

If a terminal value is zero, positivity propagation excludes
positive values at every strictly earlier point; nonnegativity
then gives zero there, and continuity includes the terminal
slice.  Restricting to $[a,t]$ proves the spatial zero-set
alternative for every $t>a$.
This realizes the local propagation underlying
\cite[Theorem 4.16, p.~70, and Theorem 4.18, p.~71]{MT2007}
by explicit barriers for continuous contact supersolutions.
\end{proof}

\begin{theorem}[Classification retaining the given flow]
\label{thm:three-shrinkers}
An admissible three-dimensional static shrinking soliton
$(N,h,f)$ with $\lambda=1/2$ generates a shrinking Ricci flow $G(t)$,
$t<0$, with $G(-1)=h$.
Every supplied shrinking flow associated to these data is, by fixed
identifications at all negative times, a compact round flow, a round
two-sphere times a line, or a quotient of that product by a free
isometric involution.
Every three-dimensional asymptotic limit from
Theorem~\ref{thm:ancient-asymptotic} has bounded curvature on each
slice and such a presentation on its original limit carrier.
\lean{PoincareMT.threeDimensionalAncientAndShrinkingSolitonClassification}
\uses{thm:surface-ancient}
\source{MT2007}
\dcref{Theorem 9.42, pp. 206--208; Proposition 9.46, pp. 209--213;
Claim 9.51, p. 215; Corollaries 9.53--9.54, pp. 217--218}
\end{theorem}
\begin{proof}[Proof for admissible static data]
\emph{Generating the flow.}
The coordinate bootstrap used for surfaces also applies to
$\operatorname{Hess}f=h/2-\operatorname{Ric}$, whose right side
is smooth; thus the given $C^2$ potential is smooth.
Bounded curvature gives a constant $C\geq0$ such that
$|\operatorname{Hess}f(v,v)|\leq C|v|^2$.
Along a local gradient trajectory $\gamma$, set
$e(s)=|\nabla f|^2(\gamma(s))$.  Then
\[
 e'(s)=2\operatorname{Hess}f(\nabla f,\nabla f),\qquad
 |e'(s)|\leq2Ce(s),\qquad
 e(s)\leq e(0)e^{2C|s|}.
\]
On a fixed finite time interval this bounds the trajectory's speed
and length.  It stays in a bounded closed ball, which is compact by
completeness.  The local existence and continuation theorem for smooth
vector fields therefore extends the trajectory in both time directions.
Uniqueness gives a smooth action $\Phi_s$ with inverse $\Phi_{-s}$.
Define
\[
          G(t)=(-t)\Phi_{-\log(-t)}^*h,\qquad t<0.
\]
Using $\mathcal L_{\nabla f}h=2\operatorname{Hess}f$,
\[
 \partial_tG(t)
    =\Phi_{-\log(-t)}^*(-h+2\operatorname{Hess}f)
    =-2\operatorname{Ric}_{G(t)}.
\]
It has $G(-1)=h$, and every slice is isometric to $(-t)h$.
Thus completeness, curvature sign and metric noncollapsing are
preserved.  If $|\operatorname{Rm}_h|\leq B$, then
$|\operatorname{Rm}_{G(t)}|\leq B/(-t)$, uniformly on every
backward interval $t\leq b<0$.
The classification argument now applies to any supplied shrinking
flow with these same slice homotheties, not just to this construction.

\emph{Ricci positivity in the compact case.}
At a maximum $p$ of $f$, its Hessian is nonpositive, so the soliton
equation gives $\operatorname{Ric}_h(p)\geq h(p)/2>0$.
We explain why a Ricci kernel elsewhere is impossible.
For a nonnegatively sectionally curved flow on a connected slab
$[a,b]$, let
\[
 \mu_k(x,t)=\min_{(e_1,\ldots,e_k)\ \text{orthonormal}}
                      \sum_{i=1}^k\operatorname{Ric}_{g(t)}(e_i,e_i).
\]
This is the sum of the $k$ smallest Ricci eigenvalues.  It is
nonnegative and continuous, including at the time endpoints:
in local orthonormal frames the metric and Ricci matrices are
continuous, and minimizing over the compact set of orthonormal
$k$-frames preserves continuity.

Each $\mu_k$ satisfies the scalar heat supersolution inequality in
the sense of lower contacts.  At a minimizing frame, extend its
vectors spatially by radial parallel transport.  Their first
covariant derivatives and diagonal second covariant derivatives
vanish at the center.  To obtain these jets, solve the linear
parallel-transport ODE along each coordinate ray, using the
full connection coefficients.  Uniqueness under radial
reparametrization gives $\nabla_uY(su)=0$.  Evaluating at zero
and differentiating along the ray give the first and diagonal
second identities; no vanishing of mixed second derivatives
is asserted.  Metric compatibility keeps the transported
frame orthonormal on one common neighborhood.  The tensor
Leibniz rule then cancels every frame derivative in the
Laplacian at the center.  The extended frame trace is an upper support
for $\mu_k$, with Laplacian equal to the corresponding trace of
$\Delta\operatorname{Ric}$.  In time use transport
$\partial_t e_i=\operatorname{Ric}^{\sharp}e_i$, which preserves
orthonormality under $\partial_tg=-2\operatorname{Ric}$.
At a fixed spatial point this is a linear ODE on one fixed
tangent space.  Its continuous coefficient is bounded on the
closed slab, so its fundamental solution $U(s)$ exists and
is invertible throughout.  For a contact time $t$, use
$e_i(s)=U(s)U(t)^{-1}e_i$.  Pairing two such solutions gives
\[
 \frac d{ds}g_s(e_i(s),e_j(s))
 =-2\operatorname{Ric}(e_i,e_j)
   +\operatorname{Ric}(e_i,e_j)+\operatorname{Ric}(e_i,e_j)=0.
\]
Only the interior contact times require ordinary time
derivatives; continuity retains the frame at both endpoints.
The two frame-derivative terms cancel the quadratic covariant-Ricci
terms in its evolution.  The remaining reaction on a vector $v$ is
\[
 2\sum_j\lambda_j\operatorname{Rm}(v,u_j,v,u_j)\geq0,
\]
where $u_j$ diagonalizes Ricci with eigenvalues $\lambda_j\geq0$.
For any smooth spatial function touching $\mu_k$ from below, its
Laplacian is at most that of the upper frame support.  The temporal
frame support has derivative at least this latter Laplacian.
These are the heat-contact inequalities needed for the scalar
strong maximum principle of Lemma~\ref{lem:ancient-heat-contacts},
without differentiating the eigenvalues.

It follows that at every $t>a$ the zero set of $\mu_k(\cdot,t)$
is either empty or the whole connected manifold.  Indeed, a positive
value at time $t$ persists to a slightly earlier time by continuity,
and strong positivity propagates it to every point at time $t$.
Since $\mu_k=0$ exactly when the Ricci kernel has dimension at least
$k$, the kernel dimension is spatially constant on every such slice.
Apply this on $[-2,-1]$ to the supplied shrinking flow.  At time $-1$
the positive Ricci form at $p$ has zero-dimensional kernel, so Ricci
is positive everywhere.  Homothety gives the same positivity at
every negative time.

\emph{The stationary Ricci quotient.}
Write $H=|\operatorname{Ric}|^2$ and $q=H/R^2$.
Under multiplication of the metric by $b>0$, the two inverse metric
contractions make $H$ scale by $b^{-2}$, while $R^2$ has the same
factor.  Thus self-similarity makes the spatial maximum of $q$
independent of time.  Fix $t<0$ and a maximizing point $x$.
For every nearby time $s$, $q(s,x)\leq\max_Nq(s,\cdot)=q(t,x)$.
Hence $\partial_tq(t,x)=0$, as well as
$\nabla q(t,x)=0$ and $\Delta q(t,x)\leq0$.
No differentiable choice of maximizing points is needed.

In an orthonormal frame put
\[
 Q=\sum_{i,j,a,b}\operatorname{Ric}_{ij}
                  \operatorname{Rm}_{iajb}\operatorname{Ric}_{ab}.
\]
Taking the two metric traces of the Ricci evolution, including the
derivatives of the inverse metrics, gives
\[
 (\partial_t-\Delta)H=-2|\nabla\operatorname{Ric}|^2+4Q,
 \qquad (\partial_t-\Delta)R=2H.
\]
The quotient rule and $H=qR^2$ give
\[
 (\partial_t-\Delta)q
   =\frac4R\langle\nabla q,\nabla R\rangle
    +\frac{-2|\nabla\operatorname{Ric}|^2
           +2q|\nabla R|^2+4Q-4H^2/R}{R^2}.
\]
At $(x,t)$, $\nabla H=2H\nabla R/R$.
Cauchy--Schwarz in
$\nabla H=2\langle\operatorname{Ric},\nabla\operatorname{Ric}\rangle$
therefore gives
$q|\nabla R|^2\leq|\nabla\operatorname{Ric}|^2$.
The stationary maximum implies
\[
                  0\leq\frac4{R^3}(RQ-H^2).
\]
For the Ricci eigenvalues $a\geq b\geq c>0$ at this point, write
$S=a+b+c=R$, $H=a^2+b^2+c^2$, and $F=H-S^2/3$.
The three-dimensional curvature identity gives
\[
 Q=\tfrac52SH-\tfrac12S^3-2(a^3+b^3+c^3).
\]
For example, this follows by contracting
\[
 R_{iajb}=\delta_{ij}\operatorname{Ric}_{ab}
  +\delta_{ab}\operatorname{Ric}_{ij}
  -\delta_{ib}\operatorname{Ric}_{aj}
  -\delta_{aj}\operatorname{Ric}_{ib}
  -\tfrac S2(\delta_{ij}\delta_{ab}-\delta_{ib}\delta_{aj})
\]
against $\operatorname{Ric}_{ij}\operatorname{Ric}_{ab}$.
Consequently the quartic polynomial $P=2H^2-2SQ$ is nonpositive
at the maximum, whereas
\[
 P=2H^2+S^4-5S^2H+4S(a^3+b^3+c^3)
   =\sum_{\rm cyc}(a-b)^2(a+b-c)^2 .
\]
The needed coercive bound can be seen without dividing by an
eigenvalue gap.  Put $u=a-b\geq0$ and $v=b-c\geq0$; expansion gives
\[
 P-3c^2F=2u^2(u^2+3uv+2uc+3v^2+4vc)\geq0,
 \qquad F=\frac{(a-b)^2+(a-c)^2+(b-c)^2}{3}\geq0.
\]
Thus $P\geq3c^2F\geq0$, forcing $F=0$ at the maximum, where $q=1/3$.
Since $q\geq1/3$ everywhere, $\operatorname{Ric}_{G(t)}=(R/3)G(t)$
everywhere on this slice.
Contracted Bianchi gives $(1/3-1/2)dR=0$.
The three-dimensional curvature identity now gives constant
sectional curvature $R/6>0$.

\emph{The noncompact positive-curvature alternative.}
Suppose the carrier is noncompact and Ricci is strictly positive.
We establish the estimates at infinity before considering the level
geometry.  Fix $p\in M$ and a unit-speed minimizing geodesic
$\gamma:[0,L]\to M$ from $p$.  A uniform curvature bound gives
$0\leq\operatorname{Ric}(\gamma',\gamma')\leq B$.
Apply the index inequality to parallel normal fields multiplied
by a cutoff which grows linearly from zero to one on the first
unit interval, is one in the middle, and decreases linearly on
the last unit interval.  The derivative integrals and discarded
end curvature terms are uniformly bounded.  Short intervals are
controlled directly by $B$; in particular
\[
 \int_0^L\operatorname{Ric}(\gamma',\gamma')\,ds\leq C,
 \qquad C=6+4B.
\]
The soliton equation gives $(f\circ\gamma)''=1/2-
\operatorname{Ric}(\gamma',\gamma')$.  Integrating on initial
subsegments and then once more yields
\[
 (f\circ\gamma)'(s)\geq s/2-|\nabla f(p)|-C,
 \qquad f(x)\geq f(p)+\tfrac14d(p,x)^2-C'd(p,x),
 \quad C'=|\nabla f(p)|+C.
\]
Every sublevel of $f$ is thus compact by completeness.  The
potential attains its minimum and is unbounded above, since $M$
is noncompact.  Connectedness makes every level above the minimum
nonempty.  The soliton identities are
\[
 dR=2\operatorname{Ric}(\nabla f,\cdot),\qquad
 R+|\nabla f|^2-f=C_0
\]
for a constant $C_0$.  The scalar upper bound implies that
$q=|\nabla f|^2$ tends uniformly to infinity on high superlevels.
High levels are therefore compact regular surfaces; we have not
assumed they are connected.

Along the complete gradient flow $\Psi_s$ constructed above,
\[
 \frac{d}{ds}f(\Psi_sx)=q(\Psi_sx),\qquad
 \frac{d}{ds}R(\Psi_sx)
       =2\operatorname{Ric}(\nabla f,\nabla f)(\Psi_sx)\geq0.
\]
Choose a nonempty compact core $\{f\leq b\}$ with $q\geq1$
on $\{f\geq b\}$.  Every trajectory reaches the core in backward
time: otherwise $f(x)-f(\Psi_sx)\geq-s$ for all $s<0$,
contradicting $f(\Psi_sx)>b$ when $s$ is sufficiently negative.
Strict Ricci positivity gives a positive minimum $c$ of $R$
on the core; scalar monotonicity propagates $R\geq c>0$ globally.

\emph{Scalar normalization at infinity.}
Translate the generated flow by $F(t)=G(t-1)$, $t<1$, so $F(0)=h$.
Homothety gives static noncollapse with the original constant and
\[
 |\mathrm{Rm}_{F(t)}|\leq\frac{B_0}{1-t},\qquad
 R_{F(t)}\geq\frac{c}{1-t}.
\]
For any sequence $x_i$ escaping compact sets, pointed compactness
gives one unscaled limit $H(t)$ on all $t<1$.
Choose its normal atlas on a reference time window containing
zero.  On each larger compact time window the displayed curvature
bound, derivative estimates and metric comparison bound all jets
in these same charts and keep the metric uniformly positive.
Diagonal extraction retains a single flow on the entire interval.
It is complete and nonnegatively curved.
Lemma~\ref{lem:rf-limit-static-noncollapse}, applied separately
at each interior time, preserves its static noncollapse constant.
Scalar curvature is a continuous expression in the inverse metric
and its first two spatial derivatives; hence the retained embeddings
pass the lower bound $c/(1-t)$ at every limit point.

Every fixed slice of this same limit contains a line through its
base.  The centers escape in each complete source slice, since a
bounded closed ball there is compact and has bounded distance in
the original metric.  In that slice put $d_i=d(p,x_i)$ and
$r_i=\sqrt{d_i+1}$; then $r_i\to\infty$ and $r_i/d_i\to0$.
Apply Lemma~\ref{lem:ancient-geometric-selected-line} to this
source slice with $Q_i=1$ and $\delta=0$.  Then
$D_i=d_i\to\infty$ and $L_i=r_i\to\infty$, and the already
retained time-slice embeddings supply its convergence hypothesis.
The resulting line lies in the already chosen limit slice.

For a selected time $T<1$, Lemma~\ref{lem:ancient-line-busemann}
gives a unit parallel Busemann gradient.  Translate $T$ to zero.
The bound $B_0/(1-T)$ holds throughout the past, so
Lemma~\ref{lem:ancient-parallel-persistence} preserves this gradient
there.  The fixed-level construction in the proof of
Lemma~\ref{lem:ancient-fixed-factor} gives its actual ancient
surface factor.  Product-ball comparison transfers static
noncollapse with constant $\kappa/2$; ambient scalar positivity
makes the surface nonflat.  Surface classification makes it
compact and round.  Thus $R_H(T,\cdot)$ is spatially constant
and $|\operatorname{Ric}_H|^2=R_H^2/2$ at every point of this slice.
This construction is performed separately at each $T$; no common
factor through future times is required.

The scalar equation on the actual ambient flow now gives $R'=R^2$,
so $R(t)^{-1}=A-t$.  Positivity for all $t<1$ forces $A\geq1$.
The inherited bound gives $c(A-t)\leq1-t$; letting $t\uparrow1$
forces $A\leq1$.  Consequently
\[
                 R_H(t)=\frac1{1-t},\qquad R_H(0)=1.
\]
Basepoint scalar convergence proves that every escaping sequence
in $(M,h)$ has a subsequence with scalar curvature tending to one.

Start any gradient trajectory above a threshold on whose entire
superlevel $q\geq1$.  Its potential grows at least linearly, so
its integer points escape every compact ball.  Strict Ricci
positivity makes $R(\Psi_sx)$ strictly increasing for $s\geq0$.
The preceding subsequence limit implies $R(\Psi_1x)\leq1$,
and hence $R(x)<R(\Psi_1x)\leq1$.  This establishes $R<1$
on a common high superlevel, not merely at selected centers.

\emph{Expansion and curvature of the high levels.}
Raise the threshold so both $q>1$ and $R<1$.
The normalized gradient $X=\nabla f/q$ satisfies $Xf=1$;
on every compact regular slab its flow identifies the levels and
their individual components.  For unit $v$, complete $v$ to an
orthonormal frame.  The difference $R/2-\operatorname{Ric}(v,v)$
is the complementary sectional curvature, and is nonnegative.
The soliton equation therefore gives
\[
             \operatorname{Hess}f\geq\tfrac{1-R}{2}h>0.
\]
For vectors tangent to a level, derivatives of $1/q$ disappear
from the Lie derivative since $df$ vanishes on those vectors.
Thus
\[
 ({\cal L}_Xh)|_{T\{f=a\}}
      =\frac{h-2\operatorname{Ric}}q\bigg|_{T\{f=a\}}\geq0.
\]
The flow expands each induced metric and its area.  Equivalently,
put $N=\nabla f/\sqrt q$ and $b=\operatorname{Ric}(N,N)$.
First variation gives
\[
 \frac{d}{da}\operatorname{Area}(\{f=a\})
       =\int_{\{f=a\}}\frac{1-R+b}{q}\,d\mu\geq0.
\]
The same pointwise calculation applies to each transported component.

Let $\mathrm{II}$ be the second fundamental form of a level.
Tracing the soliton equation tangentially gives mean curvature
$(1-R+b)/\sqrt q$.  The Gauss equation and $|\mathrm{II}|^2\geq0$
give the sufficient estimate
\[
 R_{\{f=a\}}=R-2b+(\operatorname{tr}\mathrm{II})^2-|\mathrm{II}|^2
       \leq R-2b+\frac{(1-R+b)^2}{q}<1.
\]
Indeed, for $A=1-R+b$, nonnegative Ricci curvature gives
$0\leq b\leq R$, hence $0<A\leq1<q$ and $A^2/q<A$.
The right side is strictly smaller than
$R-2b+A=1-b\leq1$.

\emph{The area contradiction.}
To select all centers on one complete trajectory, first extend the
normalized gradient smoothly across the low region.  Choose a smooth
transition $\theta$ which is zero on $(-\infty,0]$ and one on
$[1,\infty)$, with $0\leq\theta\leq1$, and put
\[
 d(q)=\theta(2q-1)q+1-\theta(2q-1),\qquad
 \widehat X=\frac{\nabla f}{d(q)}.
\]
For $q\geq0$ we have $d(q)\geq\max(q,1/2)$, and $d(q)=q$
when $q\geq1$.  The field has bounded speed:
$|\widehat X|=\sqrt q/d(q)\leq3$, using $\sqrt q\leq q+1$.
Completeness confines every finite-time trajectory to a compact
ball, so local flow continuation gives a global smooth flow
$\Phi_s$ with $\Phi_{s+t}=\Phi_s\Phi_t$.
Along it $0\leq(f\circ\Phi)'\leq1$.  If both $f(x)>a$ and
$f(x)+s>a$, where $q>1$ on $\{f>a\}$, the entire intervening
trajectory stays above $a$.  For $s\geq0$ this follows from
monotonicity; for $s<0$ it follows from
$f(\Phi_sx)\geq f(x)+s$.  There the derivative is exactly one,
so
\[
                 f(\Phi_sx)=f(x)+s.
\]
In particular $\Phi_{d-c}$ gives a diffeomorphism from level $c$
to level $d$ whenever $c,d>a$, with inverse $\Phi_{c-d}$.

Fix $x$ with $f(x)>a$ and choose centers $x_k=\Phi_kx$.
They escape compact sets because $f(x_k)=f(x)+k$.
Take an unscaled pointed limit with embeddings $E_i$, base $o$,
and source centers $x_{\nu_i}$, where $\nu_i$ increases strictly.
Set
\[
 a_i=|\nabla f(x_{\nu_i})|\longrightarrow\infty,
 \qquad b_i=\frac{f\circ E_i-f(x_{\nu_i})}{a_i}.
\]
These are functions on the actual exhaustion domains of the limit.
We next prove their smooth compact convergence along a common
subsequence, rather than infer their derivatives from a limit of
the ambient metrics alone.

A uniform Hessian bound, integrated along a minimizing geodesic
from a center $z$, gives
\[
 |f(y)-f(z)|\leq\bigl(|\nabla f(z)|+C d(z,y)\bigr)d(z,y).
\]
After division by $|\nabla f(z)|\geq1$, this bounds the normalized
values on each fixed-radius ball.  Subtracting the two Hamilton
conservation identities also gives
\[
 |\nabla f(y)|^2=|\nabla f(z)|^2+f(y)-f(z)+R(z)-R(y).
\]
The bounded scalar curvature and the value estimate therefore bound
the normalized gradients on those balls.  Every compact subset of
the complete limit lies in a fixed limit ball.  Its points can be
joined to $o$ by curves of bounded length in a larger compact ball;
the eventual bound $|dE_i(v)|\leq2|v|$ bounds their image lengths.
Thus its image lies within a fixed radius of $x_{\nu_i}$, and the
source bounds give uniform values and first derivatives for $b_i$
on every compact limit set.

Here is the higher derivative step.  In any fixed limit chart let
$g_i(t)$ be the pulled-back spacetime metric coefficients.  At
time zero the coordinate Hessian equation is
\[
 D^2b_i=Db_i\circ\Gamma(g_i(0))+
              \frac{g_i(0)+\partial_tg_i(0)}{2a_i}.
\]
The numerator is $g_i(0)-2\operatorname{Ric}_{g_i(0)}$ by the
actual Ricci flow equation.  Smooth spacetime convergence bounds
every derivative of this numerator.  Uniform positive definiteness
and metric jet convergence bound every derivative of the Christoffel
coefficients.  Differentiating the displayed equation $m$ times
bounds $D^{m+2}b_i$ by a finite Leibniz sum of coefficient derivatives
and derivatives of $b_i$ of order at most $m+1$, plus the forcing
derivative.  Induction starting with the value and gradient bounds
therefore bounds all jets on compact chart sets.
Arzela--Ascoli on a countable chart exhaustion, with one diagonal
subsequence for all derivative orders, produces a smooth function
$b$ with smooth compact convergence and $b(o)=0$.
Relabel the retained embeddings and the increasing source indices
along this common subsequence.
Passing through the same equation gives $\operatorname{Hess}b=0$,
since $a_i\to\infty$.  At $o$, the identity
$|\nabla b_i|_{g_i(0)}^2=1$ passes through the inverse metric and
first derivatives.  Hence $|\nabla b(o)|=1$; the identity
$d|\nabla b|^2=2\operatorname{Hess}b(\nabla b,\cdot)=0$
and connectedness give $|\nabla b|=1$ everywhere.

The complete unit parallel gradient gives a product isometry
\[
 e:(\Sigma\times\mathbb R,h_\Sigma+ds^2)
                 \longrightarrow (M_\infty,H(0)),
 \qquad \Sigma=b^{-1}(0),\quad b(e(y,s))=s.
\]
Indeed its flow raises $b$ at unit speed, so flowing backwards by
$b(z)$ gives the inverse projection onto the zero level; parallelism
preserves the metric and orthogonality.  The factor is connected
and complete.  Its scalar curvature equals the ambient scalar,
which has already been proved to be one.  In dimension two this
means $\operatorname{Ric}_{h_\Sigma}=h_\Sigma/2$, so Myers makes
this particular factor compact.  No identification with an earlier
choice of Busemann factor is needed.

On the fixed compact slab $\Sigma\times[-1,1]$, put
$B_i(s,y)=b_i(e(y,s))$.  Eventually the whole slab belongs to the
embedding domain.  Compact chart convergence, a finite chart cover
and metric norm comparison give uniformly
\[
 B_i(s,y)-s\to0,\qquad \partial_sB_i(s,y)\to1,
 \qquad |d_yB_i(s,y)|_{h_\Sigma}\to0.
\]
For large $i$, the endpoint values have opposite signs and
$\partial_sB_i\geq1/2$.  The intermediate value theorem and
strict monotonicity give a unique root $s=u_i(y)\in(-1,1)$.
The implicit function theorem makes $u_i$ smooth, and differentiation
of the root equation gives the quantitative bounds
\[
 |u_i(y)|\leq2|B_i(0,y)|,
 \qquad du_i=-\frac{d_yB_i(u_i(y),y)}{\partial_sB_i(u_i(y),y)}.
\]
Thus $u_i\to0$ uniformly with its first derivatives.
If $e(y_0,0)=o$, exact centering and uniqueness give $u_i(y_0)=0$.

The maps $\psi_i(y)=E_i(e(y,u_i(y)))$ are injective immersions
into the level $f=f(x_{\nu_i})$, and send $y_0$ to $x_{\nu_i}$.
The graph differential is injective because its first component
is the identity; the other two maps have invertible differentials.
Their induced metrics $h_i=\psi_i^*h$ are the actual level metrics.
In the limiting product the graph metric is exactly
$h_\Sigma+du_i\otimes du_i$.  If the ambient relative metric
error is at most $\varepsilon$ and $|du_i|\leq\eta$, then
\[
 |h_i(v,v)-h_\Sigma(v,v)|
    \leq\bigl(\varepsilon(1+\eta^2)+\eta^2\bigr)
                       h_\Sigma(v,v).
\]
Consequently the relative error tends to zero.  On a surface,
$(1-\delta)h_\Sigma\leq h_i\leq(1+\delta)h_\Sigma$
gives the same factors for area, by taking square roots of the
two-dimensional determinants.  Compactness makes the total areas
finite, so
$\operatorname{Area}(h_i)\to\operatorname{Area}(h_\Sigma)$.

Each $\psi_i(\Sigma)$ is an entire connected component of its
regular level.  It is open in that level by the inverse function
theorem in dimension two, closed by compactness, and connected.
For $i\leq j$ the flow time $\nu_j-\nu_i$ equals the difference
of their potential values and sends $x_{\nu_i}$ to $x_{\nu_j}$.
It therefore sends their entire level components onto each other.
The injective parametrizations identify this restriction with a
diffeomorphism of $\Sigma$.  The tangential metric expansion
proved above makes its pullback of $h_j$ at least $h_i$.
Area is invariant under this diffeomorphism, whence
\[
 \operatorname{Area}(h_0)\leq\operatorname{Area}(h_i)
             \leq\lim_j\operatorname{Area}(h_j)
             =\operatorname{Area}(h_\Sigma).
\]

But $R_{h_0}<1$ everywhere by the level Gauss estimate.
The continuous positive function $1-R_{h_0}$ has positive integral
on the nonempty compact surface.  Gauss--Bonnet, applied to the
two metrics on this same surface, gives
\[
 \operatorname{Area}(h_0)>
 \int_\Sigma R_{h_0}\,d\mu_{h_0}
 =\int_\Sigma R_{h_\Sigma}\,d\mu_{h_\Sigma}
 =\operatorname{Area}(h_\Sigma),
\]
contradicting the preceding area inequality.  Metric independence
of total scalar curvature holds for nonorientable surfaces as well,
so no sphere identification or prescribed value of the total
curvature is needed.
Hence Ricci has a null direction somewhere.

\emph{The null line and its global coordinate.}
Write $h=G(-1)$ and first work with the ancient flow
$F(s)=G(s-1)$, $s\leq0$.
A Ricci-null vector gives a null sectional plane, since its Ricci
curvature is a sum of nonnegative sectional curvatures.
We explain also why a null plane suffices.  Extend its orthonormal
vectors by radial parallel transport at time zero.  Its curvature
has a spatial minimum zero, so the rough curvature Laplacian on
that plane is nonnegative.  With the vectors fixed, curvature is
nonnegative for earlier times and zero at zero; its left time
derivative is nonpositive.  The curvature evolution therefore
makes the reaction on the plane nonpositive.  Polarization of the
nonnegative radial curvature forms kills every term obtained by
replacing one of the four entries of this null plane.  Thus the
fixed-vector Ricci corrections vanish, and this reaction is the
moving-orthonormal-frame reaction $2(A^2+\operatorname{adj}A)$.

In the cyclic wedge basis of an orthonormal frame the curvature
operator has the form
\[
 A=\begin{pmatrix}0&0&0\\0&a&b\\0&b&c\end{pmatrix},
 \qquad 2(ac-b^2)\leq0.
\]
The zero row follows from nonnegativity and polarization.
The corresponding two-by-two Ricci block is
$\left(\begin{smallmatrix}c&-b\\-b&a\end{smallmatrix}\right)$.
If $a=0$, its second basis vector is Ricci null.  If $a>0$,
the vector with coordinates $(a,b)$ has Ricci square
$a(ac-b^2)\leq0$.  Nonnegativity of Ricci and polarization make
this nonzero vector a member of the bilinear Ricci kernel.
The localized partial-trace maximum principle gives spatially
constant nullity on each bounded slab.  The derivative-annihilation
and transport argument in
Lemma~\ref{lem:ancient-parallel-persistence} gives backward
inclusion of these kernels.  In dimension three, two
independent Ricci-null vectors force all sectional curvatures to
vanish: planes containing either vector are null, and curvature
polarization removes the mixed components.  Nonflatness at some
point therefore makes the terminal kernel exactly one-dimensional
everywhere.

Here is a local construction of its smooth sections.  In a fixed
tangent trivialization, raise one index of Ricci using the auxiliary
Euclidean inner product, giving a smooth symmetric matrix $B(y)$.
Let $K=\ker B(p)$ and let $P$ be Euclidean orthogonal projection
onto $K$.  The matrix $B(p)+P$ is invertible: it is the identity
on $K$ and is invertible on $K^\perp$.  It remains invertible near
$p$.  The restriction of $P$ to $\ker B(y)$ is injective, since
its kernel would lie in $\ker(B(y)+P)$; constant nullity makes it
an isomorphism onto $K$.  Consequently, for $v\in K$,
\[
                 W(y)=(B(y)+P)^{-1}v
\]
belongs to $\ker B(y)$ and satisfies $W(p)=v$.
Normalize a nonzero such section using $h$ to obtain a smooth
unit null field $V$ on a neighborhood.

Its parallelism follows from an energy identity, not just from
the rank assertion.  Differentiate $\operatorname{Ric}(V,w)=0$:
\[
 (\nabla_u\operatorname{Ric})(V,w)
               =-\operatorname{Ric}(\nabla_uV,w).
\]
Differentiating once more and contracting gives at the terminal
time
\[
 (\Delta\operatorname{Ric})(V,V)
       =2\sum_i\operatorname{Ric}(\nabla_{e_i}V,\nabla_{e_i}V).
\]
The fixed-vector time derivative of $\operatorname{Ric}(V,V)$
is nonpositive there.  Its reaction vanishes, since nonnegative
sectional curvature makes every curvature tensor entry containing
a Ricci-null vector zero.  Ricci evolution now forces the displayed
sum of nonnegative terms to vanish.  Each $\nabla_uV$ lies in the
one-dimensional kernel.  Differentiating $h(V,V)=1$ shows that it
is also orthogonal to $V$, hence $\nabla V=0$.

Form the actual unit-kernel cover
\[
 \widehat M=\{(x,v):v\in\ker\operatorname{Ric}_h(x),\ h(v,v)=1\},
 \qquad \pi(x,v)=x.
\]
The local unit fields identify this space over a neighborhood
with its product with $\{+1,-1\}$; the sign is the inner product
with the chosen field.  These charts give a smooth two-sheeted
cover, equipped with $\pi^*h$.  Its free sheet reversal is
$\rho(x,v)=(x,-v)$.  The function
\[
                         r(x,v)=2\,df_x(v)
\]
is global.  On a sheet with local unit field $V$, the soliton
equation and parallelism give
\[
 dr(u)=2\nabla^2f(u,V)=h(u,V),
 \qquad |\nabla r|=1,\qquad \nabla^2r=0,
 \qquad r\circ\rho=-r.
\]

The lifted metric is complete.  Indeed a finite-sheeted cover
is proper: locally the projection is the closed projection from
a product with a finite set, and its fibers are compact.
A Cauchy sequence upstairs projects to a convergent Cauchy
sequence downstairs.  The compact set consisting of that sequence
and its limit has compact inverse image, which gives a cluster
point upstairs; the original Cauchy sequence then converges.
The same argument applies to every lifted time slice, and closed
connected components are complete.  Each component $C$ projects
onto $M$, by lifting paths from one point in its image.

The unit parallel field $\nabla r$ on $C$ has a complete flow
$\Theta_z$.  Along it, $r$ increases by $z$.  Its zero level
$\Sigma$ is therefore nonempty, and
\[
 e:\Sigma\times\mathbb R\longrightarrow C,\quad
 e(y,z)=\Theta_z(y),\qquad
 e^{-1}(q)=(\Theta_{-r(q)}q,r(q))
\]
is a diffeomorphism.  Parallelism makes this an isometry for
$h_\Sigma+dz^2$.  In particular $\Sigma$ is connected and complete.
This explicit coordinate supplies the global real line factor.

\emph{The surface factor and the supplied flow.}
Metric noncollapsing lifts with the same constant $\kappa$.
A ball upstairs maps onto the corresponding ball downstairs:
lift curves from its center whose lengths approach the distance.
Local isometry preserves curvature and length, so a curvature
bound on the upstairs ball yields the bound needed downstairs.
Projection is distance nonincreasing; the calibrated Hausdorff
volume of its image is at most that of the upstairs ball.  Thus
the downstairs lower volume bound also holds upstairs, and the
same applies within $C$.

Backward persistence of the unit parallel gradient on bounded
slabs, proved in Lemma~\ref{lem:ancient-parallel-persistence}, identifies
every slice $F(s)$, $s\leq0$, by this same map $e$ with a product
$h_\Sigma(s)+dz^2$.  The actual zero-level flow is complete,
has nonnegative curvature and bounded curvature, and is nonflat
on every slice.  For the last assertion, backward scalar positivity
from bounded-slab Harnack gives positive scalar somewhere at every
earlier time; the product formula transfers it to the surface.
The ball inclusion into a product ball times an interval of length
$2a$ gives the surface volume bound $(\kappa/2)a^2$ whenever the
required radius-$a$ parabolic curvature bound holds.  Hence this
is an ancient $\kappa/2$-solution.  Surface classification makes
$\Sigma$ compact and round.

The zero level is totally geodesic, so restriction of the original
potential gives
\[
 \operatorname{Ric}_{h_\Sigma}+\nabla^2_\Sigma(f\circ\pi)
                         =\tfrac12h_\Sigma.
\]
If $K$ is its constant sectional curvature, the Hessian is
$(1/2-K)h_\Sigma$.  At a maximum and a minimum of the restricted
potential its signs force $K=1/2$.  Thus the original scalar
curvature is one, using the product identity and surjectivity
of $C\to M$.

It remains to include $-1<t<0$ and to retain the \emph{supplied}
flow $G$, rather than choosing a new realization of the soliton.
Its slice homotheties give scalar curvature $-1/t$ and Ricci
nullity one for all $t<0$.  On any compact negative-time slab,
backward kernel inclusion and equal dimension identify the actual
kernel subspaces at all times.  The fixed field $V=\nabla r$ at
time $-1$ is consequently Ricci null throughout the slab.

To see that it stays parallel, the null-section energy argument
at each interior time makes every $\nabla_uV$ null.  Differentiating
the vanishing curvature entries then shows that $\nabla\operatorname{Rm}$
annihilates $V$ in each curvature slot.  Contraction and the second
Bianchi identity imply that $\nabla\operatorname{Ric}$ annihilates
$V$ also in its derivative slot.  In the connection variation
formula
\[
 g_t((\partial_t\nabla^t)_uV,w)
 =-(\nabla_u\operatorname{Ric})(V,w)
  -(\nabla_V\operatorname{Ric})(u,w)
  +(\nabla_w\operatorname{Ric})(u,V)
\]
every term therefore vanishes.  Continuity includes the slab
endpoints, so $\nabla^tV=0$.  The metric evolution also gives
$\partial_t g_t(V,w)=-2\operatorname{Ric}_t(V,w)=0$.
Its value is $dr(w)$ at time $-1$, hence $V$ remains the same
unit gradient of the same function $r$.

In dimension three a unit Ricci-null vector determines
\[
 \operatorname{Ric}_t=\frac{R_t}{2}(g_t-dr\otimes dr).
\]
For fixed tangent vectors the quantity
$k(t)=g_t(u,w)-dr(u)dr(w)$ thus satisfies $k'=k/t$.
Integrating on $t<0$ gives the tensor identity on the fixed cover
\[
 g_t=(-t)(g_{-1}-dr\otimes dr)+dr\otimes dr,
 \qquad e^*g_t=(-t)h_\Sigma+dz^2.
\]
Only the slice homotheties and the actual Ricci-flow equation were
used; a differentiable choice of homotheties is unnecessary.

\emph{Retaining the deck action.}
A compact connected surface of curvature $1/2$ has a smooth
round covering $q:S^2\to\Sigma$ with
$q^*h_\Sigma=2g_{S^2}$.  One construction continues inverse local
round isometries over the simply connected sphere; uniqueness of
continuation removes path dependence.  The resulting local isometry
has open image, and compactness makes its image closed, proving
surjectivity.  This is the constant-curvature covering construction
of \cite[Theorem 1.11 and its space-form consequence, pp. 8--9]{MT2007}.

Its fibers are either singletons or antipodal pairs.  To see this,
two points in a fiber and the derivative of $q$ determine an
ambient orthogonal motion with the same first-order data.
Uniqueness of local isometries from first-order data, continued
along geodesics, shows that it preserves $q$ everywhere.
These motions act freely and their orbits are exactly the fibers.
A nonidentity motion in this subgroup has no fixed unit vector
and therefore determinant $-1$ in dimension three.  The determinant
homomorphism is injective, so the subgroup has at most two elements.
Its nonidentity element squares to the identity; all its eigenvalues
are then $-1$, making it the antipodal map.

There are now two possibilities for the sheet reversal $\rho$.
If $\rho(C)$ is disjoint from $C$, each fiber of $C\to M$ is a
singleton.  This map is a diffeomorphism, giving either the sphere
product or the quotient by $(x,z)\mapsto(-x,z)$.
If $\rho(C)=C$, it restricts on $r=0$ to a free isometry of
$\Sigma$.  The antipodal surface alternative cannot have such
an isometry: every isometry lifts, by the same first-order
construction, to an orthogonal motion of $S^2$; such a motion
in dimension three has an eigenvector with eigenvalue $1$ or $-1$.
Its line is a fixed point in the antipodal quotient.
Consequently $\Sigma$ itself is a sphere in this case.
Conjugate $\rho|_C$ by the fixed product diffeomorphism and sphere
coordinates.  This gives a free involution preserving
$(-2t)g_{S^2}+dz^2$ for every $t<0$, and its orbits are exactly
the fibers over $M$.  The explicit forms of these involutions
are established below.
\end{proof}

\begin{lemma}[A singleton convex exhaustion level]
\label{lem:ancient-singleton-level}
Let $(M,g)$ be a connected, complete, noncompact three-manifold of
strictly positive sectional curvature.  Fix $o\in M$.
For a unit-speed minimizing ray $\gamma$ from $o$, put
\[
 b_\gamma(x)=\lim_{T\to\infty}(d(\gamma(T),x)-T),\qquad
 K_c=\{x:b_\gamma(x)\geq-c\text{ for every such ray}\}.
\]
There are $c\in\mathbb R$ and $p\in M$ with $K_c=\{p\}$.
This singleton is totally convex: every geodesic segment whose
endpoints equal $p$ remains at $p$, without a minimizing assumption.
\source{Petersen2016}
\dcref{Lemma 12.4.8 and the following soul argument, pp. 467--468}
\end{lemma}
\begin{proof}
We first describe the compact convex levels and their exact inner
parallel sets, as in
\cite[Lemma 3.13(i), pp. 17--19]{MaederBaumdickerSeidel2025}.
The triangle inequality makes the approximants defining
$b_\gamma$ decreasing and bounded between $-d(o,x)$ and $d(o,x)$;
their limit is one-Lipschitz.  On a geodesic $\eta$ of constant speed
$v$, Lemma~\ref{lem:rf-squared-distance-comparison} gives concavity of
\[
             d(q,\eta(t))^2-v^2t^2.
\]
Its shortened-segment upper supports include both the distance
cut locus and the point $\eta(t_0)=q$.
Divide by $2T$ after taking $q=\gamma(T)$ and subtracting $T^2$.
For each fixed $t$ the resulting expression equals
\[
 b_{\gamma,T}(\eta(t))+
   \frac{b_{\gamma,T}(\eta(t))^2-v^2t^2}{2T},\qquad
 b_{\gamma,T}(x)=d(\gamma(T),x)-T.
\]
The second term tends to zero.  Passing each concavity inequality
to the limit proves that $b_\gamma$ is concave along every geodesic.
Thus all the closed sets $K_c$ are totally convex.

For $c\geq0$, $K_c$ contains $\overline B(o,c)$ and is compact.
Indeed, if it were noncompact, take points of $K_c$ escaping all
balls and minimizing segments from $o$ to those points.  Total
convexity keeps the segments in $K_c$, and properness gives a
pointwise compact limit along one nonprincipal ultrafilter, as in
the calibrated-ray construction below.  The distance identities
make the limit a ray $\gamma$ from $o$ wholly contained in $K_c$.
But $b_\gamma(\gamma(t))=-t$, contrary to the defining inequality
at $t=c+1$.  For negative $c$, compactness follows by inclusion
in $K_0$.  Also, whenever $c<c'$, the Lipschitz bound gives
\[
                         K_c\subset\operatorname{int}K_{c'}:
\]
a ball of radius $c'-c$ about a point of $K_c$ satisfies every
inequality for $K_{c'}$.  The levels exhaust $M$ because they
contain the corresponding metric balls.

Fix $c_0>0$ and put $K=K_{c_0}$ and
$d_K(x)=\operatorname{dist}(x,M\setminus K)$.
The complement is nonempty by noncompactness of $M$.
For $x\in K$ and each ray $\gamma$ from $o$,
\[
                       d_K(x)\leq c_0+b_\gamma(x).         \tag{S1}
\]
To see this, a coray from $x$ calibrated by $b_\gamma$ reaches,
at parameter $c_0+b_\gamma(x)+\epsilon$, a point outside $K$
at exactly that distance from $x$.  Let $\epsilon\downarrow0$.
Such corays are obtained from compact limits of minimizing segments
to $\gamma(T)$; the construction and calibration squeeze are given
below and use only completeness.
Conversely, if $r\leq c_0+b_\gamma(x)$ for all rays, a point
$y\notin K$ violates at least one defining inequality.
For that ray the Lipschitz bound gives
\[
        d(x,y)\geq b_\gamma(x)-b_\gamma(y)>r.
\]
Taking the infimum over $y$ proves $r\leq d_K(x)$.
Together these two inequalities give the exact identity, for $r\geq0$,
\[
                 \{x\in K:d_K(x)\geq r\}=K_{c_0-r}.       \tag{S2}
\]

The continuous function $d_K$ attains a maximum $R>0$ on $K$;
indeed $d_K(o)\geq c_0$.
The maximal-depth set $S=\{x\in K:d_K(x)=R\}$ is nonempty,
compact, and by (S2) equals $K_{c_0-R}$.
It has empty ambient interior.  If a ball of radius $\epsilon$
about $x$ lay in $S$, choose
$0<\delta<\min\{\epsilon,R\}$ and a point $y\notin K$ with
$d(x,y)<R+\delta/2$.
On a minimizing segment to $y$, the point $z$ at distance $\delta$
from $x$ would still belong to $S$, whereas
\[
                  d_K(z)\leq d(z,y)<R-\delta/2.
\]
This contradiction proves the empty-interior assertion.

Strict positive curvature now supplies a uniform local improvement.
Fix $q\in M$.  We need one curvature bound for all nearby tangent
planes.  Choose a local trivialization
$K_y:T_qM\to T_yM$ with $K_q$ the identity.  On the compact set
$|v|_q,|w|_q\leq2$ consider
\[
 \begin{split}
 F(y,v,w)={}&\operatorname{Rm}_y(K_yv,K_yw,K_yv,K_yw)\\
 &+\bigl(|K_yv|_y^2-1\bigr)^2
   +\bigl(|K_yw|_y^2-1\bigr)^2
   +\langle K_yv,K_yw\rangle_y^2 .
 \end{split}
\]
At $q$ the curvature term is nonnegative.  If all three
normalization penalties vanish, the pair is orthonormal and
strict sectional positivity makes the curvature term positive.
Otherwise a penalty is positive.  Thus $F(q,v,w)$ has a
positive minimum $m$ on this compact set.  Continuity and a
finite subcover give $F(y,v,w)>m/2$ for all these pairs and
all $y$ in one neighborhood of $q$.
Shrink the neighborhood so that $K_y^{-1}$ has norm less than
two relative to the actual metrics.  The inverse images of
every orthonormal pair at $y$ then lie in the compact set,
and their three penalties vanish.  Hence every such plane
has curvature at least $\kappa_0=m/2>0$.
For arbitrary pairs, a change to an orthonormal basis of their
plane multiplies both the curvature numerator and Gram
determinant by the same determinant squared.  Dependent pairs
give zero on both sides.  Thus this is also the tensor bound
$\operatorname{Rm}(v,w,v,w)\geq
\kappa_0(|v|^2|w|^2-\langle v,w\rangle^2)$ needed below.
An intrinsic closed ball about $q$ lies in the neighborhood.
Shrink its radius so that every $x$ in a smaller ball has a whole
closed ball $\overline B(x,\rho)$ inside this curvature neighborhood,
with the same $\rho>0$ for all these $x$.
For a distant source $y$, let $d=d(y,x)$.
Move one quarter of the way from $y$ to $x$ along a minimizing
segment.  The shortened segment has length $C=3d/4$ and supplies
a smooth local inverse exponential branch at its endpoint,
by the reversal argument in the proof of
Lemma~\ref{lem:rf-squared-distance-comparison}.
The function $u$ equal to $d/4$ plus the branch radius is a smooth
touching upper support for $d(y,\cdot)$, and $|\nabla u(x)|=1$.

The second variation retains the curvature integral rather than
only its sign.  On the shortened segment use the affine parallel
field $Z(t)=(t/C)P_t w$ with the prescribed transverse endpoint
vector $w$.  If $J$ is the radial Jacobi field with the same boundary
values, integration of its equation makes the index cross term
$I(J,Z-J)$ vanish.  The minimizing segment gives
$I(Z-J,Z-J)\geq0$, so $I(J,J)\leq I(Z,Z)$.
The last $\rho$ units of the segment lie in the positive-curvature
ball.  Elsewhere curvature is nonnegative.  Integrating the kinetic
term and this retained terminal curvature term gives
\[
 \nabla^2u_x(w,w)\leq
 \left(\frac1C-\kappa_0 C
             \frac{1-(1-\rho/C)^3}{3}\right)
                       \bigl(|w|^2-du_x(w)^2\bigr).        \tag{S3}
\]
For a general endpoint vector the same expression follows by
separating its radial component, on which the radius Hessian
vanishes.  Equivalently, differentiate the square of the radius
and use Gauss' identity; the subtracted square in (S3) is essential.
If $C\geq\rho$ and $C\geq6/(\kappa_0\rho)$, then
\[
 \frac1C-\kappa_0 C\frac{1-(1-\rho/C)^3}{3}
                                      \leq-\kappa_0\rho/6.
\]
Indeed $1-(1-a)^3\geq a$ for $0\leq a\leq1$, and
$1/C\leq\kappa_0\rho/6$.
Choose $\theta=\min\{\kappa_0\rho/6,1\}>0$.
For all sufficiently distant sources, uniformly over endpoints
in the smaller ball, the support therefore satisfies
\[
               \nabla^2u\leq-\theta(g-du\otimes du)
                      \quad\hbox{at its touching point}.   \tag{S4}
\]

Apply this to $y=\gamma(T)$ for any ray from the fixed point $o$.
The required lower bound on $T$ is uniform in that ray, since
$d(\gamma(T),q)\geq T-d(o,q)$.
Normalize the support by $u_0=u-T$ and apply the increasing
concave transform $H(s)=1-e^{-s}$.  The chain rule gives
\[
 \nabla^2H(u_0)=e^{-u_0}
             (\nabla^2u_0-du_0\otimes du_0)
                              \leq-\theta e^{-u_0}g.       \tag{S5}
\]
The inequality uses $\theta\leq1$ to control the remaining
negative differential square.
At a touching point in a ball of radius $r$ about $q$,
$u_0=d(\gamma(T),x)-T\leq d(o,q)+r=:B$.
Thus the same constant $\delta=\theta e^{-B}/2>0$ works for
every ray and every touching point in that ball.
Along any unit-speed geodesic segment contained there, the
upper-support criterion yields concavity of
\[
                  H(b_{\gamma,T}(\eta(t)))+\delta t^2.
\]
For clarity, this criterion follows by contradiction from an
interior minimum below a chord.  For a continuous function $w$ with
the indicated upper supports and a chord $l$ on $[a,b]$, a failure
of concavity gives a negative interior value of $w-l$.
For sufficiently small $\epsilon>0$,
$w(t)-l(t)-\epsilon(t-a)(t-b)$ still has a negative minimum in
$(a,b)$ and is zero at the endpoints.  A touching upper support
at that minimum would have strictly negative second derivative
despite having a local minimum.
Taking $T\to\infty$ in each concavity inequality proves that
\[
                     H(b_\gamma(\eta(t)))+\delta t^2       \tag{S6}
\]
is concave, with the same positive $\delta$ for all rays.

If $S=K_{c_0-R}$ contained distinct points, join them by a
unit-speed minimizing segment, which stays in $S$ by total convexity.
Choose a nontrivial short subinterval $[a,b]$ about its midpoint,
contained in a ball where (S6) holds.
Write $c=c_0-R$ and $m=(a+b)/2$.
At each endpoint $b_\gamma\geq-c$, so (S6) gives, for every ray,
\[
 H(b_\gamma(\eta(m)))\geq1-e^c+\delta(b-a)^2/4.
\]
This is a common strictly positive gap.
Continuity of the exponential permits a single $\epsilon>0$ with
$1-e^{c-\epsilon}<1-e^c+\delta(b-a)^2/4$.
Monotonicity of $H$ then puts $\eta(m)$ in $K_{c-\epsilon}$,
hence in the ambient interior of $K_c=S$.
This contradicts the empty interior of $S$.
Thus $S$ is a singleton; its total convexity has been retained
throughout the construction.
\end{proof}

The singleton level also supplies directions pointing away from its
point.  We construct a complete flow in those directions before using
the topology of the ambient manifold.  The following argument develops
the flow construction discussed after Corollary 2.10 in
\cite[p. 27]{MT2007}.

\begin{lemma}[Outward flow and radial coordinates]
\label{lem:ancient-soul-coordinates}
Let $(M,g)$ be a complete connected Riemannian three-manifold with
nonnegative sectional curvature.  Suppose that one of the all-ray
levels above is $K_c=\{p\}$.  There is a smooth complete flow $\Phi$
fixing $p$ such that, for every $x\ne p$,
\[
             t\longmapsto d(p,\Phi_t(x))
\]
is strictly increasing and has image $(0,\infty)$.
The manifold is diffeomorphic to $\mathbb R^3$, with $p$ corresponding
to the origin.  There is also a homeomorphism
\[
 H:S^2\times(0,\infty)\longrightarrow M\setminus\{p\},
 \qquad d(p,H(\theta,r))=r.
\]
These conclusions concern the same point $p$.
\end{lemma}
\begin{proof}
Write $r(y)=d(y,p)$.  An inward unit direction at $y\ne p$ means the
initial tangent of any unit-speed minimizing segment from $y$ to $p$.
There may be several such directions.  Every inequality below holds
for all of them.

\emph{A common local potential.}
Fix $q\ne p$.  Since $q\notin K_c$, some defining ray has
$b_\gamma(q)<-c\leq b_\gamma(p)$.
Put $\delta=(b_\gamma(p)-b_\gamma(q))/2>0$.
For arbitrarily distant $z=\gamma(T)$ the defining limits give
\[
 d(z,p)-d(z,q)>\delta,
 \qquad d(z,q)>r(q)^2/(2\delta).
\]
Consequently
\[
 A(q):=d(z,p)^2-d(z,q)^2-r(q)^2>0:                         \tag{F1}
\]
factor the difference of squares and use
$d(z,p)+d(z,q)>2d(z,q)$.
In particular $z\ne q$.

A shortened-segment upper support for $d(z,\cdot)$, constructed
in the proof of Lemma~\ref{lem:rf-squared-distance-comparison},
is smooth near $q$,
equals $d(z,q)$ at $q$, and majorizes distance there.
It can be obtained by moving $z$ a positive distance along a selected
minimizing segment to $q$ and using its smooth radial inverse branch;
the triangle inequality gives the majorization.
Square this positive local support and extend it smoothly by a cutoff.
This gives a globally smooth function $\rho$ with
\[
 \rho(q)=d(z,q)^2,\qquad \rho(y)\geq d(z,y)^2
 \quad\hbox{near }q.
\]
Set $A(y)=d(z,p)^2-d(z,y)^2-r(y)^2$ and
$a=A(q)/(2r(q))>0$.
Choose $0<\eta\leq1$ such that throughout $B(q,\eta)$ the preceding
majorization holds and
\[
                   A(y)/r(y)>a,\qquad r(y)>r(q)/2.
\]
On the compact ball $\overline B(q,1)$ choose $B\geq0$ with
$|\nabla^2\rho(w,w)|\leq B|w|^2$.
Fix a single positive length
\[
 s=\min\{\eta/4,r(q)/4,a/[4(B+1)]\},
 \qquad Bs\leq a/4,
\]
and shrink to the open neighborhood
\[
 U=B(q,\eta/4)\cap
       \{y:\rho(y)-d(z,y)^2<as/4\}.
\]
It contains $q$ and omits $p$.

For $y\in U$ take any unit-speed minimizing segment
$\alpha:[0,L]\to M$ from $y$ to $p$, where $L=r(y)$.
Its first $s$ units stay in $B(q,\eta)$.
Nonnegative curvature gives concavity of
$d(z,\alpha(t))^2-t^2$ as in the preceding proof.
The chord between $0$ and $L$ therefore gives
\[
 d(z,\alpha(s))^2-s^2-d(z,y)^2
       \geq sA(y)/L>as.
\]
The majorization and the bound defining $U$ imply
\[
                \rho(\alpha(s))-\rho(y)>3as/4.
\]
On the other hand, integrating the Hessian bound twice along this
unit-speed geodesic gives
\[
 \rho(\alpha(s))-\rho(y)
       \leq s\,d\rho_y(\dot\alpha(0))+Bs^2/2.
\]
Since $Bs\leq a/4$, these inequalities force
$d\rho_y(\dot\alpha(0))\geq a/2$.
Thus a single smooth potential and a single positive constant
$\kappa=a/2$ work for every inward direction at every point of $U$.
No continuity of a choice of minimizing segment was used.

\emph{Gluing and completion.}
The field $-\kappa^{-1}\nabla\rho$ on $U$ pairs with every inward
unit direction by at most $-1$.
For any prescribed $R_0>0$, cover the complement of $B(p,R_0)$ by
these neighborhoods, and add $B(p,R_0)$ with the zero field.
A smooth locally finite partition of unity gives a global field $V$
with
\[
            \langle V(y),v\rangle\leq-1
               \quad(r(y)\geq R_0)
\]
for every inward $v$.
Indeed, at such a point only fields satisfying this affine inequality
have nonzero weights, and their weights sum to one.

In a normal neighborhood of $p$, the squared distance $r^2$ is smooth.
Extend it to a globally smooth function $f$ agreeing with $r^2$ on
$B(p,R)$ for some $R>0$.
For every inward unit direction in this ball,
\[
               df_y(v)=-2r(y),\qquad \nabla f(p)=0.
\]
The first identity follows by differentiating
$f(\alpha(t))=(r(y)-t)^2$ for small nonnegative $t$; the second follows
from the local minimum at $p$.
Use the preceding field with $R_0=R/4$ and choose a smooth cutoff
$0\leq\chi\leq1$, equal to one near $\overline B(p,R/4)$ and
supported in $B(p,R/2)$.
Then
\[
 X=\chi\nabla f+(1-\chi)V,\qquad
 a_0(y)=2\chi(y)r(y)+1-\chi(y)
\]
satisfy $X(p)=0$ and
\[
 \langle X(y),v\rangle\leq-a_0(y),\qquad
                   a_0(y)>0\quad(y\ne p).                \tag{F2}
\]
Where the second summand is present, $r(y)\geq R/4$; where the
first is present, the exact normal-ball derivative applies.
The margin $a_0$ is continuous.

Put $Z=X/2$.  To complete this field without changing it near $p$,
choose a compactly supported smooth cutoff $\psi$, equal to one
near $p$, and set
\[
 b=\psi+(1-\psi)(1+|Z|^2)^{-1},\qquad Y=bZ.
\]
Here $0\leq\psi\leq1$, $b>0$, and $b=1$ near $p$.
Off the compact support of $\psi$ the norm of $Y$ is at most one;
on that support it is bounded by compactness.  Thus $Y$ has bounded
speed on the complete manifold.
Every finite piece of an integral curve stays in a compact metric
ball, since its distance traveled is bounded by elapsed time times
this speed bound.  The smooth ordinary differential equation can
therefore be continued past every finite endpoint.
Uniqueness and smooth dependence give a smooth all-time flow
$\Phi_{t+s}=\Phi_t\circ\Phi_s$.
It fixes $p$, and no other orbit reaches $p$, by the inverse
$\Phi_{-t}$.
The inward pairing of $Y$ is bounded above by
$-b(y)a_0(y)/2$, a continuous strictly negative quantity off $p$.
Near $p$ we have the exact identity $Y=\tfrac12\nabla r^2$.

\emph{Distance along the flow.}
We justify monotonicity at cut points as well.
For any smooth endpoint curve $\sigma$ through $y\ne p$ and any
selected minimizing inward direction $v$, first variation gives a
touching upper support $u$ for $r\circ\sigma$ with
\[
                         u'(0)=-\langle v,\dot\sigma(0)\rangle.
\]
One construction varies the endpoint of the selected segment, uses
its energy as an upper support for squared distance, and then takes
the positive square root.  The endpoint term in first variation of
energy is $-2r(y)\langle v,\dot\sigma(0)\rangle$.

On a fixed orbit interval $[a,b]$, the positive continuous function
$y\mapsto b(y)a_0(y)/2$ has a positive minimum $\kappa$
along that orbit segment.
Apply the preceding supports to the reversed curve
$t\mapsto\Phi_{-t}(x)$.  Their derivatives are at most $-\kappa$.
The scalar upper-support mean-value inequality gives
\[
 r(\Phi_b(x))-r(\Phi_a(x))\geq\kappa(b-a).                \tag{F3}
\]
For completeness, that inequality follows directly from compact
minimization.  If a continuous function with touching differentiable
upper supports of derivative at most $C$ increased faster than $C$,
subtract a line of slope $C'>C$ smaller than its endpoint secant
slope.  The resulting function has a minimum at some point before
the right endpoint.  Its touching upper support also has a right
minimum there, yet has negative derivative, a contradiction.
Reversing time in (F3) is essential to its use with upper supports.

Every positive distance is reached on each noncentral orbit.
If a larger distance were never reached, the closure of the forward
orbit would be compact by properness, forward invariant, and would
have a positive maximum of $r$.  Advancing a maximizing point
strictly increases $r$ by (F3), a contradiction.
If a smaller positive distance were never reached, the closure of
the backward orbit would be compact, bounded away from $p$, and
backward invariant.  Moving a minimizing point backward contradicts
its positive minimum of $r$.
Continuity then gives existence of a hitting time for every positive
radius, and strict increase gives uniqueness.

\emph{A global smooth chart.}
Take a normal chart $e:B(0,\varepsilon)\to M$ with $e(0)=p$ small
enough that $Y=\tfrac12\nabla r^2$ on its image.
Gauss's lemma identifies this field with the pushforward of the
Euler field $v\mapsto v$.
Uniqueness of integral curves gives
\[
 \Phi_t(e(v))=e(e^t v)\quad(t\leq0,\ |v|<\varepsilon).
\]
The identity also holds for positive $t$ when both vectors remain
in the ball, by applying the negative-time identity and inverting
the flow.  For arbitrary $v\in\mathbb R^3$ choose $t$ so negative
that $e^t v\in B(0,\varepsilon)$ and define
\[
                       D(v)=\Phi_{-t}(e(e^t v)).          \tag{F4}
\]
The local identity and the action law make this independent of $t$.
Taking one common sufficiently negative time for two vectors proves
injectivity.  Every orbit meets the chart image, because that image
contains a positive metric ball and every orbit reaches every small
positive radius; the fixed orbit at $p$ meets it as well.
If $\Phi_t(x)=e(v)$ in that image, then $D(e^{-t}v)=x$,
which proves surjectivity.
Near any fixed vector, the same choice of $t$ works in (F4), so $D$
is smooth.  Near any fixed $x$ choose a fixed time $t$ taking it into
the chart ball.  The same time works nearby and gives
\[
                         D^{-1}(y)=e^{-t}e^{-1}(\Phi_t(y)).
\]
Thus the inverse is smooth too.  On the original chart ball $D=e$;
in particular $D(0)=p$.

\emph{Actual-distance coordinates.}
Choose $r_*>0$ so small that the entire distance sphere
$\Sigma=\{r=r_*\}$ lies in an orthonormal normal chart.
Gauss's lemma identifies it with $S^2$ by
$\theta\mapsto\exp_p(r_*\theta)$.
To ensure this covers the whole sphere, first choose a normal tangent
ball on which radial lengths equal intrinsic distance, then choose
$r_*$ smaller than both its radius and the radius of a metric ball
contained in its image.

Each orbit meets $\Sigma$ exactly once.  For $x\ne p$ let $\tau(x)$
be the time such that $r(\Phi_{\tau(x)}(x))=r_*$.
This time depends continuously on $x$: for each real $t$,
\[
 \tau(x)>t\ \Longleftrightarrow\ r(\Phi_t(x))<r_*,\qquad
 \tau(x)<t\ \Longleftrightarrow\ r(\Phi_t(x))>r_*.
\]
Both conditions define open sets.
Consequently $(y,t)\mapsto\Phi_t(y)$ is a homeomorphism
$\Sigma\times\mathbb R\to M\setminus\{p\}$, with inverse
$x\mapsto(\Phi_{\tau(x)}(x),-\tau(x))$.
Finally replace time by actual distance along each orbit.
The inverse time $T(y,u)$ for target radius $u>0$ is jointly continuous, since
\[
 T(y,u)>t\ \Longleftrightarrow\ r(\Phi_t(y))<u,
 \qquad
 T(y,u)<t\ \Longleftrightarrow\ r(\Phi_t(y))>u.
\]
Composing this change of parameter with the preceding product
homeomorphism gives $H$ and its exact distance identity.
No smoothness of the global distance spheres or of the hitting-time
function is asserted or needed.
\end{proof}

\begin{lemma}[Ambient sides of a collared sphere]
\label{lem:ancient-collar-sides}
Let $M$ be a manifold homeomorphic to $\mathbb R^3$, and let
$e:S^2\times(-r,r)\to U$ be a homeomorphism onto an open subset,
where $r>0$.  Put $S=e(S^2\times\{0\})$.
Then $M\setminus S$ has exactly two connected open components,
each with boundary $S$, and the two collar halves lie in different
components.  Under any fixed homeomorphism $M\to\mathbb R^3$,
exactly one component has bounded image and compact closure in $M$.
\end{lemma}
\begin{proof}
Define $c(t)=\max\{-1/2,\min\{1/2,t/r\}\}$.
On $U$, take the class of $c(t)$ in $\mathbb R/\mathbb Z$;
off $U$, give this map the constant value $[1/2]$.
The two clamped endpoint values have the same class, and the
nonconstant part is supported in the compact interior collar
$e(S^2\times[-r/2,r/2])$.  Thus this defines a continuous map $Q$
on $M$.  Since $M$ is simply connected and locally path connected,
$Q$ lifts to a continuous real-valued function $F$, normalized to
vanish at one point of $S$.  Existence and uniqueness of the lift
are the covering-space lifting criterion and uniqueness theorem
\cite[Propositions 1.33--1.34, pp. 61--62]{Hatcher}.
Both $F\circ e$ and $c$ lift the same map on the connected collar
and agree at that central point.  Hence
\[
                       F(e(y,t))=c(t).
\]
The zero set of $F$ is exactly $S$: the formula proves this on $U$,
and off $U$ a zero could not project to $[1/2]$.

Let $A=\{F<0\}$ and $B=\{F>0\}$.
They are disjoint nonempty open sets partitioning $M\setminus S$.
Continuity gives $\partial A,\partial B\subset S$, and approaching
each central point through the two collar halves gives the reverse
inclusions.  The sets $A\cap U$ and $B\cap U$ are the corresponding
connected half-collars.
Every connected component $C$ of $A$ meets $U$.
Otherwise $\overline C$ avoids $U$, since $U$ is open.
Any point of $\overline C\setminus A$ would lie in
$\partial A\subset U$, which is impossible.
Thus $\overline C\subset A$.  Components are closed in $A$, so
$C$ is then closed in $M$; local connectedness also makes it open
in $M$.  Connectedness of $M$ would force $C=M$, contrary to
$B\ne\varnothing$.  Since all components of $A$ meet its connected
half-collar, they coincide.  The same proof applies to $B$.

In the chosen Euclidean coordinates place the compact sphere $S$
inside a ball of positive radius $R$.
The closed exterior $\{|x|\geq R\}$ is connected, being the image
of $[R,\infty)\times S^2$ under radial multiplication.
It lies wholly in one member of the open disjoint partition $A,B$.
The other member is bounded.  Both cannot be bounded, since together
with $S$ they cover Euclidean space.  The closure of the bounded
member is compact by Heine--Borel and remains compact on transport
back to $M$.
\end{proof}

We also need the order of two disjoint such spheres $S_1,S_2$.
Write $A_i$ for the bounded side and $B_i$ for the unbounded side,
and assume $A_1\cap A_2\ne\varnothing$.
Connectedness puts $S_1$ entirely in $A_2$ or entirely in $B_2$,
and similarly for $S_2$ relative to the first partition.
If $S_1\subset A_2$, connectedness and unboundedness give
$B_2\subset B_1$: it avoids $S_1$ and cannot lie in the bounded
set $A_1$.  It follows that $\overline{A_1}\subset A_2$.
Indeed a point of $\overline{A_1}\setminus A_2$ belongs to
$\overline{B_2}\subset\overline{B_1}=B_1\cup S_1$;
the first alternative contradicts $B_1\cap\overline{A_1}=\varnothing$
and the second contradicts $S_1\subset A_2$.
The case $S_2\subset A_1$ gives the reverse strict nesting.
The only remaining possibility would be
$S_1\subset B_2$ and $S_2\subset B_1$.
Then the connected set $A_1$ avoids $S_2$ and meets $A_2$, so
$A_1\subset A_2$.  Its boundary $S_1$ would lie in
$\overline{A_2}$, contradicting $S_1\subset B_2$.
Thus the bounded closures are strictly nested.
For two disjoint entire neck carriers, their connectedness also
places each carrier on the same side of the other sphere as its
center.  These are exactly the whole-carrier inclusions needed below.

\begin{lemma}[Antipodal collar obstruction]
\label{lem:ancient-antipodal-collar}
For $a>0$, there is no smooth local diffeomorphism
$F:S^2\times(-a,a)\to\mathbb R^3$ satisfying
$F(-\theta,t)=F(\theta,t)$ everywhere.
\end{lemma}
\begin{proof}
On an annulus about the unit sphere define
\[
               f(x)=F(x/|x|,|x|-1).
\]
The radial coordinate map is a local diffeomorphism, with inverse
$(\theta,t)\mapsto(1+t)\theta$ after restricting to $|t|<1$.
Thus $Df$ would be invertible everywhere near the unit sphere.
But $f(-x)=f(x)$ implies $Df(-x)=-Df(x)$.
In dimension three the continuous function $x\mapsto\det Df(x)$
therefore takes opposite values at antipodal points.
Connectedness of the unit sphere forces it to vanish somewhere,
contradicting invertibility.
\end{proof}

\begin{lemma}[Returning from a buffered cylinder limit]
\label{lem:ancient-terminal-cylinder}
Fix $0<\varepsilon<1/2$.  There is $\eta>0$ with the following
property.  Let $H_i$ be smooth three-dimensional Ricci flows on
open time intervals containing $[-A_i,0]$, with complete slices
and nonnegative curvature operator there.  Suppose $A_i,L_i\to\infty$
and, for basepoints $q_i$,
\[
 \begin{gathered}
 R_{H_i}(0,q_i)=1,\qquad R_{H_i}(u,q_i)\leq1,\\
 R_{H_i}(u,y)\leq4\quad
 (-A_i\leq u\leq0,\ y\in B_{H_i(0)}(q_i,L_i)).
 \end{gathered}
\]
Let $0<\delta\leq\eta$, and suppose the actual shifted flows
$H_i(t-\delta)$ converge smoothly in the pointed sense on $t<\delta$
to $G(t)$, whose time-zero slice is complete.
If that slice is a round spherical cylinder, then the original
terminal metrics $H_i(0)$ contain $\varepsilon$-necks of scale one
centered at $q_i$, for all sufficiently large indices of the
convergence subsequence.
\source{MT2007,KleinerLott2013}
\dcref{Morgan--Tian, Proposition 5.14, pp. 90--91, and Corollary 9.53,
p. 217; Kleiner--Lott, Theorem D.1, p. 2847, and Appendix E, pp. 2848--2849}
\end{lemma}
\begin{proof}
We retain the spatial maps of the given pointed convergence throughout.
All time-estimate constants in this proof are fixed before $\delta$
and the convergence limit are chosen.

\emph{Scalar normalization.}
For large $i$ the controlled cylinder contains the fixed spatial
radius $64(3+8)$ and the time interval $[-2,0]$.
The local curvature-derivative estimate, applied with one unit of
elapsed time, bounds $|\nabla^2\operatorname{Rm}|$ at the basepoint
on $[-1,0]$ by a dimensional constant $C$.
Scalar evolution and nonnegative curvature give
\[
 \partial_u R=\Delta R+2|\operatorname{Ric}|^2
       \leq 3^3 C+32=:D.
\]
Here $|\Delta R|\leq3^3|\nabla^2\operatorname{Rm}|$ and
$|\operatorname{Ric}|^2\leq R^2\leq16$ suffice.
Integration and the assumed basepoint upper bound give
\[
                1-D\delta\leq R_{H_i}(-\delta,q_i)\leq1.
\]
Smooth convergence at the buffered time gives the same inequalities
for $s=R_G(0,o)$.  Choose $\eta\leq\min\{1,1/(2D)\}$ so that
$1/2\leq s\leq1$ and $1-s\leq D\delta$.
A centered parametrization $\Phi:S^2\times\mathbb R\to G(0)$ can
then be chosen so that
\[
                     s\Phi^*G(0)=g_C=2g_{S^2}+dz^2.
\]
This uses curvature $s/2$ on the round surface and rescales the line
coordinate by $s^{-1/2}$; translate that coordinate to put the
basepoint on the central sphere.

\emph{A time bound in the retained coordinates.}
Put $J=[-\varepsilon^{-1},\varepsilon^{-1}]$ and
$d=\lfloor\varepsilon^{-1}\rfloor$.
The compact slab $\Phi(S^2\times J)$ lies in one stage of the
pointed exhaustion.  Compose $\Phi$ with the stored spatial
embeddings to obtain maps $e_i$ on a neighborhood of that slab.
In centered sphere charts write $B_i(u)$ for the coefficient matrix
of $e_i^*H_i(u)$, and $B_C$ for that of $g_C$.
The finite jets of $B_i(-\delta)$ converge to those of $s^{-1}B_C$.
Since $s\in[1/2,1]$, this gives uniform initial ellipticity and
finite-order coefficient bounds depending only on $J$ and $d$.
All bounds can be taken uniformly over the centered sphere charts:
their parameter is the compact unit sphere and the round metric is
rotation invariant.

The images of the slab have bounded source distance at time
$-\delta$ by pointed convergence, hence bounded terminal distance
by $\operatorname{Ric}\geq0$.  For each evaluation point choose a
fixed-radius terminal ball around its image.  Eventually that whole
ball lies inside the radius-$L_i$ controlled ball.
Its radius determines how far along the sequence we must pass,
but not the local derivative constants.
The local curvature estimates therefore bound all covariant
curvature derivatives through order $d$ on $[-1,0]$ at these points,
uniformly in the buffer and the spatial maps.

Here is why these geometric bounds control jets in the same maps.
On $[-\delta,0]$, $0\leq\operatorname{Ric}\leq4H_i$ gives
\[
 e^{-8}B_i(-\delta)\leq B_i(u)\leq e^8 B_i(-\delta).
\]
Thus both coefficients and inverse coefficients stay bounded.
Differentiate $\partial_u B_i=-2\operatorname{Ric}_{e_i^*H_i}$
in the fixed spatial coordinates.  Expressing ordinary derivatives
of Ricci by covariant derivatives introduces Christoffel coefficients
$\Gamma=\tfrac12B_i^{-1}*\partial B_i$.
Once jets below order $j$ are bounded, every term involving the
order-$j$ metric jet is linear in that jet: derivatives of the
inverse use $\partial B_i^{-1}=-B_i^{-1}(\partial B_i)B_i^{-1}$,
and a Christoffel derivative of order $j-1$ has at most one
order-$j$ metric factor.  The curvature derivatives are already
bounded.  Consequently
\[
 \bigl|\partial_u\partial^j B_i\bigr|
                  \leq C_j(1+|\partial^j B_i|).
\]
Gronwall on an interval of length at most one bounds this jet from
its initial value.  Substitution into the displayed inequality and
integration then give a single constant $B\geq0$ with
\[
 |\partial^j B_i(0)-\partial^j B_i(-\delta)|\leq B\delta,
                      \qquad 0\leq j\leq d.               \tag{T1}
\]
This constant depends only on the initial model bounds and the
fixed derivative order, and hence is available before selecting
the buffer.

\emph{The fixed error budget.}
Let $Z\geq0$ bound the model jets through order $d$.
If the buffered coefficient error is at most $\tau$,
the triangle inequality, (T1), and
\[
               |s^{-1}-1|=(1-s)/s\leq2D\delta
\]
give terminal coefficient error at most
$\tau+(B+2DZ)\delta$ through that order.
Converting ordinary coefficient derivatives to cylinder covariant
derivatives on the fixed slab supplies $C_0\geq0$ such that the
squared metric-jet error is at most
\[
                  C_0\bigl(\tau+(B+2DZ)\delta\bigr)^2.    \tag{T2}
\]
Indeed each covariant derivative is a finite sum of ordinary
derivatives with coefficients formed from the fixed cylinder
connection and its derivatives.  Those coefficients and the
tensor-norm comparisons are bounded on the compact slab.
Set
\[
 \tau=\frac{\varepsilon}{2(C_0+1)(B+2DZ+1)}
 \quad\hbox{and require}\quad\eta\leq\tau.
\]
Convergence makes the initial error at most $\tau$ on a sufficiently
late tail.  Since $\delta\leq\tau$, (T2) is at most
$\varepsilon^2/4$; use $C_0\leq(C_0+1)^2$ to see this directly.
This is a strict margin for the required neck closeness.

Finally restrict the actual embeddings $e_i$ to the open slab
$S^2\times(-\varepsilon^{-1},\varepsilon^{-1})$.
They are smooth diffeomorphisms onto their open images, with smooth
inverses supplied by the same convergence embeddings.
Their central spheres contain $q_i$, and the terminal scalar there
is exactly one.  The scale is therefore exactly one, and the metric
and connection are those of $H_i(0)$.
The error estimate proves all the metric conditions of the claimed
necks.  No convergence at the terminal time has been assumed.
\end{proof}

\begin{proof}[Proof for an asymptotic limit]
A global slice bound must precede static classification.
Fix $t_0<0$, write $g=h(t_0)$, and suppose its scalar curvature
$R$ is unbounded.  We describe first the selection needed to produce
small necks on this same slice.

\emph{The null alternative does not need a global curvature bound.}
If a sectional plane at time $t_0$ is null, normalize the metric by
$|t_0|^{-1}$.  Its soliton constant is then $1/2$.
The null-plane calculation, unit Ricci-kernel cover and parallel
coordinate constructed above use the local curvature evolution and
the soliton equation.  Apply this construction on a connected
component of that finite cover.  It is complete and projects onto
the original connected manifold.  The zero level of the parallel
coordinate is a complete nonflat surface, and the restricted
potential satisfies the shrinking equation with constant $1/2$.
At this step use Lemma~\ref{lem:surface-shrinker} directly: it
requires neither bounded curvature nor ancient noncollapsing.
The surface has Gaussian curvature $1/2$, so the product scalar is
one.  Local isometry and surjectivity give the same scalar everywhere
on the normalized original slice.  Undoing the normalization gives
$R\equiv-1/t_0$, contradicting unboundedness.
Nonnegative sectional curvature therefore improves to strict
positivity everywhere on the slice in question.  It is noncompact,
since continuous scalar curvature is bounded on a compact manifold.
Lemmas~\ref{lem:ancient-singleton-level} and
\ref{lem:ancient-soul-coordinates} now identify its underlying manifold
smoothly with $\mathbb R^3$.

\emph{Point selection on a compact ball.}
Fix $p\in M$.  For any $x\ne p$ with $R(x)>0$, put
\[
 d=d_g(p,x),\qquad f=\sqrt R,\qquad L=d f(x)/4,
 \qquad S=\overline B_g(x,d/2).
\]
The function $f$ is bounded on this one compact set.
Start with $q=x$.  Maintain the budget
\[
       d_g(x,q)+2L/f(q)\leq d/2,\qquad f(q)\geq f(x).
                                                               \tag{P1}
\]
If $f(y)\leq2f(q)$ throughout $B_g(q,L/f(q))$, stop.
Otherwise choose a point $z$ of that ball with $f(z)>2f(q)$.
The triangle inequality puts it in $S$, and
\[
 d_g(x,z)+2L/f(z)
 <d_g(x,q)+L/f(q)+L/f(q)\leq d/2.
\]
Thus it can replace $q$ while retaining (P1).
Each replacement more than doubles $f$, so boundedness on $S$
forces termination.
At the final point put $\rho=L/f(q)$.  We have
\[
 \begin{gathered}
 Q:=R(q)\geq R(x),\qquad d_g(p,q)\geq d/2,\qquad \rho>0,\\
 \rho\sqrt Q=d\sqrt{R(x)}/4,\qquad
 B_g(q,\rho)\subset B_g(x,d/2),\qquad
 R\leq4Q\quad\hbox{on }B_g(q,\rho).
 \end{gathered}                                                \tag{P2}
\]
This is the compact-ball doubling argument in
\cite[Corollary 9.38 and Proposition 9.39, pp. 204--205]{MT2007}.

Choose $x_i$ with $R(x_i)>\max\{i,16\}$ and outside
$\overline B_g(p,2(i+1)^2)$.
Such a choice follows by taking $R(x_i)$ larger also than the
maximum of $R$ on that compact ball.
Perform (P2) and then shrink its radius to $r_i=\rho_i/(i+1)$.
With $q_i$ the selected points, $Q_i=R(q_i)$ and
$d_i=d_g(p,q_i)$, the conclusions are
\[
 \begin{gathered}
 d_i\longrightarrow\infty,\quad Q_i\longrightarrow\infty,
 \quad r_i\sqrt{Q_i}\longrightarrow\infty,\quad
 d_i\sqrt{Q_i}\longrightarrow\infty,\\
 r_i/d_i\longrightarrow0,\qquad
 R\leq4Q_i\quad\hbox{on }B_g(q_i,r_i).
 \end{gathered}                                                \tag{P3}
\]
For the less immediate claims, (P2) and $\sqrt{R(x_i)}>4$ give
$\rho_i\sqrt{Q_i}>2(i+1)^2$, while
$\rho_i\leq d_g(p,x_i)/4\leq d_i/2$.
Shrinking by $i+1$ therefore gives both the diverging normalized
radius and the vanishing ratio to the distance from $p$.
Scalar time monotonicity on the retained limit implies, for every
$s\leq t_0$ and $y$ in this same spatial ball,
\[
 R_h(s,y)\leq R_h(t_0,y)\leq4Q_i,
             \qquad |\operatorname{Rm}_{h(s)}|(y)\leq36Q_i.     \tag{P4}
\]
The last inequality uses the dimension-three contraction bound
$|\operatorname{Rm}|\leq9R$ for the adopted norm and nonnegative
curvature operator.  The coarse constant is sufficient here.

In the actual rescaled flows
\[
                    H_i(u)=Q_i h(t_0+u/Q_i)
\]
the terminal base scalar is one, and scalar is at most four on
the terminal radius-$r_i\sqrt{Q_i}$ ball at every $u\leq0$.
These expanding controlled cylinders, rather than a global bound
on the original slice, are the inputs to the buffered compactness
and line-splitting argument.  Here the noncollapse needed for
compactness also precedes the global curvature bound.
On any radius-$r$ ball with terminal tensor bound $r^{-2}$,
scalar contraction, time monotonicity and the reverse tensor
comparison bound the past tensor by $81r^{-2}$.
The concentric radius-$r/9$ ball therefore satisfies the required
parabolic curvature bound throughout its radius-squared time interval.
The original parabolic noncollapse and volume inclusion give the
static constant $\kappa/729$ on the larger ball.
This property is unchanged by the rescaling.
In particular, the terminal radius-$1/6$ controlled ball in $H_i$
has volume at least $(\kappa/729)6^{-3}$, since its curvature is at
most $36$.  The terminal unit ball has the same lower bound.

Choose the tolerance $\eta$ from
Lemma~\ref{lem:ancient-terminal-cylinder} before forming the limit.
Buffered compactness supplies a shift $0<\delta<\eta$ and an
actual pointed limit of $H_i(t-\delta)$ on $t<\delta$.
The limit is complete, has nonnegative curvature operator and a
global tensor bound $36$, since every fixed controlled radius is
eventually present.  The scalar estimate in that lemma gives
$1/2\leq R_G(0,o)\leq1$, so it is nonflat.
Lemma~\ref{lem:rf-limit-static-noncollapse} passes static
$\kappa/729$ noncollapse to this same complete limit at each
time.  A parabolic curvature bound includes its terminal
static bound, so the restriction to $t\leq0$ has the same
parabolic noncollapse constant.
Nonnegative Ricci and the scalar bound four give
$\operatorname{Ric}_{H_i}\leq4H_i$ on the controlled terminal
balls throughout the buffer.  Apply
Lemma~\ref{lem:ancient-geometric-selected-line} with $\Lambda=4$
and the data (P3) to obtain a line in this same $G(0)$.
Splitting by that line and persistence in the bounded ancient flow
give one product $\Sigma\times\mathbb R$ for all $t\leq0$.
The product-ball volume estimate gives the surface noncollapse
constant $\kappa/1458$.
Theorem~\ref{thm:surface-ancient} makes $\Sigma$ compact and round.

\emph{Excluding the antipodal factor on the actual limit.}
The round-surface covering dichotomy developed above gives a smooth
local isometry $q:S^2\to\Sigma$ after scalar normalization.
Either $q$ is a diffeomorphism or its fibers are exactly antipodal
pairs.  In the latter case the product parametrization gives a
local diffeomorphism $\Phi:S^2\times\mathbb R\to G(0)$ satisfying
$\Phi(-\theta,z)=\Phi(\theta,z)$.
Its compact slab $\Phi(S^2\times[-1,1])$ lies in one exhaustion
stage of the supplied pointed convergence.  A sufficiently late
stored embedding maps that slab into the original manifold $M$.
Compose it with the Euclidean diffeomorphism of $(M,g)$ obtained
before point selection.  On the open slab the composite is a local
diffeomorphism to $\mathbb R^3$ and still identifies antipodal pairs.
This contradicts Lemma~\ref{lem:ancient-antipodal-collar}.
Thus the factor is spherical.  Only the supplied convergence map
on one compact slab was needed for this exclusion.

Normalize its cylinder coordinates by the actual base scalar
$s=R_G(0,o)$ and center them at $o$, as in
Lemma~\ref{lem:ancient-terminal-cylinder}.
That lemma gives scale-one $\varepsilon$-necks in the retained
$H_i(0)$ for all sufficiently large indices of its subsequence.
Rescaling back gives necks in the original fixed metric $g$ with
scale exactly $Q_i^{-1/2}$, which tends to zero.
In particular, for any prescribed $\rho>0$ choose an index with
$Q_i>\rho^{-2}$; the resulting neck has scale less than $\rho$.
This preserves the needed quantifier order: choose the fixed neck
accuracy, then its buffer tolerance, then the limit, and only then
the sufficiently late index for the requested small scale.

The geometric obstruction is the following.  There exists
$\varepsilon_*>0$ such that for every complete strictly positively
curved three-manifold and every $0<\varepsilon\leq\varepsilon_*$,
its $\varepsilon$-neck scales have a positive lower bound depending
on the metric.  A scale-$s$ neck has, after multiplication by $s^{-2}$,
the prescribed $C^{\lfloor1/\varepsilon\rfloor}$ closeness to a
scalar-normalized round cylinder
\[
 g_C=2g_{S^2}+dz^2,\qquad |z|<\varepsilon^{-1}.
\]
Here $g_{S^2}$ has sectional curvature one and the neck center has
scalar curvature $s^{-2}$.  In the compact case the maximum of scalar
curvature therefore bounds every neck scale below.

We give the flux part of the noncompact argument, specifying the
geometric inputs that precede it.  A point soul supplies a point $p$
and a ray $\gamma$ from $p$.  Each sufficiently accurate neck sphere
separates a relatively compact open side $A$ containing $p$ from an
unbounded side.  Its two half-necks lie on the corresponding sides.
If scales have no positive lower bound, choose a neck $N_1$ and then
a neck $N_2$ of scale $s_2<s_1/8$ whose whole carrier avoids the
closures of both $N_1$ and $A_1$.  The separation argument gives
\[
 \overline{A_1}\subset A_2,\qquad N_1\subset A_2,
 \qquad \overline{A_1}\cap N_2=\varnothing.
\]
In these inclusions $N_i$ denotes the neck carrier.  Common containment
of the soul orders the two bounded sides; disjointness of the second
carrier from the first bounded side rules out the reverse order.
Here are the metric reasons for the two additional assertions.
An axial path followed by a spherical path gives
\[
 d(x,\text{center}(N))\leq
                  (2\pi+2\varepsilon^{-1})s
                         \quad(x\in N).
\]
At fixed $\varepsilon$, a sequence of scales tending to zero therefore
has carrier diameters tending to zero.  Its centers escape every
compact set, since their scalar curvatures equal $s^{-2}$.
Given a compact set $K$, choose a compact closed $\delta$-neighborhood
of $K$ with $\delta>0$.  Eventually the center lies outside this
neighborhood and the displayed radius is less than $\delta$;
then the whole carrier avoids $K$.  Completeness also makes the
closure of each fixed neck compact, since that closure lies in the
corresponding closed metric ball.  We may thus take
$K=\overline{N_1}\cup\overline{A_1}$ when choosing $N_2$.

To identify which side contains $p$, the point-soul construction
supplies a homeomorphism $S^2\times(0,\infty)\to M\setminus\{p\}$
whose second coordinate is exactly $d(p,\cdot)$.
Radial distance consequently has no local
maximum away from $p$: increase the radius slightly in this
homeomorphism.  Suppose that a relatively compact open side $A$
omits $p$, and that its frontier has diameter at most $D$.
For $x\in A$, the maximum of radial distance on $\overline A$
is attained at some $z\in\partial A$, since an interior maximizer
would be a local maximum away from $p$.
A minimizing segment from $p$ to $x$ meets $\partial A$ at a point
$y$, and therefore
\[
 d(x,y)=d(p,x)-d(p,y)
       \leq d(p,z)-d(p,y)\leq d(y,z)\leq D.
\]
For a neck central sphere we can take $D=2\pi s$.
On the other hand, a point on either half-neck slice
$z=\pm(2\varepsilon)^{-1}$ has distance greater than $2\pi s$
from every point of the central sphere, once $\varepsilon$ is
below a universal threshold.  To obtain an ambient distance estimate,
fix a smooth bump $q$ equal to one near zero and supported in
$(-1/4,1/4)$, with $|q'|\leq C_D$ for a constant $C_D\geq1$.
For $\tau\in\{-1,1\}$ set
$q_{\varepsilon,\tau}(z)=q(\varepsilon z-\tau/2)$.
It vanishes at zero, equals one at $\tau/(2\varepsilon)$,
is supported in $|z|\leq3/(4\varepsilon)$, and has derivative
bounded by $C_D\varepsilon$.  Extend its axial pullback by zero
off the neck.  The support is a compact interior slab, so this
extension is smooth.  The axial covector estimate gives the global
differential bound
\[
 |d q_{\varepsilon,\tau}(v)|\leq
       \frac{C_D\varepsilon}{s\sqrt{1-\varepsilon}}|v|_g.
\]
Integration along every piecewise smooth path, followed by the
length infimum, gives the same Lipschitz constant for ambient distance.
For a point $x$ on the indicated half-neck slice and $y$ on the
central sphere, it follows that
\[
 d(x,y)\geq\frac{s\sqrt{1-\varepsilon}}{C_D\varepsilon}
                                                  >2\pi s
\]
when $\varepsilon\leq\min\{1/4,(8\pi C_D)^{-1}\}$.
The half-neck contained in $A$ supplies such a deep point,
a contradiction.  Thus $p\in A$.
The ambient complementary regions and their order are given by
Lemma~\ref{lem:ancient-collar-sides} and its following argument.
The remaining point-soul input supplies the common Euclidean and
radial parametrizations used here.

Use the convention
\[
 b(x)=\lim_{T\to\infty}\bigl(d(\gamma(T),x)-T\bigr).
\]
The approximants decrease with $T$, lie between
$-d(p,x)$ and $d(p,x)$, and are one-Lipschitz.  Consequently the limit
exists and is one-Lipschitz.  The weak sign needed here is
$\Delta b\leq0$.  To see it, fix a nonnegative smooth compactly
supported test $\eta$ and a ball $B(p,R)$ containing its support.
For $T>R+H$, distance to $\gamma(T)$ is at least $H$ on that support.
Weak Laplacian comparison in dimension three gives
\[
 \int_M \bigl(d(\gamma(T),x)-T\bigr)\Delta\eta\,d\mu
       \leq \frac2H\int_M\eta\,d\mu.
\]
Subtracting $T$ is harmless since $\int\Delta\eta=0$.
The compact bound $d(p,x)|\Delta\eta|$ permits dominated convergence
as $T\to\infty$; next let $H\to\infty$.
The weak Green identity for a locally Lipschitz function then yields
\[
        \int_M db(\nabla\eta)\,d\mu\geq0.                 \tag{B}
\]
Differentiability almost everywhere follows from the Lipschitz bound,
and $|\nabla b|\leq1$ there.

Choose a smooth increasing profile $\phi$ with values in $[0,1]$,
equal to zero for $z\leq-1$ and one for $z\geq1$.
Thus $\phi'\geq0$ and $\int_{\mathbb R}\phi'=1$.
For each neck form its outward transition $T_i$: it is zero on the
compact side away from the neck, one on the unbounded side away
from the neck, and either $\phi(z)$ or $1-\phi(z)$ on the neck,
according to which half-neck is compact.  Constancy near the two
ends makes this a smooth global function.  Its gradient is supported
in the compact slab $|z|\leq1$, and $1-T_i$ has compact support.
The ordering above gives $T_1\geq T_2$: on $N_1$ the second transition
is zero, on $N_2$ the first is one, and away from both necks this is
the inclusion of their compact sides.  Apply (B) to $T_1-T_2$.
Writing the finite integrals as
\[
 F_i=\int_{N_i}db(\nabla T_i)\,d\mu,
                  \qquad F_1-F_2\geq0,                    \tag{C}
\]
retains the sign of the negative outward flux.

Let $A_C>0$ be the area of a cross section of $g_C$.
The zeroth-order neck comparison supplies, in its coordinates,
\[
 s^2(1-\varepsilon)g_C\leq g\leq
 s^2(1+\varepsilon)g_C.
\]
Hence the norm of the axial covector is at most
$[s\sqrt{1-\varepsilon}]^{-1}$, and the volume density relative
to the cylinder lies between
$[s\sqrt{1-\varepsilon}]^3$ and
$[s\sqrt{1+\varepsilon}]^3$.
Product integration gives
$\int_{S^2\times\mathbb R}\phi'(z)\,d\mu_C=A_C$.
The Lipschitz bound alone therefore gives the outer estimate
\[
 |F_2|\leq
 [s_2\sqrt{1+\varepsilon}]^3
       \frac{A_C}{s_2\sqrt{1-\varepsilon}}.                \tag{D}
\]

The inner estimate uses a stronger geometric fact.
Put $\sigma=-1$ when the negative half-neck is compact and
$\sigma=1$ when the positive half-neck is compact, and put
$a=s^{-1}\partial_z$ in neck coordinates.  On the central slab,
almost everywhere,
\[
                    |\sigma\nabla b-a|_g\leq\tfrac14.
                                                               \tag{E}
\]
We now construct the segment and estimate that give (E).
For any point $x$, take unit-speed minimizing segments $c_k$ from
$x$ to $\gamma(k)$.  Their lengths tend to infinity by the triangle
inequality.  At each fixed nonnegative parameter $t$ their values
eventually lie in the compact closed ball $\overline B(x,t)$.
Choose one nonprincipal ultrafilter on the integers and take these
compact limits for every $t$, obtaining a curve $c$.
The same ultrafilter for all parameters lets the distance identities
pass to the limit:
\[
 c(0)=x,\qquad d(c(t),c(u))=|t-u|\quad(t,u\geq0).
\]
In particular $c$ is a continuous minimizing ray.
The defining infimum for $b$, evaluated at $c_k(t)$, gives
\[
 b(c_k(t))\leq d(\gamma(k),c_k(t))-k
                    =d(\gamma(k),x)-k-t.
\]
Pass to the limit and use continuity of $b$ to get
$b(c(t))\leq b(x)-t$.  The reverse inequality is its Lipschitz
bound, so equality holds.

For any prescribed $T>0$, join $x$ to $c(T)$ by a smooth unit-speed
minimizing geodesic $\zeta:[0,T]\to M$.
This geodesic is defined smoothly across the endpoints by completeness.
The Lipschitz inequalities from each endpoint squeeze its values:
\[
 b(x)-t\leq b(\zeta(t))
       \leq b(c(T))+(T-t)=b(x)-t.
\]
At a differentiability point $x$ of $b$, the right derivative at zero
therefore gives $db_x(\dot\zeta(0))=-1$.
Since $|\dot\zeta(0)|=1$ and $|\nabla b(x)|\leq1$, equality in
Cauchy--Schwarz yields $|\nabla b(x)|=1$ and
$\dot\zeta(0)=-\nabla b(x)$.
Taking $T$ larger than the radius of a compact set about $x$ puts
$\zeta(T)$ outside that set.

Apply this with $|z(x)|\leq1$ and with the compact set consisting
of $\overline A$ together with the closed half-neck slab
$|z|\leq(2\varepsilon)^{-1}$.
The first exit from the open half-neck slab occurs at a time $L>0$
on one of its two boundary slices.  Compactness of the closed slab
and continuity justify this first-exit assertion, and the entire
initial segment remains in the neck.  On that segment,
\[
 |(z\circ\zeta)'|\leq[s\sqrt{1-\varepsilon}]^{-1}.
\]
For $\varepsilon\leq1/4$ the axial displacement is at least
$(4\varepsilon)^{-1}$, so
\[
                  L\geq\frac{s}{8\varepsilon}
                           >\frac{s}{100\varepsilon}.
\]
This exit must be outward.  Indeed every point $y$ on the central
sphere satisfies $d(x,y)\leq(2\pi+2)s$, by a spherical path and
the short axial path from $x$.
If $\zeta(L)$ were on the compact-side half-neck slice, the remaining
segment to the exterior endpoint would meet the central sphere
at some later parameter $u$.  Minimality would imply
\[
 d(\zeta(L),\zeta(u))=u-L\leq(2\pi+2)s-L,
\]
because $u=d(x,\zeta(u))$.
But the global bump depth estimate makes this distance greater than
$(2\pi+2)s$ when
$\varepsilon\leq[4C_D(2\pi+2)]^{-1}$.
The contradiction fixes the outward axial sign.

We prove the required direction estimate for any unit-speed neck
geodesic $\zeta:[0,L]\to N$ that minimizes intrinsically on every
subinterval and has increasing endpoint axial coordinate.
Write $f=z\circ\zeta$.  The first neck metric derivative gives a
universal constant $K$ such that
\[
                       |f''|\leq K\varepsilon/s^2.          \tag{G}
\]
Here is the scaling and coordinate justification.
Use a spherical stereographic chart centered at the point in question,
paired with the axial coordinate.  The spherical model coefficient
is a multiple of $16/(|y|^2+4)^2$, whose derivative vanishes at
$y=0$; the axial coefficient is constant.  Thus model Christoffel
symbols vanish there.  The first covariant metric derivative in the
neck bound is consequently the ordinary derivative of the normalized
coefficients at that point.  Its tensor norm bounds each trilinear
evaluation by Cauchy--Schwarz.  Expanding three inputs in the three
coordinate directions bounds the complete derivative operator by a
universal multiple of $\varepsilon$.
After a fixed linear identification with Euclidean three-space, the
actual coefficient form $B$ therefore has bounds
\[
 B(v,v)\geq a_0s^2|v|^2,\qquad
                 |DB|\leq d_0s^2\varepsilon,
\]
with universal $a_0>0$ and $d_0\geq0$.
The inverse form has norm at most $(a_0s^2)^{-1}$.
The three terms of the Koszul formula give
$|\Gamma|\leq3d_0\varepsilon/(2a_0)$.
A unit tangent has coordinate norm squared at most
$(a_0s^2)^{-1}$, so the geodesic equation and the fixed axial
linear functional $\ell$ give (G) with
$K=3|\ell|d_0/(2a_0^2)$.  The centered charts may change from
point to point; the scalar $f''$ and the constant $K$ do not.

For two points of the model cylinder, choose a spherical path of
length at most $\pi+1$ in the unit sphere while linearly interpolating
their axial coordinates.  Its cylinder length is at most
$|\Delta z|+C_0$, where $C_0=\sqrt2(\pi+1)$.
The whole path remains in the axial interval.  Pulling it into the
neck and using the upper metric bound gives, for $0\leq a\leq b\leq L$,
\[
 b-a\leq d_N(\zeta(a),\zeta(b))
       \leq s\sqrt{1+\varepsilon}
                         (|f(b)-f(a)|+C_0).                \tag{H}
\]
Here $d_N$ is the infimum of lengths of paths staying in the neck.
Ambient minimizing geodesics satisfy the first inequality as well,
since allowing more paths can only decrease that infimum.

Set $h=s/(200\sqrt\varepsilon)$ and
$\delta=K\varepsilon/s^2$.
The length assumption $L>s/(100\varepsilon)$ gives $2h\leq L$.
At any $t\in[0,L]$, choose a subinterval $[a,b]$ of length $h$
containing $t$, using the right or left side when necessary.
Equation (G) gives $|f'(u)-f'(t)|\leq\delta h$ throughout it,
and hence $|f(b)-f(a)|\leq h(|f'(t)|+\delta h)$.
Substitution in (H) gives
\[
 s|f'(t)|\geq Q(\varepsilon):=
 (1+\varepsilon)^{-1/2}
       -(200C_0+K/200)\sqrt\varepsilon.
\]
The right side tends to one as $\varepsilon\downarrow0$.
Once it is positive, $f'$ never vanishes.  Continuity and the
increasing endpoint coordinate force its sign to be positive.
Thus for any fixed $\beta>0$, a universal smaller neck threshold
gives $sf'(t)\geq1-\beta$ at every parameter, including both endpoints.

Let $a=s^{-1}\partial_z$ as above and $v=\dot\zeta(t)$.
The zeroth-order bilinear metric error yields
\[
 \bigl||a|_g^2-1\bigr|\leq\varepsilon,\qquad
       |\langle v,a\rangle_g-sf'(t)|\leq2\varepsilon.
\]
For the second bound, the coordinate vector representing $v$ has
scaled model length at most $2$ by the lower metric comparison;
apply the bilinear error to it and the axial vector.
Consequently
\[
 |v-a|_g^2=1+|a|_g^2-2\langle v,a\rangle_g
                                      \leq2\beta+5\varepsilon.
\]
Choosing $\beta$ and then $\varepsilon$ makes this smaller than
any prescribed positive squared tolerance, independently of $s$.
For a decreasing endpoint coordinate, reverse the parametrization
and obtain the same estimate for $-\dot\zeta$.
The outward calibrated segment has axial sign $-\sigma$ and
initial tangent $-\nabla b$.  With tolerance $1/4$, the estimate
therefore gives (E).  Include all thresholds above in the universal
choice of $\varepsilon_*$.

Since $dz(a)=s^{-1}$ and $dT=-\sigma\phi'(z)\,dz$, (E) gives
\[
 -db(\nabla T)\geq
 \phi'(z)\left(s^{-1}
           -\frac{1}{4s\sqrt{1-\varepsilon}}\right).
\]
The right side is nonnegative for $\varepsilon<1/2$.
Integrating with the lower volume comparison yields
\[
 -F_1\geq[s_1\sqrt{1-\varepsilon}]^3 A_C
       \left(s_1^{-1}
            -\frac{1}{4s_1\sqrt{1-\varepsilon}}\right).
                                                               \tag{F}
\]
For completeness the coarse constants can be read directly from
these factors.  Set $u=\sqrt{1-\varepsilon}$ and
$v=\sqrt{1+\varepsilon}$.  Then $u\geq1/2$, $u^2\geq1/2$,
and $u^2(u-1/4)\geq1/8$, while $v\leq2$, $v^3\leq4$,
and $v^3/u\leq8$.  Equations (C)--(F) imply
\[
 \frac{A_Cs_1^2}{8}\leq-F_1\leq-F_2\leq|F_2|
                            \leq8A_Cs_2^2.
\]
Cancel the positive area and take nonnegative square roots to obtain
$s_2\geq s_1/8$, contrary to the chosen scales.
With the separating-side, escape and direction inputs supplied,
this proves the neck obstruction.  Applied to the small necks from
the preceding reduction, it gives the required curvature bound on
the original slice.

At time $-1$ equation (A) now gives admissible static data.
We must identify the entire retained flow, since the equation on
one slice alone does not do this.  Fix a finite slab
$[a,b]\subset(-\infty,0)$ and an anchor $s\in[a,b]$.
Scalar monotonicity and $|\operatorname{Rm}|\leq R$ turn the
right-endpoint slice bound into a single full curvature bound $B$
on this slab.  Curvature contraction gives a Ricci quadratic bound
$|\operatorname{Ric}(v,v)|\leq C_R h(t)(v,v)$; for example the
coarse constant $C_R=27B$ suffices in dimension three.
Equation (A) then gives
\[
 |\nabla^2 f_t(v,v)|\leq C_Hh(t)(v,v),\qquad
 C_H=C_R+\frac1{-2b},\quad f_t=f(\cdot,t).
\]
Integrating the metric equation on fixed tangent vectors yields
uniform norm comparison on the slab:
\[
 |v|_{h(t)}\leq D|v|_{h(r)},\qquad
 D=e^{C_R(b-a)},\quad r,t\in[a,b].
\]
It also compares the intrinsic distances by the same factor.
At the retained basepoint $p$, smoothness of $f$ and of the inverse
metric makes $t\mapsto|\nabla^{h(t)}f_t(p)|_{h(t)}$ continuous.
Compactness of $[a,b]$ bounds it by a number $G$.

The spatial gradient estimate includes points where $df_t=0$.
The symmetric Hessian endomorphism has operator norm at most $C_H$
by its quadratic bound.  Apply the differentiated-norm inequality
to the smooth function
$w_t=(|\nabla^{h(t)}f_t|_{h(t)}^2+1)^{1/2}$.
The sum of the squared covariant derivatives of the gradient is
at most $3C_H^2$, so $|\nabla w_t|\leq4C_H$ is a valid uniform
bound.  Integrating along a minimizing geodesic in the complete
slice and using $w_t(p)\leq G+1$ gives
\[
 |\nabla^{h(t)}f_t(x)|_{h(t)}
           \leq G+1+4C_Hd_{h(t)}(p,x).
\]
Convert both the vector norm and distance to the fixed complete
reference metric $h(s)$.  With
$A=D(G+1)+4C_HD^2+1>0$ we obtain, for
$X_t=-\nabla^{h(t)}f_t$,
\[
                 |X_t(x)|_{h(s)}
                    \leq A(1+d_{h(s)}(p,x)).
\]
No bound on the time derivative of the potential has been assumed.

This estimate prevents a partial trajectory from escaping in finite
time in either direction.  For a trajectory $\gamma$ through $x$
at time $s$, reparametrize any segment from $s$ to $t$ onto $[0,1]$.
Let $v(r)$ denote its speed in $h(s)$ and set
\[
 q(r)=1+d_{h(s)}(p,x)+\int_0^r v(u)\,du.
\]
The length bound gives $1+d_{h(s)}(p,\gamma(r))\leq q(r)$,
and the growth estimate gives
$q'\leq A|t-s|q$.  Gronwall therefore confines every such partial
trajectory to the same closed reference ball of radius
\[
          (1+d_{h(s)}(p,x))e^{A(b-a)}.
\]
The ball is compact by completeness.  This use of an integrated
speed majorant avoids differentiating distance at its cut locus.

The vector field is jointly smooth: in local coordinates its
components are $-h^{ij}(t,x)\partial_j f(t,x)$.
Local ordinary differential equation theory therefore gives smooth
initial-value families.  On the compact product of the time slab
and the confinement ball, finitely many such neighborhoods give
a common positive continuation time.  If a partial family stopped
before a prescribed interior time, restart it close enough to its
endpoint with this common time and shrink the neighborhood of its
initial point so that the nearby trajectories enter the restart
chart.  Uniqueness glues the families on their overlap and extends
past the proposed endpoint.  The same argument works backwards.
Using slightly larger slabs within negative time handles both
endpoints of $[a,b]$.

Uniqueness also makes the constructions on different slabs agree.
We obtain evolution maps $E_{s,t}$ for all $s,t<0$, smooth in the
endpoint and initial point, with
\[
 E_{s,s}=\mathrm{id},\qquad
 E_{t,r}\circ E_{s,t}=E_{s,r},\qquad
 E_{s,t}^{-1}=E_{t,s}.
\]
Thus every spatial map is a diffeomorphism.  Put $E_t=E_{-1,t}$.
Differentiating a transported tangent vector gives
$\nabla_t(dE_tu)=\nabla_{dE_tu}X_t$; metric compatibility then
gives the pullback variation
\[
 \frac{d}{dt}E_t^*h(t)
   =E_t^*(\partial_th+\mathcal L_{X_t}h)
   =E_t^*(-2\operatorname{Ric}-2\nabla^2 f_t)
   =\frac1t E_t^*h(t).
\]
For each pair of fixed tangent vectors, division by $t$ makes the
derivative zero.  Since $E_{-1}=\mathrm{id}$, integration gives
\[
 E_t^*h(t)=(-t)h(-1),\qquad
 h(t)=(-t)(E_t^{-1})^*h(-1).
\]
These inverses supply the homotheties of the actual retained flow.
The static classification therefore applies on its own carrier,
with the original flow and all of its negative-time slices retained.
\end{proof}

\begin{remark}[Geometric inputs]
\label{rem:ancient-neck-input}
The singleton-level and outward-flow lemmas above construct both the
Euclidean and actual-distance radial parametrizations at the same point.
The signed weak flux, numerical scale comparison, whole-carrier escape,
radial maximum argument, ambient sphere separation, calibrated segments
and axial window estimate are also developed above.  Their shared
normal-coordinate and metric-jet inputs are the local geometric
estimates used in Theorem~\ref{thm:rf-compactness}.
The high-level graph construction above uses the same pointed
compactness and fixed-level product inputs as the later splitting
arguments; their hypotheses must be retained at both uses.
\end{remark}

\section{Asymptotic volume and a universal constant}

For a complete nonnegatively Ricci-curved metric define
\[
 \Theta_g(p,r)=\frac{\mu_g(B(p,r))}{r^n},\qquad
 \operatorname{AVR}(g)=\inf_{r>0}\Theta_g(p,r).
\]
There is no division by $\omega_n$.

\begin{proposition}[Volume-ratio limit]
\label{prop:ancient-avr}
For an ancient slice, $\Theta_g(p,r)$ is nonincreasing in $r$, tends
to its infimum as $r\to\infty$, and has basepoint-independent infimum.
It is finite and agrees with the extended-nonnegative-real definition
using calibrated measure.
\lean{PoincareMT.asymptoticVolumeRatioBishopGromov}
\source{MT2007}
\dcref{Theorem 1.34, pp. 19--20; Section 9.6, pp. 221--222}
\end{proposition}
\begin{proof}
Nonflatness excludes dimension zero.  Completeness makes bounded closed
sets compact, so finite-radius balls have finite volume.
The Ricci lower bound makes the polar Jacobian divided by $r^{n-1}$
nonincreasing before the cut time.  Extend it by zero afterward and
integrate the polar volume formula to obtain Bishop--Gromov.
The small-radius value is $\omega_n$, so the infimum is finite and
equals the limit at infinity.  Calibration identifies the measure.
If $D=d(p,q)$ and $r>D$, then
$B(p,r-D)\subseteq B(q,r)\subseteq B(p,r+D)$.
Divide by $r^n$ and let $r\to\infty$; the factors
$(1\pm D/r)^n$ tend to one.  This proves basepoint independence,
also after embedding finite values in the extended nonnegative reals.
\end{proof}

\begin{proposition}[Bounded ancient volume theorem]
\label{prop:bounded-ancient-zero-volume}
Let $n\geq2$ and let $g(t)$, $t\leq0$, be a Ricci flow on a connected
noncompact $n$-manifold.  Suppose every slice is complete with
nonnegative curvature operator, and there is a single $C\geq0$
with $A(t,x)\leq C$ for every $t\leq0$ and $x$.
Suppose some $q$ satisfies $R(0,q)>0$ and a single
$\kappa>0$ satisfies the parabolic volume implication for every
terminal ball and every closed cylinder $[t-r^2,t]$.
Then for every $t\leq0$ and every $p$,
\[
 \operatorname{AVR}(g(t))=0,\qquad
 \lim_{r\to\infty}\frac{\mu_{g(t)}(B_t(p,r))}{r^n}=0 .
\]
\lean{PoincareMT.RicciFlow.zero_asymptoticVolumeRatio_of_bounded_ancient}
\dcref{Kleiner--Lott, corrected 2013, Theorem 41.2, pp. 2674--2677,
and Proposition 41.13, p. 2678}
\end{proposition}
\begin{proof}
We use the intermediate results proved in the following subsections.
Lemma~\ref{lem:ancient-infinite-scalar-ratio} gives unbounded scalar
ratio in every dimension $n\geq2$ under the stated hypotheses.
Suppose the volume ratio at some $t_1\leq0$ were positive.
Lemma~\ref{lem:ancient-selected-line} constructs a complete bounded
ancient limit $G$ with positive scalar curvature at its terminal
basepoint, positive terminal asymptotic volume, and a minimizing line.
Lemma~\ref{lem:ancient-line-busemann} supplies a unit parallel gradient
on that terminal slice.

We argue by induction on $n\geq2$.  When $n=2$, a nonzero parallel
vector is curvature-null.  The surface identity
\[
 \mathrm{Rm}(a,b,c,d)=\tfrac R2
          \bigl(g(a,c)g(b,d)-g(a,d)g(b,c)\bigr)
\]
then forces $R=0$, by pairing the unit parallel vector with a
perpendicular unit vector.  This contradicts $R_G(0,o)>0$.
For $n>2$, Lemma~\ref{lem:ancient-parallel-persistence} preserves
the same gradient through the ancient past, and
Lemma~\ref{lem:ancient-fixed-factor} constructs a connected
noncompact factor of dimension $n-1$ satisfying every hypothesis
of the proposition, with a single positive noncollapse constant.
Its terminal asymptotic volume is positive, contradicting the
induction hypothesis.  The original ratio is therefore not positive.
Its nonnegativity and Proposition~\ref{prop:ancient-avr} give both
asserted conclusions.
\end{proof}

\subsection{The scalar-ratio alternative}

Throughout the next three lemmas, $g(t)$ satisfies the hypotheses
of Proposition~\ref{prop:bounded-ancient-zero-volume}; its conclusion
is not assumed.  We use bounded ancient Harnack, available directly
under the uniform curvature bound.  The extension to flows with
only slice-wise bounds is not used here.

\begin{lemma}[Infinite scalar ratio]
\label{lem:ancient-infinite-scalar-ratio}
For every $t_0\leq0$, $p$, and real $L,D$, there is $x$ with
\[
 d_{t_0}(p,x)>L,\qquad R(t_0,x)d_{t_0}(p,x)^2>D.
\]
\end{lemma}
\begin{proof}
If the displayed unboundedness fails, put
$f(x)=R(t_0,x)d_{t_0}(p,x)^2$ on the offending slice.
It is nonnegative and bounded outside a compact ball.
The decreasing tail suprema therefore have a finite limit
$\Lambda\geq0$.  Choose $q_k$ outside the radius-$k$ ball whose
$f$-value is within $1/k$ of its tail supremum.
Then $d_{t_0}(p,q_k)\to\infty$ and $f(q_k)\to\Lambda$;
every $C>\Lambda$ is a uniform bound outside some ball.
There is a positive scalar value somewhere on every earlier slice.
Otherwise nonnegative scalar curvature would vanish identically on
one slice.  The uniform curvature bound and
$\partial_tR=\Delta R+2|\operatorname{Ric}|^2\leq\Delta R+CR$
allow forward scalar comparison, keeping it zero through time zero,
contrary to the assumed positive terminal value.  Apply this at
$t_0-1$; integrated bounded ancient Harnack along a path from that
positive point to $q_k$ gives $R(t_0,q_k)>0$.
A positive $\Lambda$ contradicts
Lemma~\ref{lem:ancient-positive-ratio-obstruction}; $\Lambda=0$
contradicts Lemma~\ref{lem:ancient-zero-ratio-obstruction}.
\end{proof}

\begin{lemma}[Obstruction at a finite positive ratio]
\label{lem:ancient-positive-ratio-obstruction}
Fix $t_0\leq0$ and $p$.  Suppose the function
$R(t_0,x)d_{t_0}(p,x)^2$ is bounded outside a compact ball.
There is no escaping sequence $q_k$ for which its values converge
to a finite positive number $\Lambda$.
\end{lemma}
\begin{proof}
Suppose such a sequence exists.  Write $Q_k=R(t_0,q_k)>0$, after
discarding finitely many terms, so that
$Q_kd_{t_0}(p,q_k)^2\to\Lambda\in(0,\infty)$.
Then $Q_k\to0$.  Rescale by
$g_k(s)=Q_kg(t_0+s/Q_k)$, $s\leq0$.
The normalized base scalar is one and
$d_{g_k(0)}(p,q_k)\to\sqrt\Lambda$.
We describe the annular normalization and the retention of terminal
time in the local limit.  Write the finite-ratio tail bound as
$R(t_0,x)d_{t_0}(p,x)^2\leq C$ when $d_{t_0}(p,x)\geq L$.
Choose $b,r>0$ such that
\[
 b+2r<\sqrt\Lambda,\qquad n^2C/b^2\leq r^{-2}.
\]
For example first take $b=\sqrt\Lambda/2$, then decrease $r$.
Eventually $\sqrt{Q_k}L<b$ and
$d_{g_k(0)}(p,q_k)>b+2r$.
Every point in the terminal $2r$-ball about $q_k$ then has
normalized distance at least $b$ from $p$.  The scalar-distance
product is invariant under this rescaling, so its terminal scalar
curvature is at most $C/b^2$.  Bounded ancient Harnack at a fixed
spatial point gives $R_k(s,x)\leq R_k(0,x)$ for every $s\leq0$.
Nonnegative curvature operator therefore bounds the full curvature
norm by $n^2C/b^2$ on the same fixed ball for the entire past.

Apply parabolic noncollapsing to the terminal $r$-ball: the required
curvature bound holds on its whole backward cylinder.
Rescaling volume by $Q_k^{n/2}$ gives
$\mu_{g_k(0)}(B(q_k,r))\geq\kappa r^n$.
The local injectivity estimate from curvature on the $2r$-ball
and volume on the $r$-ball supplies a common normal radius
$S\in(0,r)$.  Here is the local mechanism of that estimate.
For a metric with $|\operatorname{Rm}|\leq K$ on $B(p,2r)$
and $\mu(B(p,r))\geq v>0$, choose $s\leq r$ within the
Jacobi comparison radius.  The exponential $e$ from an orthonormal
tangent frame is defined on $B(0,r)$ by compact confinement.
On $B(0,s)$ its differential satisfies
\begin{equation}
\label{eq:ancient-local-injectivity}
 \tfrac12|w|\leq |de_xw|_g\leq\tfrac32|w|.
\end{equation}
Indeed the radial Jacobi equation in a parallel frame has integral
form $J(t)=tw-\int_0^t(t-u)\mathcal R(u)J(u)\,du$;
decreasing the comparison radius makes the error at $t=1$
at most $|w|/2$.  Relative volume comparison gives
\[
 \mu(B(p,s/4))\geq
 w_0:=v\,\frac{V_{-K}(s/4)}{V_{-K}(r)}>0,
\]
where $V_{-K}$ is the volume of a ball in the simply connected
space of constant sectional curvature $-K$.
The total pullback volume available in $B(0,s)$ is at most
$P=n!(3/2)^n\omega_ns^n$ by the differential bound.
Choose an integer $N\geq1$ with $(N+1)w_0>P$.

Failure of injectivity on a sufficiently small closed tangent ball
produces a nonzero radial return $e(v_*)=p$ of length
$\ell=|v_*|$ at most twice that ball's radius.  To see the geometric
alternative, take a first collision before the conjugate radius.
Compactness gives a minimizing collision pair: local injectivity
prevents such pairs from accumulating on the diagonal.  Both radial
segments minimize and have the same length.  In the two local inverse
branches, the gradients of half the squared radial norm are the
terminal radial velocities and have equal norm.  Unless they are
opposite, their negative sum decreases both radial norms, contradicting
minimality of the collision.  Continuing one radial geodesic along
the reverse of the other gives the return.
Lift successive traversals of this short loop through the same
local exponential.  The inverse differential bound is $2$, so
their endpoints $y_i$, $0\leq i\leq N$, satisfy
\[
 e(y_i)=p,\qquad |y_i|\leq2i\ell.
\]
The lift continues as long as its length bound leaves it strictly
inside $B(0,s)$: at a putative last time it has a limit in a
compact inner ball, and a local inverse continues it.
Uniqueness follows from uniqueness on local inverse neighborhoods.

One must prove these endpoints distinct.  Lifting radial paths
from the returning point defines a local isometry $d$ of the
pullback metric, with $e\circ d=e$, $d(0)=v_*$ and
$|d(x)-v_*|\leq2|x|$.  The confined radial lifts depend continuously
on their endpoint by local-inverse continuation.  They therefore
agree locally with smooth inverse branches, making $d$ smooth;
differentiating $e\circ d=e$ proves the local-isometry assertion.
It has no fixed point: uniqueness of the
reverse lift of a radial path would otherwise imply $v_*=0$.
A repeated loop endpoint would, by the same uniqueness and
continuation, give $d^m=\mathrm{id}$ on a smaller connected ball
for some $1\leq m\leq N$.
For concreteness it suffices to require
\[
 \ell\leq\frac{s}{1024(N+1)^4 4^N},\qquad
 a=\frac{s}{16\,2^N},\quad r_*=a/2,\quad\delta=2N\ell.
\]
All the iterates used below stay in $B(0,s/4)$, while
$\delta<r_*/16$ and $N\delta^2<(r_*/16)^2$.
The energy $E(x)=\sum_{i=0}^{m-1}|d^i(x)|^2$ attains its
minimum on $\overline B(0,r_*)$.  Its identity term and
$E(0)\leq m\delta^2$ put a minimizer $y$ in $B(0,r_*/16)$.
The displacement bound puts $d(y)$ in $B(0,3r_*/16)$;
periodicity makes it another minimizer.

These two points can be joined by a minimizing geodesic confined
to $B(0,r_*)$.  One may extend the pullback metric outside its
controlled ball, preserving its Gauss identity and the bounds
$\tfrac12|w|\leq|w|_g\leq\tfrac32|w|$; comparison with the
straight segment then confines the minimizing geodesic to the
original domain.  On that domain radial Jacobi comparison gives
\[
 \frac{d^2}{dt^2}|\gamma(t)|^2\geq|\dot\gamma(t)|^2.
\]
More explicitly, the radial Jacobi field $J(u)=u w$ satisfies
$\langle D_uJ(1),J(1)\rangle\geq |w|^2/2$ at the chosen
comparison scale; Gauss' identity identifies twice this pairing
with the Hessian of the squared radial norm.  Each
$d^i\circ\gamma$ is a geodesic because $d^i$ is a local
isometry.  Their squared norms are convex, and the identity term
is strictly convex when $y\ne d(y)$: a nonconstant affine
geodesic has nonzero constant intrinsic speed.  Hence $E$ cannot
have both endpoint minima.  This excludes finite periods and
proves the $N+1$ endpoints distinct.

Lift a minimizing radial path to each point of $B(p,s/4)$
from each $y_i$.  The margins $|y_i|+s/2<s$ keep all lifts in
the tangent ball, and uniqueness keeps their endpoints distinct.
Thus every point of this intrinsic ball has at least $N+1$
preimages.  The change-of-variables formula on local inverse
pieces now gives the volume bound: choose a countable cover by
injective neighborhoods, disjointify it into measurable pieces,
and add their change-of-variables identities by Tonelli.
Every image point is counted at least $N+1$ times, so
\[
 (N+1)\mu(B(p,s/4))\leq P,
\]
contrary to the choice of $N$.  This proves a positive uniform
injectivity radius depending only on $n,K,r,v$.
Minimizing-geodesic coverage and injectivity identify the image
of each smaller tangent ball with the intrinsic ball of the same
radius, and identify radial length with distance.
Thus there are time-independent terminal exponential
charts $\Phi_k:B(0,S)\to B_{g_k(0)}(q_k,S)$, normalized to the
Euclidean metric at zero.  Their radial geodesics have exactly
their Euclidean parameter lengths.  All earlier curvature norms
on the terminal $2S$-ball are bounded by one $K_0>0$.

The local derivative estimates can be used at every earlier time
on the fixed terminal $S$-ball.  Center an estimate at the point
and time being estimated.  Its earlier-time ball of a fixed
smaller radius lies in the corresponding terminal ball, since
distances decrease in time.  That ball lies in the controlled
$2S$-ball.  The ancient time range supplies a full backward buffer.
Consequently every $|\nabla^j\operatorname{Rm}|$ has a uniform
bound there, independent of the time and the sequence index.
To obtain terminal coefficient bounds, parallel transport an
orthonormal frame along each coordinate ray.  Denote the transport
matrix by $T(x)$ and its inverse coframe by $A(x)$.
The radial variation in this parallel frame is
$J(t,x)=tA(tx)$; it satisfies
\[
 J''+\mathcal R(t,x)J=0,\qquad J(0,x)=0,\quad J'(0,x)=I,
\]
where $\mathcal R(t,x)$ is the transported curvature endomorphism
contracted twice with the initial ray velocity $x$.
Differentiating its integral equation in $x$ bounds the
parameter derivatives of $J$ by induction and Gronwall.
At a given derivative order the only term containing the highest
derivative of $J$ is $\mathcal R D_x^mJ$; the other terms use
lower derivatives of $J$ and derivatives of $\mathcal R$.

The curvature coefficient bounds themselves follow by a coupled
induction.  A coordinate derivative of a component of
$\nabla^l\operatorname{Rm}$ in the radial frame is a component
of $\nabla^{l+1}\operatorname{Rm}$ multiplied by the coframe,
plus one connection term for each of its $4+l$ slots.
The radial connection formula follows from a transverse variation
of parallel transport.  For constant coordinate vectors $u,v$, put
$Y_v(y)=T(y)v$ and
$H(t)=T(tx)^{-1}(t\nabla_uY_v(tx))$.  Radial parallelism and
the covariant-derivative commutator give
\[
 H(0)=0,\qquad
 H'(t)=\mathcal K(tx)\bigl(x,tA(tx)u\bigr)v,
\]
where $\mathcal K$ is the curvature operator in the parallel frame.
The ray velocity is parallel, so $T(tx)x=x$ and $A(tx)x=x$;
this explains the first argument $x$ in the displayed curvature term.
The initial value is zero because the frame is smooth at the center.
Integrating gives the connection matrix $\omega_x(u)v=H(1)$:
\begin{equation}
\label{eq:ancient-radial-connection}
 \omega_x(u)v=\int_0^1
       \mathcal K(tx)\bigl(x,tA(tx)u\bigr)v\,dt.
\end{equation}
Each component is thus the radial integral of a sum of terms
$y_iA_{kj}(y)\mathcal K_{ik}(y)$ evaluated at $y=tx$.
Differentiation $m$ times in $x$ contributes $t^m$; a uniform
bound for the differentiated integrand gives a bound divided by
$m+1$.  Differentiation under this integral is justified on a
compact smaller ball, including $t=0$.  Leibniz' rule uses coframe
and curvature derivatives only through order $m$.
At induction order $m+1$, the previously
bounded order-$m$ coframe and connection derivatives therefore
bound order-$m+1$ curvature components.  The differentiated
Jacobi equation next bounds the order-$m+1$ coframe, and the
connection integral bounds its connection derivatives.
Only covariant curvature derivatives through the requested
order are used.  Finally the terminal metric has coefficients
$\sum_a A_{ai}A_{aj}$, so Leibniz' rule bounds every metric
derivative on a smaller chart.  Write the time-dependent
coefficients in these fixed charts as $B_k(s,x)$.

One must choose the spatial neighborhood before the length of a
past slab.  Let $Z\geq0$ bound the first spatial derivative of
$B_k(0,\cdot)$ on $B(0,S/4)$ and put
\[
 \rho=\min\{S/8,1/[2(Z+1)]\}.
\]
Normalization and the mean-value estimate give
$B_k(0,x)\geq\tfrac12 I$ on $\overline B(0,\rho)$.
The metric decreases forward in time, so
$B_k(s,x)\geq B_k(0,x)\geq\tfrac12 I$ for every $s\leq0$.
For a fixed $a>0$, metric comparison gives
\[
 \|B_k(s,x)\|\leq e^{2n^3K_0a}\|B_k(0,x)\|
                  \quad(-a\leq s\leq0).
\]
This upper bound may depend on $a$; the chosen radius does not.

Spatial derivatives propagate backward inductively.  Once lower
metric derivatives are bounded, the coordinate form of
$\partial_sg=-2\operatorname{Ric}$ and the bounds on covariant
curvature derivatives give, at each fixed coordinate point,
\[
 \|\partial_s D_x^{j+1}B_k\|
       \leq C_j\bigl(1+\|D_x^{j+1}B_k\|\bigr).
\]
The highest metric derivative enters linearly; all other terms use
lower derivatives and bounded inverse metrics.  Here this
linearity is essential and can be checked directly.
Differentiating $B^{-1}B=I$ gives
\[
 D^m(B^{-1})=-B^{-1}(D^mB)B^{-1}
                 +\text{terms in derivatives of order less than }m.
\]
Likewise $\Gamma=\tfrac12 B^{-1}*(DB)$ has its order-$(m-1)$
derivative affine in $D^mB$.  Repeatedly using
$D\operatorname{Rm}=\nabla\operatorname{Rm}
+\Gamma*\operatorname{Rm}$ shows that $D^m\operatorname{Rm}$
is bounded by $C(1+|D^mB|)$ when lower metric derivatives and
covariant curvature derivatives through $m$ are bounded.
Here $*$ denotes a finite sum of contractions in the fixed
coordinate basis.  Finally
$\operatorname{Ric}_{ij}=\sum_{ab}B^{ab}\operatorname{Rm}_{iajb}$.
In its order-$m$ Leibniz expansion the only highest-order
terms differentiate exactly one factor $m$ times; all intermediate
terms involve bounded lower derivatives.  Thus no quadratic
highest-jet term occurs.  Integrating this
inequality from the terminal value gives a bound on the entire
closed slab by Gronwall.  In coordinates the Ricci-flow operator
is a smooth function of the metric, its inverse, and its first
two spatial derivatives.  Differentiate this equation repeatedly:
mixed time and space derivatives are controlled by finitely many
higher spatial derivatives.  The finite collections of these
derivatives range in compact sets with metric lower bound
$\tfrac12 I$, so all derivatives of that coordinate operator are
uniformly bounded there.  This gives all mixed derivative bounds
on each finite past slab.

These bounds retain $s=0$.  Each source coefficient is smooth up
to that boundary from the past; continuity extends its interior
derivative bounds to the closed half-cylinder
$(-\infty,0]\times\overline B(0,\rho)$.
On a compact convex part of this domain, a bound on derivatives
of order $j+1$ makes the order-$j$ derivatives equicontinuous by
integration along line segments.  Arzela--Ascoli and one diagonal
subsequence give uniform limits for every derivative on every
compact subdomain, including the terminal boundary.  In the
interior, integration along segments identifies consecutive
limits as derivatives.  Continuity extends that identity to the
boundary, again along segments in the closed convex domain.
Thus the coefficient limit is smooth up to time zero from the
past, with all of the limiting derivatives just obtained.

It remains positive on the fixed smaller open ball by the lower
bound.  The Ricci-flow equation passes to the limit using one time
and two spatial derivatives, and curvature converges using the
same metric derivatives and inverse coefficients.  The resulting
local ancient flow has nonnegative curvature operator and scalar
one at its terminal center.  Completeness of this annular region
is neither asserted nor needed; no future-time extension is used.

The actual terminal radial-square functions in those charts are
\[
 F_k(x)=\tfrac12d_{g_k(0)}(p,\Phi_k(x))^2.
\]
On a fixed coordinate closed ball of radius $r$, normal charts
give $d_{g_k(0)}(q_k,\Phi_k(x))=|x|$.
Choose $D$ bounding the distances from $p$ to $q_k$, and let
$b$ bound the norm of the pulled-back metric coefficients on a
slightly larger ball.  Straight coordinate segments have length
at most $\sqrt b\,|x-y|$.  The triangle inequality gives
$d(p,\Phi_k(x))\leq D+r$, so
\[
 |F_k(x)-F_k(y)|\leq(D+r)\sqrt b\,|x-y|,
 \qquad 0\leq F_k(x)\leq\tfrac12(D+r)^2.
\]
Arzela--Ascoli supplies a uniformly convergent subsequence with
Lipschitz limit $F$ and $F(0)=\Lambda/2>0$; the already convergent
metric jets retain their limits on this subsequence.

We will use actual ambient distances between chart images.
Coefficient convergence alone must be localized before it controls
those distances.  Choose $R$ inside the common normal-coordinate
radius and then $r>0$ so that $3r<R$ and the limit intrinsic
$3r$-ball is contained in the coordinate closed $R$-ball.
Restrict endpoints to the intersection of the coordinate and limit
intrinsic $r$-balls.  Almost minimizing limit paths stay in the
intrinsic $3r$-ball.  Almost minimizing source paths stay in the
source $3r$-ball about $q_k$; the exact target and radial properties
of $\Phi_k$ put their inverse images in the coordinate $R$-ball.
Uniform metric convergence there compares lengths in both directions
by every fixed factor $C>1$ for sufficiently large $k$.
The distances are bounded by a fixed multiple of $r$; letting
$C\downarrow1$ proves uniform convergence of ambient source
distances to the retained limit metric distances on this smaller
region.  Shrink once more so that $F>0$ throughout it.

Here is the radial variation that identifies $F$, without assuming
a global cone map.  In the original terminal metric, any escaping
sequence $x_k$ has, after extraction, a minimizing ray $\alpha$
from $p$ such that
\[
 \frac{d(x_k,\alpha(d(p,x_k)))}{d(p,x_k)}\longrightarrow0.
\]
Indeed minimizing segments to $x_k$, parametrized by unit speed,
have limits on compact parameter intervals by properness.
One may choose a common ultrafilter before the parameter, giving
a ray by passage of the exact distance identity to the limit,
then extract a subsequence converging at parameter one.
Writing $r_k=d(p,x_k)$ and $\gamma_k$ for these segments,
corresponding-side Toponogov comparison gives
\[
 \frac{d(x_k,\alpha(r_k))}{r_k}
       \leq d(\gamma_k(1),\alpha(1))\longrightarrow0.
\]
Successive extraction supplies this approximation simultaneously
for any fixed finite family of escaping sequences.

For two such rays $\alpha,\beta$, the comparison cosine
\[
 c(a,b)=\frac{a^2+b^2-d(\alpha(a),\beta(b))^2}{2ab},
                  \qquad a,b>0,
\]
lies in $[-1,1]$ by the triangle inequalities and is nondecreasing
in each variable by the same corresponding-side comparison.
Its diagonal has a limit $c_\infty$.  Squeezing between the
diagonal values at $\min(u,v)L$ and $\max(u,v)L$ gives
\[
 \lim_{L\to\infty}\frac{d(\alpha(uL),\beta(vL))^2}{L^2}
       =u^2+v^2-2uvc_\infty\quad(u,v>0).
\]
The relative ray errors above, multiplied by
$\lambda_k=\sqrt{Q_k}$, tend to zero because the normalized
radii have positive finite limits.  Varying radii can be replaced
by their limiting multiples of $\lambda_k^{-1}$: the error along
a unit-speed ray is exactly the change in its parameter.
Thus the same distance formula holds for the original source
sequences and their radial perturbations.

Fix three limit points $x,z,y$ in the smaller region and approximate
their source images simultaneously by rays.  Move the ray for $z$
from its original radial parameter to $a$ times that parameter,
where $a$ is close to one.  The moved source points remain in a
fixed smaller normal ball: their additional normalized displacement
is bounded by $|a-1|\sqrt{2F(z)}$ plus a vanishing ray error.
Their inverse coordinates consequently have a convergent subsequence.
Uniform distance convergence gives a limit point $w_a$ such that,
for $v=x,y$,
\[
 d(v,w_a)^2=P_v(a),\qquad
 P_v(a)=2a^2F(z)+2F(v)
           -a\bigl(2F(z)+2F(v)-d(v,z)^2\bigr).
\]
This formula follows by eliminating $c_\infty$ with the unperturbed
distance at $a=1$.  The same moved point works for both endpoints;
no continuity of the choice $a\mapsto w_a$ is required.

Suppose $z$ is strictly between $x$ and $y$ on a minimizing segment.
The triangle inequality makes
$\sqrt{P_x(a)}+\sqrt{P_y(a)}$ minimal at $a=1$, where it equals
$d(x,z)+d(z,y)=d(x,y)$.  Both square roots are differentiable
there because the two subsegment lengths are positive.  Its
derivative being zero gives
\[
 \frac{2F(z)-2F(x)+d(x,z)^2}{2d(x,z)}+
 \frac{2F(z)-2F(y)+d(y,z)^2}{2d(y,z)}=0.
\]
Substitute $d(x,z)=s\,d(x,y)$ and
$d(z,y)=(1-s)d(x,y)$ and rearrange.  Endpoint and constant-segment
cases hold directly.  Thus every short minimizing segment
$\gamma:[0,1]\to N$ in the annulus satisfies
\[
 F(\gamma(s))=(1-s)F(\gamma(0))+sF(\gamma(1))
       -\tfrac12s(1-s)d(\gamma(0),\gamma(1))^2 .
\]
We give the regularity and rigidity steps explicitly.  Every affine
geodesic is minimizing on a neighborhood of each parameter value.
On such an interval the displayed identity says that
$F(\gamma(s))-|\gamma'|^2s^2/2$ is affine.  The speed is constant;
the affine formulas agree on overlapping intervals.  Connectedness
of the parameter interval therefore extends the quadratic identity
along every geodesic contained in the coordinate region.

Local Lipschitzness does not initially give differentiability at a
specified point $x$.  Instead use joint normal coordinates.
The map $(q,v)\mapsto(q,\exp_qv)$ has an invertible differential
at $(x,0)$: its first component is $q$ and its velocity derivative
in the second component is the identity.  The inverse function
theorem supplies a common endpoint neighborhood for nearby centers
$q$, all containing $x$.  Rademacher's theorem supplies a center
$q$ in this open set at which $F$ is differentiable.  Differentiating
the quadratic formula along $s\mapsto\exp_q(sv)$ at $s=0$ gives
\[
 F(\exp_qv)=F(q)+dF_q(v)+\tfrac12 h_q(v,v).
\]
The right side is smooth in $v$, and the normal endpoint inverse is
smooth near $x$.  Thus $F$ is smooth near every $x$ in the smaller
annulus.  Differentiating twice along a geodesic of arbitrary initial
velocity gives $\operatorname{Hess}F(v,v)=h(v,v)$.
Polarization gives $\operatorname{Hess}F=h$, and metric
compatibility yields the smooth field $V=\nabla F$ with
$\nabla_XV=X$.

The curvature commutator applied to $V$ is
\[
 \mathcal R(X,Y)V
   =\nabla_XY-\nabla_YX-[X,Y]=0,
\]
by vanishing torsion.  The curvature symmetries then make every
slot containing $V$ vanish.  Write $T=\operatorname{Rm}$ for the
terminal covariant tensor, with $T(u,w,u,w)$ nonnegative on
sectional planes.  Differentiating a null slot gives
\[
 (\nabla_uT)(a,b,c,V)=-T(a,b,c,u).
\]
The analogous identities in the other slots imply
$(\nabla_uT)(V,a,V,b)=0$: after differentiating either occurrence
of $V$, one null slot remains.  Differentiate this last identity
once more in direction $u$.  At the point in question choose
extensions with vanishing covariant derivatives for $u,a,b$.
The two differentiated $V$ slots give
\[
 (\nabla^2_{u,u}T)(V,w,V,w)=2T(u,w,u,w).
\]
Tracing in an orthonormal frame proves
$(\Delta T)(V,w,V,w)=2\operatorname{Ric}(w,w)$.
The Laplacian here acts on the tensor before evaluation in $V$;
it is not the scalar Laplacian of the identically zero evaluation.

Now fix $x,V(x),w$ as time-independent tangent vectors and consider
$f(t)=\operatorname{Rm}_{g(t)}(V,w,V,w)$ for $t\leq0$.
Nonnegative curvature operator gives $f(t)\geq0=f(0)$, hence its
left derivative at zero is nonpositive.  The fixed-input curvature
evolution identifies this derivative with the tensor Laplacian plus
the quadratic curvature reaction and its Ricci metric corrections.
Every quadratic summand contains a curvature factor with a $V$
slot, and the Ricci corrections vanish because
$\operatorname{Ric}(V,\cdot)=0$ and all curvature slots containing
$V$ vanish.  The entire reaction is therefore zero.
Consequently $2\operatorname{Ric}(w,w)\leq0$ for every $w$.
Nonnegative curvature operator gives the reverse inequality, so
Ricci and scalar curvature vanish.  A nonnegative curvature operator
with zero trace is zero, and the full tensor vanishes on the annulus.
This contradicts its scalar-one center.  The argument uses the
actual ancient annular limit, not merely an abstract metric cone.
The local differential obstruction is
\cite[Lemma 41.4 and Theorem 41.2, Case 1, equation (41.9),
pp.~2674--2676]{KleinerLott2013}; it requires neither completeness
of this coordinate region nor an extension beyond terminal time.
\end{proof}

In the finite-ratio construction, the scalar-tail estimate first
chooses the annular margin and normal-chart radius, parabolic noncollapsing
retains a whole-past curvature bound, and the closed half-cylinder compactness
theorem retains all mixed jets at time zero.  The
limiting coefficient field defines a Ricci flow on an open coordinate set;
smooth convergence transfers the curvature operator and the scalar-one
terminal normalization.  Thus the contradiction applies to a local ancient
flow with the original terminal radial charts.

The zero-ratio branch instead uses a distance-normalized flat annular
limit and a global backward enclosing atlas.  Compatible charts give
a smooth isometric immersion and a lower-distance estimate.  The
annular and whole-ball approximation arguments below then identify
the cone link and prove Euclidean volume rigidity.
The corresponding-side Toponogov comparison remains an intrinsic input to
the radial variation and link construction, not a consequence of coefficient
convergence.

\subsection{The zero-ratio cone and its link}

\begin{lemma}[Obstruction at zero scalar ratio]
\label{lem:ancient-zero-ratio-obstruction}
There is no slice $t_0\leq0$ and center $p$ for which, for every
$\varepsilon>0$, the inequality
$R(t_0,x)d_{t_0}(p,x)^2\leq\varepsilon$ holds outside a
sufficiently large ball.
\end{lemma}
\begin{proof}
Suppose such a slice exists.  Thus for each $\varepsilon>0$,
$R(t_0,x)d_{t_0}(p,x)^2\leq\varepsilon$ outside a sufficiently
large ball.  We first construct the metric link, then obtain uniformly
sized smooth charts around a dense family of its directions.  For
unit-speed minimizing rays from $p$ define
\[
 d_\infty(\alpha,\beta)
   =\lim_{L\to\infty}\frac{d_{t_0}(\alpha(L),\beta(L))}{L}
   =\sqrt{2-2c_\infty(\alpha,\beta)}.
\]
The comparison-cosine argument above proves existence.  Limits of
metric triangle inequalities make this a pseudometric; identify
two rays when their distance is zero, and call the metric quotient
$S$.  This is its \emph{chord} metric, denoted $d_S^{\rm chord}$.
Corresponding-side comparison gives
$d_\infty(\alpha,\beta)\leq d_{t_0}(\alpha(1),\beta(1))$.
The space of rays is compact in the topology of convergence on
bounded parameter intervals: their values lie in compact balls,
they are $1$-Lipschitz, and the exact minimizing-distance identity
passes to the limit.  The last inequality makes its projection to
$S$ continuous.  Hence $S$ is compact.

Define the metric cone $C(S)$ by adjoining a radius $r\geq0$
to each point of $S$ and identifying all points with radius zero.
Its distance is
\[
 d_C((r,\xi),(s,\eta))^2
       =(r-s)^2+rs\,d_S^{\rm chord}(\xi,\eta)^2.
\]
The homogeneous ray limits proved above identify this formula
with limits of actual distances divided by $L$; this proves its
triangle inequality, including at the vertex.
Radius is $1$-Lipschitz.  Closed bounded radial regions are compact,
as images of a compact interval times $S$.
At this stage no Riemannian cone metric or manifold link is assumed.

The approximation by rays is uniform at infinity.  Otherwise a
sequence of escaping points would have a fixed positive relative
distance from every ray, contradicting the minimizing-segment
extraction and relative-error estimate already proved.
Also $d(\alpha(L),\beta(L))/L$ decreases to $d_\infty$ uniformly
over all ray pairs.  Indeed each finite-radius function is
continuous, its limit is continuous by the quotient projection,
and the ray-pair space is compact; Dini's theorem applies.
The same argument with radii in a fixed compact positive interval
gives uniform convergence of the homogeneous distance formula.
Consequently each fixed rescaled annulus has a relation to the
corresponding annulus of $C(S)$ that covers both spaces and whose
distance distortion tends to zero: relate a point to the ray
approximating it, retaining its normalized radial coordinate.
The triangle inequality bounds the extra error by the two
ray-approximation errors.

Choose a countable dense family of represented directions
$\eta_j$ in $S$, and use the actual centers
$q_{k,j}=\eta_j(L_k)$, where $L_k=k+1$, with rescaled flows
\[
 G_k(s)=L_k^{-2}g(t_0+L_k^2s).
\]
Their terminal distance from $p$ is exactly one.
On every fixed annular region separated from $p$, the zero-ratio
condition gives $\sup R_{G_k}(0,\cdot)\to0$.
Bounded ancient Harnack extends this decay through the whole past
on the fixed terminal region.  The radius can be chosen before
choosing a ray or extracting a subsequence.  Indeed, use the
zero-ratio bound with $A=1/(n^2+1)$.  Eventually its exterior
threshold, divided by $L_k$, is below $1/2$.  Every center of
normalized radius at least $3/4$ then has a terminal $1/4$-ball
on which the whole-past curvature bound is
$4n^2/(n^2+1)\leq4$.  On its closed backward cylinder of radius
$1/8$ this is smaller than $(1/8)^{-2}$, so noncollapsing gives
volume at least $\kappa(1/8)^n$ to the terminal $1/8$-ball.
The local injectivity argument above, with these fixed curvature,
radius and volume bounds, gives one $S_0\in(0,1/8)$ for every
such center.  Discard a single finite initial set of indices.
The resulting normalized exponential charts satisfy
\[
 \Phi_{k,j}:B(0,S_0)\longrightarrow B_{G_k(0)}(q_{k,j},S_0),
 \qquad d_{G_k(0)}(q_{k,j},\Phi_{k,j}(x))=|x|.
\]

Radial Jacobi estimates choose one $r_0>0$, with
$2r_0<S_0/4$, on which all pulled-back quadratic forms lie between
$\tfrac14|v|^2$ and $\tfrac94|v|^2$.
The terminal normal-coordinate derivative bounds apply with the
same curvature bound $5$ and radius $S_0$.  A diagonal over
the countably many charts, derivative orders and compact
subballs gives one subsequence with smooth coefficient limits
$B_j$ on $B(0,S_0/4)$.  They are positive on $B(0,2r_0)$,
normalized at zero, and satisfy
$B_j(x)(x,v)=\langle x,v\rangle$, by passage to the limit in
Gauss' identity.  Curvature convergence makes them flat there.

Choose a smooth cutoff $\chi$ equal to one near
$\overline B(0,r_0)$ and supported in $B(0,2r_0)$, and extend by
\[
 h_j=\chi B_j+(1-\chi)\delta .
\]
This is smooth globally, with
$\tfrac12|v|\leq|v|_{h_j}\leq\tfrac32|v|$.  It also retains
Gauss' identity globally:
\[
 h_j(x)(x,v)
 =\chi(x)\langle x,v\rangle+(1-\chi(x))\langle x,v\rangle
 =\langle x,v\rangle .
\]
Consequently $d_{h_j}(0,x)=|x|$: the radial segment has length
$|x|$, while Cauchy--Schwarz applied to this identity bounds
the derivative of Euclidean radius along any curve by its
$h_j$-speed.  Flatness and agreement with the coefficient
limit are retained on the inner ball; flatness of the transition
annulus is not needed.

On one common smaller spatial ball, the \emph{same} charts
converge to $h_j$ uniformly for $s\in[-1,0]$, with every
spatial derivative.  To prove this,
the flow equation and an upper coefficient bound $B$ give
\[
 \|\Phi_{k,j}^*G_k(s)-\Phi_{k,j}^*G_k(0)\|
       \leq 2n^3B\sup|\operatorname{Rm}_{G_k}|\longrightarrow0.
\]
All spatial derivatives are uniformly bounded by the backward
metric-jet argument.  If convergence of one derivative were not
uniform in $s$, choose offending times $s_k$ and points in a
compact inner ball.  Smooth compactness applied to those selected
spatial functions gives a subsequence converging smoothly; its
zeroth-order limit must be $h_j$ by the displayed estimate.
Uniqueness of derivatives contradicts the offending bound.
This proves static backward convergence without choosing a
different coordinate system at each time.

For clarity, we describe how these local limits fit together.
After a further common subsequence, all cross-chart distances have
locally uniform limits
\[
 d_{ij}(x,y)=\lim_k
       d_{G_k(0)}(\Phi_{k,i}(x),\Phi_{k,j}(y)).
\]
The distances are equicontinuous in both arguments by the chart
Lipschitz bounds, and are bounded because their centers have
radius one.  For $i=j$ the ambient-distance comparison on nested
normal balls identifies $d_{ii}$ with $d_{h_i}$.
Identify $(i,x)$ with $(j,y)$ exactly when $d_{ij}(x,y)=0$.
Triangle inequalities give transitivity, while the lower
within-chart distance bound gives uniqueness of $y$ in each chart.

This relation produces actual open overlap maps $T_{ij}$.
If $d_{ij}(x,y)=0$, keep $y$ in a compact ball strictly inside
its chart.  The lower distance bound separates the image of that
ball's boundary from its center.  Connectedness of intrinsic
balls and openness of the chart image put a uniform intrinsic
ball inside that image.  For large $k$, $\Phi_{k,i}(x)$ and
nearby points lie in this ball.  Thus
$T_{k,ij}=\Phi_{k,j}^{-1}\Phi_{k,i}$ is defined on a fixed
neighborhood and stays in a compact part of the target chart.
Any limit of its values must have zero $d_{ij}$-distance, so it
equals $T_{ij}$.  Uniform inverse bounds give continuity and
the inverse map; triangle inequalities give the cocycle identity.

The transition maps are smooth as well.  They are isometries
between the source coefficient metrics.  Ellipticity bounds
$DT_{k,ij}$, and their connection transformation law reads
\[
 \partial_a\partial_b T^\ell
   =\Gamma^c_{i,ab}\partial_cT^\ell
       -\Gamma^\ell_{j,uv}(T)\partial_aT^u\partial_bT^v.
\]
Here are the bounds needed to differentiate this identity.
Enlarge a compact source set slightly and retain all its target
images in one compact subset of the target chart.  Positivity
and coefficient convergence give one lower eigenvalue bound
$\beta>0$ and upper bound $C$ on these neighborhoods.  The
pullback identity gives
$\beta|DT_{k,ij}v|^2\leq C|v|^2$, hence
$|DT_{k,ij}|\leq\sqrt{C/\beta}$.
The coefficient matrices lie in a compact set of invertible
matrices, so all derivatives of matrix inversion are bounded
there.  Composition gives bounds for all inverse-metric jets.
The Koszul formula writes each Christoffel symbol as the inverse
metric contracted with first coefficient derivatives; it
therefore bounds all Christoffel jets.  Differentiating the
displayed equation $m$ times bounds derivatives of $T_{k,ij}$
of order $m+2$ using its already bounded derivatives through
order $m+1$.  This induction supplies every higher bound.
Smooth compactness and
the already unique value limit identify the smooth limiting
transition, its inverse, and its metric-preserving property.
The quotient of the chart pieces therefore has a smooth atlas.

We now realize these charts in the actual metric cone.  Choose a
common $0<a<\min(r_0,S_0,1/8)$ also smaller than the retained
distance-comparison and backward-jet radii.  On $\overline B(0,a)$ retain the distance
of the global metric $h_j$.  The normal radial identity places
all source images in the rescaled annulus $[1/2,2]$.
Choose cone points related to those source images by the annular
relations above.  Their within-chart distances tend uniformly
to $d_{h_j}$, and their cross-chart distances tend to $d_{ij}$.
Compactness of the cone annulus gives simultaneous pointwise
limits along a common free ultrafilter.  These limits are
isometries $e_j$ on the compact coordinate balls.  Convergence
is uniform on each ball: a finite net and the vanishing
distance distortion bound the error away from the net.
Thus one obtains
\[
 d_C(e_i(x),e_j(y))=d_{ij}(x,y),\qquad
 e_j(0)=(1,[\eta_j]).
\]
The second identity follows from the relation at the actual
source center on $\eta_j$, whose approximation error tends to
zero.  The relations retain normalized source radius exactly.
Since cone radius is $1$-Lipschitz, the actual source-radius
functions converge uniformly along that filter to $r\circ e_j$.
For each integer
$m$, retain an index making these errors less than $1/(m+1)$
on the first $m+1$ charts.  Indices can be chosen strictly
increasing because every such finite intersection belongs to
the free ultrafilter.  This yields one ordinary subsequence
retaining the radius limits and all earlier metric-jet limits.

Each $e_j$ covers the cone ball $B_C(e_j(0),a)$.  To see this,
fix $z$ in that ball and choose source points $y_k$ related to
$z$.  Use the original common ultrafilter for this coverage
argument.  The full annular distortion bound and convergence
at the center make $d_{G_k(0)}(q_{k,j},y_k)<a$ eventually
along that filter.  The full normal chart
therefore gives $y_k=\Phi_{k,j}(x_k)$ with $|x_k|<a$.
The cone approximations of these points tend to $z$, since
they are related to the same source points.  Uniform
convergence to $e_j$, and closedness of its compact image, put
$z$ in that image.  This argument takes place in the full cone:
$a<1/8$ and the $1$-Lipschitz radius keep the ball strictly
inside the annulus $[1/2,2]$.  If the backward-jet radius requires
a further common restriction, the equality
$d_{h_j}(0,x)=|x|$ retains this coverage at the smaller radius.

Use the open pieces $U_j=B(0,a/4)$ in the overlap quotient.
Their realizations are open embeddings and cover
$B_C(e_j(0),a/8)$.  Indeed, any point in this latter ball has
a preimage $x$ in the closed coordinate ball, and
$|x|/2\leq d_{h_j}(0,x)<a/8$ puts it in $U_j$.
For openness at $e_j(x)$, with $|x|<a/4$, choose a small cone
ball still inside $B_C(e_j(0),a)$.  Its preimages $y$ satisfy
$|y-x|\leq2d_C(e_j(y),e_j(x))$; shrinking that ball keeps
$|y|<a/4$.  This proves openness at every point of the image.

The cross-distance identity makes the $e_j$ descend to an
injective open map $f$ from the overlap quotient into $C(S)$,
with its existing quotient topology.  Thus $f$ is an open
embedding.  The centers $(1,[\eta_j])$ are dense in the unit
slice, and the same covered radius $a/8$ works for every
center, so the range contains the entire unit slice.
Its preimage $K$ is compact: it is the image of that compact
slice under the continuous inverse homeomorphism from
the range of $f$ to the quotient.  The quotient is now an
open subspace of a metric space, with its smooth local
charts; finitely many relatively compact chart neighborhoods
give a neighborhood $V$ of $K$ with compact closure.  The
following finite patching constructs actual source embeddings
on a possibly smaller such neighborhood.

We next need actual embeddings of a neighborhood of $K$ into
the source manifolds.  A finite chart cover suffices, but charts
cannot simply be identified using their limiting transitions:
the source transitions differ slightly.  Suppose $A_k$ is
already defined on an old neighborhood and a new chart has
quotient map $q$.  On their overlap its coordinate readout
$r_k=\Phi_k^{-1}A_kq$ tends smoothly to the identity.
Choose nested compact old regions $A\subset\operatorname{int}A'$
and a compact new chart piece.  Let $\chi$ be a fixed smooth
cutoff, equal to one near their overlap with $A'$ and supported
inside the readout domain.  The correction
\[
             a_k(x)=x+\chi(x)(r_k(x)-x)
\]
is eventually a diffeomorphism: its displacement has derivative
norm less than one, giving injectivity, and compact support
gives properness; the inverse function theorem makes its image
both open and closed.  All its derivatives converge to those
of the identity, and it preserves the chart domain because it
is the identity outside a compact subset of that domain.
Keep the old map on $\operatorname{int}A$ and use
$\Phi_k a_k$ on the new piece.  They agree on the overlap.

There are no extra collisions between the old and new pieces.
On their compact product, source distances converge uniformly
to quotient distances.  A hypothetical collision therefore
must approach the zero set, which lies inside the larger old
region $A'$.  In that region the corrected chart agrees with
the old map exactly, and its injectivity identifies precisely
the same quotient point.  Finite induction gives smooth open
embeddings $A_k$ on a fixed relatively compact neighborhood $V$
of $K$.  In each retained chart, their inverse-coordinate
readouts converge smoothly to the identity.  Pulling back by
these readouts retains the static metric convergence on
$[-1,0]$, uniformly on every compact subset of $V$.
The same local formulas and a finite cover give convergence
of normalized basepoint radii and all pairwise distances on
compact subsets to the corresponding radius and distance in
$C(S)$.

These embeddings contain a whole thin source annulus.
To check this stronger coverage, first take a finite
$\varepsilon/2$-net of ray directions in $S$.
Uniform sphere-to-link distortion makes the same finite rays
an $\varepsilon$-net of every sufficiently large source sphere
after division by its radius.  If
$|d(p,x)/L-1|\leq\delta$, use the ray net at radius $d(p,x)$
and then move along that ray to radius $L$; the extra normalized
error is at most $\delta$.  Choosing
$\delta\leq\min(1/2,\varepsilon/4)$ gives a finite net for
the entire rescaled annulus, with centers on the radius-$L$
ray evaluations.
Compactness of $K$ supplies one coordinate radius inside any
chosen smaller neighborhood $W$ of $K$ in $V$.  The lower
chart distance bound separates its boundary, and the corrected
maps are uniformly close to the source charts.  Their open
images therefore contain a fixed smaller ball about each of
the finitely many source centers.  Choose the annular net finer
than that radius.  For all sufficiently large $k$,
\[
 \{x:|d_{G_k(0)}(p,x)-1|\leq\delta\}\subset A_k(W).
\]
This prevents loss of part of the unit slice during patching.

We next give $S$ its smooth metric.  In every retained cone
chart the function $u=r^2/2$ has the geodesic quadratic
interpolation identity proved above.  The same Rademacher and
normal-coordinate argument makes $u$ smooth with
$\nabla^2u=h_j$.  The radial curve $a\mapsto(a,\xi)$ has unit
speed in the cone distance and realizes the local radial
variation.  The metric-arc argument below makes this curve a
smooth geodesic inside each retained chart, before any smooth
dilation action is assumed.  Differentiating $u$ along it, and using that radius
is $1$-Lipschitz, gives $|\nabla u|^2=2u$.
In particular $du\ne0$ near $u=1/2$.  The implicit function
theorem gives smooth $(n-1)$-dimensional charts on this level,
whose images in the metric cone are exactly open subsets of $S$.

The transitions are smooth for a reason independent of any
preassigned smooth structure on $S$.  Ambient chart transitions
preserve Riemannian distances, since both compute the same cone
distance.  Here is the regularity argument for such a map
$F$ between open Riemannian regions.  Intrinsic balls give
the manifold topology, so distance preservation first gives
continuity.  On sufficiently small coordinate neighborhoods,
freezing each positive metric gives two-sided comparisons
between intrinsic distance and a fixed Euclidean norm.
The upper comparison follows by a straight coordinate path;
for the lower one, paths leaving a larger coordinate ball
already have a fixed positive length before their first exit.
Shrinking the inner ball therefore retains the lower bound
even for the unrestricted intrinsic distance.  Combining the
two comparisons makes $F$ Euclidean locally Lipschitz.

Fix a point $p$ where smoothness is required.  A smooth
geodesic flow near zero initial velocity gives an endpoint
map $E(q,v)$ and a smooth map
\[
             (q,v)\longmapsto(q,E(q,v)).
\]
Its derivative in $v$ at $(p,0)$ is the identity.  Its full
derivative has triangular blocks with identity diagonal,
so the inverse function theorem gives a smooth local inverse
near $(p,p)$.  The chosen geodesics remain inside the domain
of $F$ on a common time interval containing $[0,1]$.
Rademacher provides a differentiability point $q$ close enough
to $p$ that $(q,p)$ lies in this inverse neighborhood.
The inverse then writes every nearby $y$ as $E(q,v(y))$,
with $v(y)$ smooth.

Distance preservation sends each short affine geodesic to a
smooth affine geodesic; the normal-coordinate argument for
metric arcs given below proves this without assuming $F$
smooth.  The image geodesic has initial point $F(q)$ and
initial velocity $dF_qv$.  Smooth dependence on these initial
data therefore makes its time-one endpoint smooth in $v$.
That endpoint is $F(E(q,v))$.  Composing with $v(y)$ proves
$F$ smooth near $p$, including when $F$ was not initially
known differentiable at $p$.  Apply the same argument to
the inverse transition.  Finally, preservation of distances
along short minimizing geodesics preserves their speeds,
so $|dF_pv|=|v|$; polarization gives tensor preservation.
These ambient
transitions preserve $u$, and their restrictions to its regular
level are therefore smooth and preserve the induced metric.
The local level metrics consequently agree on overlaps and
define a smooth metric $g_S$ on the actual compact space $S$.
For a positive-radius point $a$, normalize to $a/r(a)$ and
choose one of the unit-slice radial charts $H$ there.
Then $\delta_{r(a)}\circ H$ is an open chart around $a$;
its metric and potential are multiplied by $r(a)^2$.
These charts cover the positive cone.  Their distance-preserving
transitions identify the local tensors and radial potentials,
giving a smooth metric on that cone.  Its radial
potential satisfies
\[
                 \nabla^2u=g_C,\qquad |\nabla u|^2=2u.
\]
The realization of the patched neighborhood into this positive
cone is a local diffeomorphism preserving these metrics: compare
its retained distance charts with the radial charts just used.
Thus the inverse image $b:S\to V$ is a smooth immersion with
compact image $K$, $b^*g_V=g_S$, and $u\circ b=1/2$.
Restriction to any smaller open neighborhood retains the
Hessian and gradient identities.

We will need that $S$ is connected, including when $n=2$.
First extend the approximation relations through the vertex.
Fix a normalized radial bound $R$ and an error $\varepsilon>0$.
For a source point with $d(p,x)/L<\varepsilon/16$, use any
fixed ray at parameter $d(p,x)$; its normalized error is at most
$2d(p,x)/L<\varepsilon/8$.  Outside this small region use
uniform ray approximation with error $\varepsilon/4$.
Uniform homogeneous ray-distance convergence on $[0,R]^2$,
including its zero-radius edges, bounds the remaining
pairwise error by $\varepsilon/2$.  The sum of the two
approximation errors and this distance error is less than
$\varepsilon$.  Every cone point has an exact ray representative
in the source at scale $L$.  Thus these relations cover both
whole bounded radial balls and retain radius exactly.

The cone now has actual minimizing metric segments.  Approximate
two cone points $x,y$ by source points and split a minimizing
source geodesic at distance $Lr$, where $0<r<d_C(x,y)$.
Small enough distortion ensures that this distance is available.
The split point has normalized radius at most $r(x)+r$,
so all its cone approximations lie in one compact radial region.
Their distances from the endpoints tend to $r$ and
$d_C(x,y)-r$.  A compact limit gives this exact distance split;
the endpoint cases are immediate.

Repeated splitting constructs an $(N+1)$-point chain from
$x$ to $y$ with every pairwise distance equal to its index
difference times $d_C(x,y)/N$.  The triangle inequality gives
both bounds for this pairwise assertion from the endpoint
distance and the adjacent step lengths.  At time $t\in[0,1]$
take the chain point with index $\lfloor Nt\rfloor$.
These step functions all lie in the compact closed ball about
$x$ of radius $d_C(x,y)$.  Take their pointwise limits along
one free ultrafilter.  Since $\lfloor Nt\rfloor/N\to t$, the
limits satisfy
$d_C(\gamma(s),\gamma(t))=|s-t|d_C(x,y)$ and have endpoints
$x,y$.  In particular $\gamma$ is continuous and minimizing.
For unit endpoints $x,y$, the radial variation identity gives
\[
       r(\gamma(t))^2=1-t(1-t)d_S^{\rm chord}(x,y)^2.
\]
If $d_S^{\rm chord}(x,y)<2$, its radius is bounded below by
$\sqrt{1-d_S^{\rm chord}(x,y)^2/4}>0$.
Radially normalizing $\gamma$ therefore gives a continuous path
in $S$ from $x$ to $y$.

For antipodal endpoints, represented by rays $\alpha,\beta$,
the limiting chord distance two and corresponding-side
comparison force
$d(\alpha(L),\beta(L))=2L$ at every $L>0$.
Thus $p$ is a minimizing midpoint.  Squared-distance
semiconcavity from Toponogov gives, for any third ray $\zeta$,
\[
 d(\zeta(L),\alpha(L))^2+d(\zeta(L),\beta(L))^2
       \leq 2d(\zeta(L),p)^2+2L^2=4L^2.
\]
Divide by $L^2$ and pass to the limit.  Since $S$ has positive
dimension, a coordinate neighborhood of $x$ contains a point
$z\ne x$ with $d_S^{\rm chord}(x,z)<2$.  The resulting inequality
\[
 d_S^{\rm chord}(x,z)^2+d_S^{\rm chord}(y,z)^2\leq4
\]
forces $d_S^{\rm chord}(y,z)<2$.  The two normalized segment
paths join $x$ to $y$ through $z$.  Hence $S$ is path connected.
This uses comparison in the original nonnegatively curved
manifold; it is not a property of every abstract metric cone.

To prevent a flat quotient link from being mistaken for a sphere,
one also uses backward enclosing normal charts.  They yield a
Euclidean realization $\iota:S\to\mathbb R^n$ with a smooth unit
normal $\nu$ satisfying $d\nu=d\iota$ and
$\iota^*g_{\rm Eucl}=g_S$.  Terminal distance convergence and
backward metric monotonicity retain separation:
$d_S^{\rm chord}(x,y)\leq|\iota(x)-\iota(y)|$.
We construct this realization from the complete bounded ancient flow,
nonnegative curvature operator, fixed parabolic noncollapse, and
zero-ratio tail bound.  Here the source dimension is $n\geq2$,
the unit slice has dimension $n-1$, and the target is $\mathbb R^n$.
The enclosing step must retain the whole link.
Write $g_V$ for the metric on the fixed annular neighborhood and
$g_k=A_k^*G_k(-1)$.  The established local coefficient convergence
gives $g_k\leq2g_V$ on $K=b(S)$ for all sufficiently large $k$.
Indeed, on each compact coordinate subset the least eigenvalue
of $g_V$ is positive, so sufficiently small coefficient error
implies this inequality; take a finite cover of $K$.
Since $b^*g_V=g_S$, the maps $A_kb$ have uniformly bounded
differential with respect to $g_S$.  Integrating along paths in
the connected compact Riemannian manifold $(S,g_S)$ bounds the
earlier-time distances between all their images by a fixed
multiple of its finite intrinsic diameter.

This bound extends to a fixed open neighborhood $W$ of $K$.
Around each point of $K$ choose a smaller convex coordinate
ball with compact closure.  Uniform upper coefficient bounds
give a uniform upper length bound for every coordinate segment
from its center.  Cover $K$ by finitely many such balls and
add these bounds to the preceding image-diameter bound.
Shrinking within $V$ if necessary, we obtain a relatively
compact $W$ and a number $D>0$ such that, with a fixed $z_0\in S$,
\[
 d_{G_k(-1)}(A_kx,A_kb(z_0))\leq D\qquad(x\in W)
\]
eventually, with one threshold for all $x$.

Set $q_k=A_kb(z_0)$ and $\epsilon_k=R_{G_k}(0,q_k)\to0$.
For $a<b\leq0$, bounded ancient Harnack gives the precise
earlier-distance estimate
\[
 |\operatorname{Rm}_{G_k(a)}|(x)
 \leq n^2R_{G_k}(b,q_k)
       \exp\!\left(\frac{d_{G_k(a)}(x,q_k)^2}{2(b-a)}\right).
\]
Each flow separately has bounded curvature on the whole past;
no common global curvature bound for this sequence is required.
Taking $a=-1$, $b=0$ proves uniform curvature decay on every
fixed radius ball in $G_k(-1)$.  Scalar time monotonicity and
the nonnegative curvature operator extend the bound to all
earlier times on that same spatial ball.  Thus the curvature
antecedent of parabolic noncollapse holds there eventually.
Fix the positive normal-coordinate radius $\rho<1$ determined
by dimension $n$, curvature bound one on a radius-two ball,
and volume lower bound $\kappa$ on its concentric unit ball.
These are fixed numerical inputs to the local injectivity
estimate; the resulting radius is independent of the scale
chosen next.  A further fixed scale factor
\[
                 c=\left(\frac{\rho}{16(|D|+1)}\right)^2
\]
places the entire earlier image of $W$ inside the radius-$\rho/8$
ball for the metric $cG_k(-1)$.
More explicitly use $H_k(s)=cG_k(-1+s/c)$; its time-zero
radius-two ball corresponds to the fixed earlier radius
$2/\sqrt c$ before this additional scaling.  The preceding
whole-past decay therefore supplies curvature at most one on
that ball for all sufficiently large $k$.  Noncollapse and the
normal-chart construction give charts $\Psi_k$ on the Euclidean
ball of radius $\rho$, whose images are exactly the intrinsic
balls of that radius.

Their terminal coefficients converge smoothly on radius
$\rho/4$ to the Euclidean coefficients.  To identify the limit,
use the normalized radial Jacobi fields, not just flatness of
the limiting metric.  If terminal curvature is at most
$\varepsilon$ along a radial segment to $x$, the Jacobi integral
equation gives the error bound
\[
 |(\Psi_k^*H_k(0))_x(w,w)-|w|^2|
 \leq\delta(2+\delta)|w|^2,\qquad
 \delta=\frac{\varepsilon|x|^2
                 e^{\max(1,\varepsilon|x|^2)}}6.
\]
The already established derivative bounds give a smooth
coefficient subsequence.  Letting $\varepsilon\downarrow0$
identifies all its diagonal values with the Euclidean ones;
polarization identifies the bilinear form.  Smooth convergence
then identifies every derivative with the derivative of this
constant coefficient field.

Consider the actual maps $J_k=\Psi_k^{-1}A_k$ on $W$.
The exact radial-distance identity confines all their images
to $B(0,\rho/8)$, a compactly contained subset of the region
where the target coefficients converge.  In any source chart,
$J_k$ preserves the two actual pullback metrics.
Ellipticity bounds $DJ_k$, and the connection transformation
formula bounds its second derivative.  Differentiating that
formula inductively bounds every higher derivative using the
convergent source and target metric jets.  A diagonal over a
countable source atlas and compact subsets yields smooth local
limits.  On overlaps these limits agree, since for each $k$
they are restrictions of the single map $J_k$.
They therefore give one smooth map $J:W\to\mathbb R^n$ with
\[
                   J^*g_{\rm Eucl}=c g_V.
\]
Positivity makes its differential injective; equal dimensions
and the inverse function theorem make
$j=c^{-1/2}J$ a local isometry.

The normal is explicitly
$\nu(z)=dj_{b(z)}(\nabla u(b(z)))$, and
$\iota=j\circ b$.  The identities for $u$ give
$|\nu|=1$ and $\nabla_v\nabla u=v$.
Preservation of the Levi-Civita connection by the local isometry
then gives $d\nu=d\iota$.
Moreover $du(db(v))=0$, so $\nu$ is orthogonal to the image
of $d\iota$; and $\iota^*g_{\rm Eucl}=g_S$.
This construction uses local inverses of $j$ and does not
presuppose global injectivity.

For injectivity fix $x,y\in S$.  Their terminal source distances
tend to $d_S^{\rm chord}(x,y)$, and Ricci nonnegativity gives
\[
 \sqrt c\,d_{G_k(0)}(A_kb(x),A_kb(y))
       \leq d_{cG_k(-1)}(A_kb(x),A_kb(y)).
\]
The coordinate endpoints lie in $B(0,\rho/8)$.
For each $\eta>0$, eventual coefficient convergence bounds
the metric by $(1+\eta)g_{\rm Eucl}$ on the closed ball of
radius $\rho/6$.  The straight segment between the coordinate
endpoints lies in that ball, so the right side is at most
\[
 \sqrt{1+\eta}\,|J_kb(x)-J_kb(y)|.
\]
Taking limits and then $\eta\downarrow0$, and dividing by
$\sqrt c$, proves
$d_S^{\rm chord}(x,y)\leq|\iota(x)-\iota(y)|$.
Only this explicit comparison path is needed; no confinement
of an ambient minimizing geodesic is assumed.
Thus $\iota$ is injective.  Since $d(\iota-\nu)=0$,
$\iota-\nu$ is constant on the connected link.
The map $\nu:S\to S^{n-1}$ is a local diffeomorphism, has closed
image by compactness and open image by its invertible differential,
and is injective because $\iota$ is injective.  It is therefore
a metric-preserving diffeomorphism onto the whole unit sphere.
This works also for the one-dimensional link, without assuming
simple connectedness.

It remains to identify the chord metric, since the preceding
isometry concerns the induced Riemannian metric.  Put
\[
 \alpha(x,y)=\arccos(1-d_S^{\rm chord}(x,y)^2/2).
\]
This function satisfies the triangle inequality.  Write
$A=\alpha(x,z)$, $B=\alpha(z,y)$.  If either is zero or
$A+B\geq\pi$ the claim is immediate.  Otherwise use the
intermediate cone point $(r,z)$, where
\[
             r=\frac{\sin(A+B)}{\sin A+\sin B}>0.
\]
The cone cosine law computes the sum of its distances to
$(1,x)$ and $(1,y)$ as $\sqrt{2-2\cos(A+B)}$.
The cone triangle inequality bounds the chord $xy$ by this
number, and monotonicity of cosine on $[0,\pi]$ gives the claim.
Also $\alpha=2\arcsin(d_S^{\rm chord}/2)$, so for every
$c_1>1$ there is $\delta>0$ such that
$\alpha\leq c_1d_S^{\rm chord}$ on chords below $\delta$.
Chord distance is bounded by smooth link-path length:
subdivide a path into regular-level charts, where chord
distance is ambient distance and the level metric is induced,
and sum the local inequalities.  Subdivide further until all
chords are below $\delta$.  The angular triangle inequality
bounds the endpoint angle by $c_1$ times the path length.
Let $c_1\downarrow1$ and take the infimum over paths to obtain
$\alpha(x,y)\leq d_{g_S}(x,y)$.

For the reverse inequality we first justify the differential
geometry of the actual cone coordinates.  An arc $\gamma$
in a smooth Riemannian chart satisfying
\[
 d(\gamma(s),\gamma(t))=c|s-t|
\]
on an open interval is a smooth geodesic, even when the chart
metric is incomplete.  For $c=0$ it is constant.  Otherwise
take normal coordinates at $\gamma(0)$ on a smaller interval
and write $v(t)=\exp_{\gamma(0)}^{-1}\gamma(t)$.
The radial distance formula gives $|v(t)|=c|t|$.
For $s<0<t$, the two radial legs have total length exactly
$d(\gamma(s),\gamma(t))$.  Their unit initial directions must
be opposite.  Indeed distance additivity and the triangle
inequality force the endpoints of equally short radial
sublegs of length $\eta$ to have distance $2\eta$.
For each $\theta>1$, continuity of the normal-coordinate
metric bounds this distance above by
$\theta\eta|u_2-u_1|$ when $\eta$ is sufficiently small,
by integrating along the coordinate chord.  Cancel $\eta$
and let $\theta\downarrow1$.  Since $u_1,u_2$ are unit,
$|u_2-u_1|\geq2$ implies $|u_1+u_2|^2=0$.
Fixing one negative parameter makes all positive directions
equal; fixing one positive parameter does the same on the
negative side.  Thus $v(t)=tw$ for one $w$, and
$\gamma(t)=\exp_{\gamma(0)}(tw)$.  Translation of the parameter
proves smoothness and the geodesic equation throughout the
interval.  This applies to interior subarcs of the minimizing
cone segment in every retained chart.

Write $\delta_a$ for cone dilation by $a>0$.  In one retained
chart consider $O(t,x)=\delta_{1+t}x$ for $t$ near zero,
interpreted in that chart.  Continuity of dilation gives an
open neighborhood $U$ of any chosen $x$ and a common
$\epsilon>0$ on which $O$ remains in the chart for
$(t,y)\in(-\epsilon,\epsilon)\times U$.
The cone distance formula and $u=r^2/2$ give
\[
 d(O(s,y),O(t,y))=r(y)|s-t|,\qquad
 u(O(t,y))=(1+t)^2u(y).
\]
The preceding arc argument makes each orbit a smooth geodesic.
If $v=\partial_tO(0,y)$, the normal-coordinate distance formula
gives
$|v|=\lim_{t\to0}d(y,O(t,y))/|t|=r(y)$.
Differentiating the smooth potential identity then gives
\[
 |v|^2=2u(y),\qquad
 \langle\nabla u,v\rangle=2u(y)=|\nabla u|^2.
\]
Consequently $|v-\nabla u|^2=0$, so the initial velocity is
exactly $\nabla u(y)$.  Joint smoothness now follows from the
geodesic equation: for small $(t,y)$ the curve
$s\mapsto O(ts,y)$ is defined on $-1<s<2$ and has initial
data $(y,t\nabla u(y))$, which depend smoothly on $(t,y)$.
Uniqueness identifies it with the smooth geodesic solution
for those data; evaluating at $s=1$ proves smoothness of $O$.
At $(0,y)$ its joint differential is
$(a,v)\mapsto a\nabla u(y)+v$.

Restrict the starting point to the regular level $S=\{u=1/2\}$
and set $P(r,z)=\delta_r z$.  At $(1,z)$ the radial derivative
is $\nabla u(z)$ and the angular derivative is the level
inclusion.  The eikonal identity gives unit radial norm,
and $du=0$ on level tangents gives orthogonality.
For any fixed $c>0$, dilation is a smooth homothety between
retained charts: give the source chart the metric $c^2g$;
the dilation map preserves distances, so the local
distance-isometry argument used for chart transitions gives
smoothness and equality of tensors.  The identity
$P(r,z)=\delta_cP(r/c,z)$ transfers joint smoothness from
radius one to radius $c$; its radial derivative has the
factor $1/c$, which cancels the homothety's factor $c$.
Thus at $(c,z)$,
\[
 (P^*g)_{(c,z)}((a,v),(b,w))=ab+c^2g_S(v,w).
\]
This positive form makes $dP$ injective, hence invertible
in equal dimensions.  The inverse function theorem gives
smooth local polar coordinates about every positive radius.
Their inverse is exactly the cone radius and
$\delta_{1/r}$-normalization, by the dilation identities.
These inverses agree on overlaps, so the tensor is
$dr^2+r^2g_S$ wherever these polar coordinates are used.

Now let $D=d_S^{\rm chord}(x,y)<2$ and take the minimizing
cone segment between the unit-radius endpoints.

Write $\theta(t)$ for the normalized segment.  Its radius
satisfies $r(t)^2=1-t(1-t)D^2$, and its speed is $D$.
Differentiation and the polar tensor give
\[
 |\theta'(t)|_{g_S}
       =\frac{D\sqrt{1-D^2/4}}{1-t(1-t)D^2}.
\]
The integral from zero to one is
\[
 2\arctan\!\left(\frac{D}{2\sqrt{1-D^2/4}}\right)
               =\arccos(1-D^2/2)=\alpha(x,y).
\]
Use interior subarcs and send their endpoints to zero and one.
Continuity in the link and the displayed integrable speed
suffice; no endpoint differentiability is assumed.
This gives $d_{g_S}(x,y)\leq\alpha(x,y)$.
For $D=2$, choose $z_j\ne y$ converging to $y$ in a level
chart.  The antipodal comparison forces
$d_S^{\rm chord}(x,z_j)<2$; the bounds
$d_{g_S}(x,z_j)\leq\pi$ pass to the limit in the manifold
topology.  Thus $d_{g_S}(x,y)\leq\pi=\alpha(x,y)$ as well.
We have proved the angular-distance formula for every pair.

The spherical isometry preserves intrinsic distances, whose
round-sphere formula is the same arccosine of chord distance.
Taking cosine recovers equality of the nonnegative chords.
The radial extension $(r,\xi)\mapsto r\nu(\xi)$ therefore
preserves cone distances.  It is surjective by normalizing
any nonzero Euclidean vector, and maps the vertex to zero.
Hence the actual metric cone is Euclidean and its radial
unit ball has normalized Hausdorff measure $\omega_n$.

For comparison back to source volume, define a projection on the
ray range.  For $L>0$ its value is
$\pi_L(\gamma(rL))=(r,[\gamma])$.  Corresponding-side comparison gives
\[
 d_C(\pi_L(a),\pi_L(b))\leq L^{-1}d_{t_0}(a,b),
\]
which also makes the map independent of any coincident ray
representations.  Its restriction to the ray union inside
$B(p,L)$ maps onto the whole radial cone unit ball.

Here normalized Hausdorff measure agrees with the coordinate
Riemannian volume form.  To see the normalization and the local
comparison, freeze the metric coefficients at a chart point
$z_0$ and choose a linear map $T$ with
$|Tv|^2=g_{z_0}(v,v)$.  Then
$|\det T|=\sqrt{\det g_{z_0}}$.
For every $\lambda>1$, shrink the coordinate neighborhood so
that intrinsic distances are within factors $\lambda$ of
$|T(z-w)|$, and the volume density is within factors
$\lambda$ of $|\det T|$.  The lower distance bound includes
paths leaving the chart, by the first-exit estimate on a
larger coordinate neighborhood.  The Lipschitz inequalities
for $n$-dimensional Hausdorff measure and linear change of
variables compare the two measures within factors
$\lambda^{n+1}$ on every measurable subset of this
neighborhood.  A countable disjoint refinement of these
neighborhoods gives the same comparison globally.
Letting $\lambda\downarrow1$ proves equality.  The Euclidean
normalization is Lebesgue measure, so its unit ball has
measure $\omega_n$.

The Lipschitz inequality for Hausdorff outer measure applies to
this subset without requiring it to be measurable or of full
measure.  That inequality, monotonicity under inclusion in the
source ball, and equality
of normalized Hausdorff and Riemannian volume in the source,
give
\[
 \omega_n=\mathcal H_C^n(\{r<1\})
       \leq L^{-n}\operatorname{Vol}_{g(t_0)}B(p,L).
\]
Bishop--Gromov gives the reverse upper bound.  Changing centers
shifts ball radii by at most their fixed distance and hence
preserves the asymptotic ratio.  Monotonicity between the
Euclidean small-radius and infinite-radius limits now gives
$\operatorname{Vol}B(x,r)=\omega_n r^n$ for every $x,r>0$.

We now use $n\geq2$, completeness and the nonnegative curvature
operator to deduce local rigidity from these exact ball volumes.
In a normalized normal neighborhood of $x$, let
$\rho(v)>0$ be the exponential volume density; $\rho(0)=1$.
Along $\gamma(t)=\exp_x(t\theta)$, with $|\theta|=1$, use
parallel orthonormal frames.  The radial variation differential
has matrix $J(t)=P_t^{-1}\,t\,d\exp_x|_{t\theta}$, where
$P_t$ is parallel transport from $x$; it satisfies
$J''=-\mathcal R J$, $J(0)=0$, and $J'(0)=I$.
The radial curvature operator is symmetric and annihilates
the radial vector, so restricting to its orthogonal complement
retains the Jacobi equation and the full Ricci trace.
Write $A(t)$ for this transverse Jacobi matrix, of dimension
$m=n-1>0$.  For $t>0$ before singularity it is invertible and
\[
 |\det A(t)|=t^m\rho(t\theta).
\]
Indeed the full differential has radial eigenvalue $t$;
factoring it from the determinant leaves this transverse
power.  Absolute determinant avoids any orientation choice.

The Wronskian $A^TA'-(A')^TA$ has zero derivative, by the
Jacobi equation and symmetry of curvature, and vanishes at
zero.  Thus $B=A'A^{-1}$ is symmetric.  Differentiating the
inverse identity gives $B'=-B^2-\mathcal R_\perp$.
Put $h=\operatorname{tr}B$ and
$y=|\det A|^{1/m}=t\rho(t\theta)^{1/m}$.  Determinant
differentiation gives $y'/y=h/m$, hence
$y''/y=h'/m+(h/m)^2$.
Since
$h'=-\operatorname{tr}(B^2)-\operatorname{Ric}(\dot\gamma,\dot\gamma)$
and $\operatorname{tr}(B^2)\geq h^2/m$, we obtain
\[
 \frac{d^2}{dt^2}\bigl(t\rho(t\theta)^{1/m}\bigr)
 \leq-\frac{\operatorname{Ric}(\dot\gamma,\dot\gamma)}m
                           t\rho(t\theta)^{1/m}.
\]
The expression $y=t\rho(t\theta)^{1/m}$ is smooth through
zero, with $y(0)=0$ and $y'(0)=1$; no fractional power of a
vanishing determinant is differentiated there.  Ricci
nonnegativity gives $y''\leq0$.  More explicitly,
$(ty'-y)'=ty''\leq0$, so $y/t$ is nonincreasing for $t>0$,
and its limit one at zero yields $\rho\leq1$.
Choose a normal ball on which the exponential map is injective
and covers the intrinsic ball of the same radius.
Change of variables and exact ball volume give
$\int(1-\rho)=0$.  If $\rho<1$ at a point, continuity gives
a positive deficit on an open Euclidean subset of positive
measure, a contradiction.  Hence $\rho=1$ throughout.
With $\rho=1$ the left side of the displayed differential
inequality is zero, so radial Ricci is
nonpositive.  Let $t\downarrow0$ and use Ricci nonnegativity:
Ricci vanishes in every direction at $x$.  Its scalar trace
is zero, and the nonnegative curvature operator is therefore
zero.  The slice is flat, contradicting positivity propagated
by bounded Harnack.
\end{proof}

\subsection{A nonflat ancient limit with a line}

\begin{lemma}[Selected line limit]
\label{lem:ancient-selected-line}
Let $g(t)$ satisfy the hypotheses of
Proposition~\ref{prop:bounded-ancient-zero-volume}, and suppose
$\operatorname{AVR}(g(t_1))>0$ for some $t_1\leq0$.
Then there are $\delta>0$ and a Ricci flow $G(s)$, $s<\delta$,
on a connected $n$-manifold, with complete slices, nonnegative
curvature operator, $|\mathrm{Rm}_G|\leq4n^2$, and a point $o$
such that $R_G(0,o)>0$ and $\operatorname{AVR}(G(0))>0$.
The terminal slice $G(0)$ contains an isometric line
$\gamma:\mathbb R\to M_\infty$ with $\gamma(0)=o$.
\end{lemma}
\begin{proof}
Set $v=\operatorname{AVR}(g(t_1))>0$ and $t_0=t_1-1<0$.
Lemma~\ref{lem:ancient-infinite-scalar-ratio} gives unbounded
distance-squared scalar curvature at $t_0$.  We perform the spatial
selection at $t_0$ and keep the later positive volume at $t_1$.
All unmarked distances in the next paragraph are for $g(t_0)$.

Write $R(x)=R(t_0,x)$ and choose $C\geq1$ with $0\leq R\leq C$.
For each integer $i\geq1$, unboundedness supplies $x_i$ with
$d(p,x_i)\sqrt{R(x_i)}>4(i+1)^2\sqrt C$.
Apply doubling selection to $\sqrt R$, with displacement budget
$d(p,x_i)/2$.  Explicitly, if the required bound fails at a
candidate, move by less than $L_i^*/\sqrt R$ to a point where
$\sqrt R$ more than doubles, with
$L_i^*=d(p,x_i)\sqrt{R(x_i)}/4$ held fixed.
The sum of these displacements is bounded by the geometric
series $2L_i^*/\sqrt{R(x_i)}=d(p,x_i)/2$.
Global boundedness forces this procedure to terminate.
It yields $q_i$ and $\rho_i>0$ with
\[
 R(q_i)\geq R(x_i),\quad d(p,q_i)\geq\tfrac12d(p,x_i),\quad
 \rho_i\sqrt{R(q_i)}=L_i^*,\quad
 R\leq4R(q_i)\ \hbox{on }B(q_i,\rho_i).
\]
The displacement budget also gives $\rho_i\leq d(p,q_i)$.
Reduce the selected radius to $r_i=\rho_i/(i+1)$ and set
\[
 Q_i=R(q_i)>0,\quad d_i=d(p,q_i),\quad
 L_i=r_i\sqrt{Q_i},\quad D_i=d_i\sqrt{Q_i}.
\]
Then $d_i,r_i,L_i,D_i\to\infty$, $r_i/d_i\to0$, and $Q_i\leq C$.
The scalar values themselves need not diverge.
Bounded ancient Harnack gives $R(s,x)\leq R(t_0,x)$ for
$s\leq t_0$; nonnegative curvature operator gives
$|\mathrm{Rm}|\leq n^2R$.  Thus the actual rescalings
\[
 H_i(s)=Q_i g(t_0+s/Q_i),\qquad s<-t_0Q_i,
\]
satisfy $R_{H_i}(0,q_i)=1$ and, on the fixed terminal balls
$B_{H_i(0)}(q_i,L_i)$ throughout $s\leq0$,
\[
 0\leq R_{H_i}(s)\leq4,\qquad
 |\mathrm{Rm}_{H_i}(s)|\leq4n^2.
\]
Take $\eta=(2n)^{-1}\leq1$.  The bound $4n^2=\eta^{-2}$
verifies the parabolic antecedent on the terminal $\eta$-ball,
so the terminal unit ball has volume at least
$\kappa\eta^n$.  The expanding-cylinder compactness construction
from the Ricci-flow chapter applies, with arbitrarily long past
slabs and $L_i\to\infty$.  Its scalar derivative estimate at
the base chooses a common $0<\delta<1$ so that the earlier
basepoint curvature stays positive.  It gives a complete pointed
limit $G(s)$ of $H_i(s-\delta)$ for $s<\delta$, with
nonnegative curvature operator, $|\mathrm{Rm}_G|\leq4n^2$,
and $R_G(0,o)>0$.

Since nonnegative Ricci and scalar at most four imply
$\operatorname{Ric}_{H_i}\leq4H_i$ on these controlled balls,
Lemma~\ref{lem:ancient-geometric-selected-line} applies to the
selected $g(t_0)$, $Q_i,r_i,q_i$ and the same buffer $\delta$.
It gives an isometric line through $o$ in this retained $G(0)$.

The same pointed limit has positive asymptotic volume.
The bounded ancient volume comparison is antitone in time,
so at every source time $t_0-\delta/Q_i\leq t_1$ the volume
ratio is at least $v$.  Bishop--Gromov and center comparison
give, at every radius $r>0$,
\[
 \operatorname{Vol}_{H_i(-\delta)}B(q_i,r)
                    \geq (v/2^n)r^n.
\]
The complete-limit ball coverage and determinant comparison
developed in the Ricci-flow chapter pass this lower bound to
$G(0)$.  In particular $\operatorname{AVR}(G(0))>0$.
No further limit is chosen for this volume argument.
\end{proof}

\subsection{A line gives a parallel gradient}

\begin{lemma}[Busemann splitting function]
\label{lem:ancient-line-busemann}
Let $(M,g)$ be a connected complete Riemannian manifold of
dimension $n\geq2$, with nonnegative Ricci curvature, and let
$\gamma:\mathbb R\to M$ be an isometric line.  Then
\[
 f(x)=\lim_{T\to\infty}\bigl(T-d(x,\gamma(T))\bigr)
\]
is smooth, $f(\gamma(0))=0$, $|\nabla f|=1$ and $\nabla^2f=0$.
\end{lemma}
\begin{proof}
Put $o=\gamma(0)$ and define the two oriented limits
\[
 b_\pm(x)=\lim_{T\to\infty}\bigl(T-d(x,\gamma(\pm T))\bigr).
\]
The approximants increase with $T$ and are bounded in absolute
value by $d(x,o)$; triangle inequalities prove both assertions.
Their limits are $1$-Lipschitz, satisfy
$b_+(\gamma(s))=s$ and $b_-(\gamma(s))=-s$, and obey
$b_++b_-\leq0$ everywhere.  Distance Laplacian comparison
gives subharmonicity in two useful forms.  Moving the distant
pole along a minimizing segment gives smooth lower supports
for the distance approximants with Laplacian at least
$-2(n-1)/d(x,\gamma(\pm T))$.  On a fixed compact set this
error tends uniformly to zero.  Add the two actual supports
to treat the sum of approximants.  If a smooth function has
strictly negative Laplacian and bounds the limiting sum on
the boundary of a compact set, it also bounds each approximant
there, because the approximants lie below the limit.
For sufficiently large $T$, a positive interior maximum of
the approximant minus that function contradicts the support
Laplacian inequality.  Passing to the limit gives compact
comparison.  Finally, an upper test with negative Laplacian
could be lowered by a sufficiently small smooth bump near
its contact point while retaining negative Laplacian and the
boundary bound, contradicting comparison.  Thus every smooth
upper test of the sum has nonnegative Laplacian.

The strong maximum argument for this continuous sum is local.
On a compact coordinate annulus, the squared coordinate radius
$\rho$ has $|\nabla\rho|^2$ bounded below positively.  For
large $\alpha$,
\[
 \Delta e^{-\alpha\rho}
   =e^{-\alpha\rho}(\alpha^2|\nabla\rho|^2-\alpha\Delta\rho)>0.
\]
If a nonpositive function with the upper-test property were
strictly negative on an inner ball, comparison with a small
negative multiple of this exponential, adjusted to vanish on
the outer boundary, would propagate strict negativity across
the annulus.  Comparison follows by taking a positive interior
maximum of the difference, which would supply a forbidden
upper test with negative Laplacian.  Therefore every zero has
a neighborhood of zeros.  The zero set is also closed and
contains $o$; connectedness gives $b_-=-b_+$.

We also need distance comparison in integral form, across cut
points.  Write $d_p=d(p,\cdot)$ and take a nonnegative compact
smooth test $\phi$ whose support is at least distance $A>0$
from $p$.  First, coordinate Lipschitzness and integration by
parts give
\[
 \int d_p\Delta\phi\,d\mu
      =-\int (\mathrm d d_p)(\nabla\phi)\,d\mu.
\]
For this identity, extend each local coordinate representative
of $d_p$ to a Lipschitz function on Euclidean space.
Integration by parts along coordinate lines and Rademacher's
theorem identify its weak derivatives with its differential
almost everywhere.  Apply this to the smooth compact flux
$\sqrt{\det g}\,g^{ij}\partial_j\phi$ and sum over a
partition of unity near the test support.  The resulting
differential pairing is integrable, so this is a signed
integral identity, not a formal differentiation of distance.

Use normalized exponential polar coordinates in a ball
containing the support, keeping only nonterminal minimizing
vectors.  Along a unit direction $\theta$ they form an interval
$0<t<c(\theta)$.  Here minimizing vectors form a measurable
star-shaped set.  A nonzero vector is terminal exactly when
no larger rational multiple is still minimizing within the
exponential domain; this gives a measurable terminal set
meeting each positive ray at most once.  Polar Fubini makes
that set null.  The smooth exponential sends it to a null
set, using local Lipschitz bounds on a countable cover.
On the retained set a minimizing extension and the no-corner
property imply uniqueness of the minimizing segment, hence
injectivity of the exponential map.  Change of variables
therefore applies to signed integrable functions as well as
volumes.  Write the polar Jacobian as
$J(t,\theta)=t^{n-1}H(t,\theta)$.
Jacobi comparison makes $H$ nonincreasing before the terminal
radius.  The raw exponential density is nonnegative and
continuous on every closed radial interval inside its domain,
even where the exponential differential becomes singular.
At a differentiability point of distance, its gradient is the
unit radial velocity: the Lipschitz bound gives norm at most
one, and differentiation of $d_p(\exp_p(t\theta))=t$
gives inner product one.  Thus the right-hand integrand in
Green's identity is the radial derivative of the test.

Put $m=n-1>0$ and choose a finite radius $b$ whose open ball
contains the test support, strictly below the radius of the
exponential domain.  Truncate each ray at
$d=\min(c(\theta),b)$.  For $q(t)=\phi(\exp_p(t\theta))$,
the signed one-dimensional comparison is
\[
 -\int_0^d Jq'\,dt
 \leq\int_0^d\frac{m}{t}Jq\,dt-q(d)J(d)
 \leq\int_0^d\frac{m}{t}Jq\,dt.
\]
Indeed $J'\leq mJ/t$, and the product $qJ$ tends to zero
at zero and to its displayed terminal value at $d$.
Integrating this derivative upper bound proves the inequality;
it requires no separate integrability of $J'$.
The integrands are integrable because $J=O(t^m)$ and
$J/t=O(t^{m-1})$.  Integrate over directions and use
$t\geq A$ wherever $q\ne0$.  We obtain
\[
 \int d_p(x)\Delta\phi(x)\,d\mu
       \leq\frac{n-1}{A}\int\phi\,d\mu.
\]
Now take $p=\gamma(T)$.  Every fixed compact test support
eventually lies at distance at least $A$ from this pole.
Use $\int\Delta\phi=0$, pass to the Busemann limit on the
compact support by domination with
$d(o,x)|\Delta\phi(x)|$, and then let $A\to\infty$.
This gives $\int b_+\Delta\phi\geq0$, and the reversed line
gives the opposite inequality.  Signed tests follow by writing
$\phi=(\phi+C\chi)-C\chi$, with a nonnegative compact smooth
cutoff $\chi=1$ near its support and $C\geq\sup|\phi|$.
Thus $b_+$ is weakly harmonic.  On a small coordinate ball,
bounded chart differential and straight segments turn its
intrinsic Lipschitz bound into a Euclidean Lipschitz bound.
A Lipschitz extension and coordinate-line integration by
parts construct weak first derivatives in $L^\infty$, hence
in $L^2$ on every compact subball.  Put $u=b_+\circ e$ and
$a^{ij}=\sqrt{\det g}\,g^{ij}$ in these coordinates.
Transporting a compact smooth coordinate test $\psi$ to the
manifold, with zero extension outside the chart, gives
\[
 \int u\,\partial_i(a^{ij}\partial_j\psi)=0.
\]
The flux $a^{ij}\partial_j\psi$ is itself a smooth compact
test.  Weak integration by parts and symmetry of $a$ therefore
give $\int a^{ij}\partial_j u\,\partial_i\psi=0$.
The smooth positive definite coefficients are uniformly
elliptic on each compact subball.  Density of smooth compact
tests and continuity of the energy pairing extend the equation
to compactly supported $H^1$ tests on smaller balls.
The local bootstrap in
Section~\ref{sec:rf-weak-harmonic-replacements} gives a smooth
harmonic representative.  It equals $u$ everywhere: both
are continuous and equality almost everywhere excludes a
nonempty open set of disagreement.  Hence $b_+$ itself is
smooth and harmonic.

Completeness also gives, for every $x$ and $r>0$, a point $y$
on the radius-$r$ sphere with $b_+(y)=b_+(x)+r$.
Indeed maximize $b_+$ on the compact closed ball and compare
its maximum with points at distance almost $r$ along almost
shortest paths to $\gamma(T)$.  Let the path error tend to zero
and then $T\to\infty$; the Lipschitz upper bound gives equality
and places the maximizer on the sphere.  For small $r$, express
these calibrated points as $\exp_x(rv_r)$ with $|v_r|=1$.
A convergent subsequence of $v_r$ and differentiability give
$db_+(v)=1$, while the Lipschitz bound gives $|\nabla b_+|\leq1$.
Hence $|\nabla b_+|=1$ everywhere.  Bochner now yields
\[
 0=\tfrac12\Delta|\nabla b_+|^2
   =|\nabla^2b_+|^2+\operatorname{Ric}(\nabla b_+,\nabla b_+),
\]
The identity follows by differentiating the squared gradient
norm twice in a normal frame: its quadratic term is
$|\nabla^2b_+|^2$, and commuting the traced third derivatives
gives $\langle\nabla\Delta b_+,\nabla b_+\rangle+
\operatorname{Ric}(\nabla b_+,\nabla b_+)$.
Both remaining terms in the display are nonnegative, so the
Hessian vanishes.  Set $f=b_+$; then $f(\gamma(0))=0$,
$|\nabla f|=1$, and $\nabla^2f=0$.
Our Busemann sign is opposite to that in
\cite[Definition 2.2 and Proposition 2.3, p.~22, and Lemma 2.14,
pp.~28--29]{MT2007}.
\end{proof}

\subsection{Backward persistence and the fixed factor}

\begin{lemma}[Persistence of a parallel gradient]
\label{lem:ancient-parallel-persistence}
Let $n\geq2$ and let $G(t)$, $t\leq0$, be a Ricci flow on a connected $n$-manifold,
with complete slices, nonnegative curvature operator and a uniform
full-curvature bound.  Suppose a smooth function $f$ satisfies
$|\nabla^{G(0)}f|=1$ and $\nabla^{2,G(0)}f=0$.
Then for every $t\leq0$,
\[
 \nabla^{G(t)}f=\nabla^{G(0)}f,\qquad
 |\nabla^{G(t)}f|=1,\qquad \nabla^{2,G(t)}f=0.
\]
\end{lemma}
\begin{proof}
We preserve the same function and vector field on the fixed
manifold.  On any closed slab
$[a,0]$, let $\sigma_k$ be the sum of the $k$ smallest Ricci
eigenvalues, equivalently the minimum Ricci trace on orthonormal
$k$-frames.  These functions are continuous and nonnegative.
A minimizing frame extended by radial parallel transport gives
a spatial upper support whose scalar Laplacian is the trace of
the rough Ricci Laplacian at the center.  Transporting the frame
in time by $\partial_t e_j=\operatorname{Ric}^{\sharp}e_j$
preserves orthonormality and cancels the fixed-frame Ricci-square
terms.  Its remaining reaction is
\[
 2\sum_{j=1}^k\sum_{\ell}
       \lambda_\ell\,\mathrm{Rm}(e_j,v_\ell,e_j,v_\ell)\geq0,
\]
where $v_\ell$ is a Ricci eigenbasis and $\lambda_\ell\geq0$.
For a smooth spatial lower contact to $\sigma_k$, comparison
with the spatial upper support bounds its Laplacian by the
derivative of this time support.  Lemma~\ref{lem:ancient-heat-contacts}
therefore propagates positivity of every $\sigma_k$ forward
on the connected manifold.  Since
$\sigma_k=0$ exactly when $\dim\ker\operatorname{Ric}\geq k$,
Ricci nullity is spatially constant for $t>a$ and nonincreasing
in time on $[a,0]$.

Constant spatial rank supplies smooth local null sections $W$.
Here is the construction through any prescribed null vector.
In a fixed tangent trivialization, represent Ricci by a smooth
symmetric matrix $A(y)$ using a fixed Euclidean inner product.
At the base point $x$ let $P$ be orthogonal projection onto
$\ker A(x)$.  The matrix $A(x)+P$ is invertible, hence
$A(y)+P$ remains invertible nearby.  Projection $P$ maps
$\ker A(y)$ injectively to $\ker A(x)$: a vector in its
kernel would also lie in the kernel of $A(y)+P$.
Equal dimensions make this projection surjective.
Thus for $c\in\ker A(x)$ the smooth section
$W(y)=(A(y)+P)^{-1}c$ lies in $\ker A(y)$ and has $W(x)=c$.
Transport it back through the trivialization.  No smooth
choice of Ricci eigenvectors is needed.
We also record why a Ricci-null vector $v$ annihilates full
curvature under nonnegative sectional curvature.  The symmetric
form $(u,w)\mapsto\operatorname{Rm}(u,v,w,v)$ is nonnegative
and its orthonormal trace is $\operatorname{Ric}(v,v)=0$,
so the form is zero.  For fixed $y$, the nonnegative form
$(u,w)\mapsto\operatorname{Rm}(u,y,w,y)$ consequently has
$v$ in its kernel.  Thus $\operatorname{Rm}(v,y,w,y)=0$.
Polarizing in $y$ gives
\[
 \operatorname{Rm}(v,y,w,z)+\operatorname{Rm}(v,z,w,y)=0.
\]
Together with pair symmetry and first Bianchi, these identities
give $\operatorname{Rm}(v,y,w,z)=0$ in every slot.
Twice differentiating $\operatorname{Ric}(W,\cdot)=0$ gives
\[
 (\Delta\operatorname{Ric})(W,W)
     =2\sum_j\operatorname{Ric}(\nabla_{e_j}W,\nabla_{e_j}W).
\]
At any time $t>a$, hold $W(x)$ fixed in time and use the earlier
slab $[a,t]$.  Its nonnegative Ricci pairing is zero at the
right endpoint, so its left derivative there is nonpositive.
Nonnegative sectional curvature and Ricci nullity annihilate
every curvature component with a $W$ slot by the preceding
algebra, hence also the
Ricci reaction.  The evolution equation makes the displayed
nonnegative sum nonpositive.  Every $\nabla_{e_j}W$ is null.
Differentiating the null pairing now gives
$(\nabla\operatorname{Ric})(\cdot,W,\cdot)=0$; differentiating
once more gives $(\Delta\operatorname{Ric})(W,\cdot)=0$.
Indeed the first identity holds for every local null section,
so the correction containing $\nabla W$, itself null, also
vanishes.
Thus the time derivative of Ricci annihilates its current kernel.

This last assertion preserves the actual kernel vectors in a
fixed tangent space, including through changes of rank.
To see the necessary linear algebra, represent Ricci there by
a symmetric matrix $A(t)$ using one fixed Euclidean inner
product.  At a reference time let $P$ be orthogonal projection
onto its kernel.  The operator $A+P$ is invertible, since $A$
is invertible on the orthogonal complement and $P$ is identity
on the kernel.  It stays invertible nearby.  On an interval of
constant nullity, for every reference null vector $c$ the
section $w(t)=(A(t)+P)^{-1}c$ lies in $\ker A(t)$: projection
from that kernel onto the reference kernel is injective and,
by equal dimensions, surjective.  Differentiating
$(A+P)w=c$ and using $A'w=0$ gives $w'=0$.
Continuity makes nullity upper semicontinuous; combined with
its time antitonicity, this gives constant nullity immediately
to the left of each time.  Closedness of $\{t:A(t)c=0\}$ and
continuation by these left neighborhoods give
\[
 \ker\operatorname{Ric}_{G(t)}(x)
          \supseteq\ker\operatorname{Ric}_{G(0)}(x),\qquad a\leq t\leq0.
\]
Continuity includes $a$, where no null-derivative assertion
was required.

Apply this to the fixed terminal field $V=\nabla^{G(0)}f$.
Its parallelness implies $\operatorname{Ric}_{G(0)}(V,\cdot)=0$,
so it is Ricci-null at every earlier time.  The null-section
calculation annihilates the two tensor slots of
$\nabla\operatorname{Ric}$ containing $V$.
Differentiate $\operatorname{Rm}(V,\cdot,\cdot,\cdot)=0$:
the correction containing $\nabla V$ vanishes because every
covariant derivative of a local null section remains null.
Thus $\nabla\operatorname{Rm}$ vanishes whenever a curvature
slot is $V$.  Contracted second Bianchi gives
\[
 \sum_i(\nabla_{e_i}\operatorname{Rm})(X,e_i,Y,V)
   =(\nabla_V\operatorname{Ric})(X,Y)
       -(\nabla_Y\operatorname{Ric})(X,V).
\]
The left side and the last term vanish, so the derivative
slot vanishes as well.
The connection variation formula
\[
 G(t)((\partial_t\nabla)_X V,Y)
 =-(\nabla_X\operatorname{Ric})(V,Y)
  -(\nabla_V\operatorname{Ric})(X,Y)
  +(\nabla_Y\operatorname{Ric})(X,V)
\]
is consequently zero.  Thus $\nabla^{G(t)}V=0$ on the whole
slab, with endpoints included by continuity.  Finally
$\partial_tG(t)(V,\cdot)=-2\operatorname{Ric}(V,\cdot)=0$,
so $G(t)(V,\cdot)=df$, $|V|_{G(t)}=1$, and
$\nabla^{2,G(t)}f=0$.  Taking arbitrary $a<0$ proves persistence
throughout the ancient past.
\end{proof}

\begin{lemma}[A fixed ancient factor with positive volume]
\label{lem:ancient-fixed-factor}
Let $n>2$ and let $G(t)$, $t\leq0$, be a Ricci flow on a
connected $n$-manifold $M$, with complete slices, nonnegative curvature operator and
$|\mathrm{Rm}_{G(t)}|\leq K$ for one $K\geq0$.
Suppose $R_G(0,o)>0$ and $w=\operatorname{AVR}(G(0))>0$.
Let $f$ be smooth, with $f(o)=0$, and suppose the same field
$V=\nabla^{G(t)}f$ is parallel and unit for every $t\leq0$.
Then the induced metrics $h(t)$ on $N=f^{-1}(0)$ form a
Ricci flow on a connected noncompact $(n-1)$-manifold, with
complete slices, nonnegative curvature operator,
$|\mathrm{Rm}_h|\leq K$ and $R_h(0,o)>0$.
For every $t\leq0$, $y\in N$ and $r>0$,
\[
 \operatorname{Vol}_{h(t)}B(y,r)\geq\frac{w}{2^{n+1}}r^{n-1}.
\]
In particular this one constant gives parabolic noncollapsing
and positive asymptotic volume at every time.
\end{lemma}
\begin{proof}
Set $m=n-1$.  The smooth regular-level
atlas is fixed using $G(0)$; let $h(t)$ be the induced metric
for every $t\leq0$.  The complete unit parallel field has a
global flow $\Phi_s$.  Along it $f(\Phi_sx)=f(x)+s$, and
$\mathcal L_VG(t)=0$.  Consequently
\[
 (y,s)\longmapsto\Phi_s(y),\qquad
 (N\times\mathbb R,h(t)+ds^2)\longrightarrow(M,G(t))
\]
is a diffeomorphism and an isometry: its inverse sends $x$ to
$(\Phi_{-f(x)}x,f(x))$; tangential derivatives preserve the metric
and are orthogonal to the unit vector $V$.
The factor is connected because the map
$x\mapsto\Phi_{-f(x)}x$ is a continuous surjection from the
connected ambient manifold onto $N$.  It is complete:
inclusion contracts intrinsic distances, so a factor Cauchy
sequence is ambient Cauchy; its ambient limit lies in the
closed level $N$.  The embedded-level topology agrees with
the intrinsic metric topology, giving convergence in $N$.
Regular-level coordinates identify $TN=\ker df=V^\perp$.
The induced connection is the tangential ambient connection,
as follows from the Koszul formula.  Its second fundamental
form satisfies
$\langle\nabla_XY,V\rangle=-\langle Y,\nabla_XV\rangle=0$.
Parallelness also gives $\mathcal R(X,Y)V=0$;
curvature symmetries and Bianchi annihilate every normal slot.
Gauss therefore identifies intrinsic curvature with the
tangential ambient tensor.  Adjoin $V$ to a factor orthonormal
basis: all additional terms in the Ricci, scalar and full
curvature-norm sums vanish.  In particular factor Ricci is
exactly ambient Ricci restricted to tangent vectors.  Hence
$\partial_th=-2\operatorname{Ric}_h$ on this fixed manifold.
The same identities preserve nonnegative curvature operator,
the bound $|\mathrm{Rm}_h|\leq K$, and positive scalar
curvature at the factor basepoint.

Let $w=\operatorname{AVR}(G(0))>0$.  Ancient volume monotonicity
and center comparison give the uniform all-past bound
$\operatorname{Vol}_{G(t)}B(x,r)\geq(w/2^n)r^n$ for all
$t\leq0$, $x$, and $r>0$.
Both projections of the product isometry contract lengths,
so an ambient $r$-ball about $(y,0)$ lies in
$B_{h(t)}(y,r)\times(-r,r)$.
In product coordinates the metric matrix is block diagonal
with last entry one, so its density is exactly the factor
density.  Tonelli's theorem gives product volume on measurable
chart rectangles.  A countable chart cover extends this identity
to the whole cylinder, without any assumption of finite total
volume.
It follows that
\[
 \frac{w}{2^n}r^n
 \leq\operatorname{Vol}_{G(t)}B((y,0),r)
 \leq2r\operatorname{Vol}_{h(t)}B(y,r),\qquad
 \operatorname{Vol}_{h(t)}B(y,r)\geq\kappa' r^m,
 \quad\kappa'=\frac{w}{2^{n+1}}>0.
\]
This single constant works for every past time, center and radius.
It implies parabolic noncollapsing, and Bishop--Gromov gives
positive factor asymptotic volume.  The factor is noncompact:
otherwise its finite total volume would bound these ball
volumes while $\kappa'r^m\to\infty$.
\end{proof}

\begin{remark}[The bounded-flow dependency]
\label{rem:ancient-cone-gap}
The annular construction above develops its compact-domain
coverage, overlap maps and finite patching, and constructs the
connected smooth link.  The backward enclosing realization,
angular-distance formula and volume-rigidity argument are now
developed as well.  The dimension-reduction argument uses the
fixed-factor geometry and product-volume calculation just proved.
The bounded ancient proposition is the precise input needed in the
Ricci-flow chapter's volume-to-curvature argument.
\end{remark}

\begin{lemma}[A compact round backward limit]
\label{lem:ancient-compact-round-limit}
If a prescribed backward rescaling sequence of a three-dimensional
ancient $\kappa$-solution has a compact round limit, then every
slice of the original solution has constant positive sectional
curvature, including time zero.
\end{lemma}
\begin{proof}
Eventually the convergence embeddings are defined on the whole
compact limit.  Their images are open, since they are local
diffeomorphisms, and closed, since they are compact.  The source
is connected, so each image is the whole source.  Thus the original
carrier is compact and convergence controls the whole slice.

Let $k>0$ be the sectional curvature of the limiting slice at $-1$.
For the ordered source curvature eigenvalues $\lambda\geq\mu\geq\nu$,
our dimension-three contractions are
\[
 R=2(\lambda+\mu+\nu),\qquad
 |\mathrm{Rm}|^2=4(\lambda^2+\mu^2+\nu^2).
\]
Uniform convergence of these contractions to $6k$ and $12k^2$
implies
\[
 (\lambda-k)^2+(\mu-k)^2+(\nu-k)^2
       =\tfrac14|\mathrm{Rm}|^2-kR+3k^2\longrightarrow0
\]
uniformly on the whole slice.  For each $c>1$, it follows that
$0<\nu\leq\mu\leq\lambda\leq c\nu$ at all sufficiently
late rescalings.  This ratio is invariant under scaling, so the
original flow has this pinching at arbitrarily early times.

We recall why the pinching persists.  In dimension three the
symmetric form $B=(R/2)g-\operatorname{Ric}$ has these same three
eigenvalues.  The conditions $B\geq0$ and
$B(v,v)\leq c B(w,w)$ for all unit $v,w$ define a closed convex
cone invariant under orthonormal changes of frame.
In a time-dependent orthonormal frame the curvature reaction is
\[
 \lambda'=2(\lambda^2+\mu\nu),\quad
 \mu'=2(\mu^2+\lambda\nu),\quad
 \nu'=2(\nu^2+\lambda\mu).
\]
Order is preserved because differences satisfy linear equations.
At $\nu=0$ its reaction is nonnegative, and on $\lambda=c\nu$
the reaction of $\lambda-c\nu$ is
\[
 2(c-1)\nu\bigl(c\nu-(c+1)\mu\bigr)
                    \leq-2(c-1)\nu^2\leq0.
\]
These inequalities preserve the cone for the reaction equation.
For the full parabolic equation, transport the tensor by the
Ricci equation for orthonormal frames into a fixed metric fiber.
At a spatial maximum of its distance to the cone, a supporting
unit functional has nonpositive pairing with the rough Laplacian:
extend the support by spatial parallel transport and apply the
scalar second derivative test.  The quadratic reaction is
Lipschitz on the bounded tensor range of any compact time slab.
Its pairing at a closest cone point is nonpositive, so the upper
derivative of the maximum distance is at most a constant times
that distance.  An exponential integrating factor, or its strict
barrier approximation, keeps an initially zero distance zero.
This is the compact tensor maximum principle in the fixed cone.
\cite[Theorem 4.8, p. 66; Corollary 9.44, pp. 208--209]{MT2007}

Fix $t<0$ and apply this preservation from an earlier pinched
slice to $t$, for every $c>1$.  Letting $c\downarrow1$ makes
the curvature spectrum scalar.  Equivalently,
$\operatorname{Ric}=(R/3)g$.  Contracted Bianchi gives
$\tfrac13dR=\tfrac12dR$, hence $R$ is spatially constant.
Nonnegative curvature and nonflatness make it positive, and
the dimension-three curvature identity gives sectional curvature
$R/6$.  Continuity of Ricci, scalar curvature and the metric
extends the Einstein identity to zero; the same argument proves
roundness there.
\end{proof}

\begin{theorem}[Universal noncollapsing and zero volume]
\label{thm:ancient-universal}
There is $\kappa_0>0$, independent of the ancient solution and its
original constant, such that every three-dimensional ancient
$\kappa$-solution is either round or parabolically
$\kappa_0$-noncollapsed on all scales.
Separately, every ancient $\kappa$-solution in every dimension
has $\operatorname{AVR}(g(t))=0$ for all $t\leq0$.
The latter conclusion includes round solutions.
\lean{PoincareMT.m22UniversalNoncollapsingAndZeroAVR}
\uses{thm:three-shrinkers,prop:ancient-avr,prop:bounded-ancient-zero-volume}
\source{MT2007}
\dcref{Corollary 9.44, pp. 208--209; Proposition 9.58, pp. 220--221;
Theorem 9.59, pp. 222--225}
\end{theorem}
\begin{proof}
Set $T=2$, $L_0=10(1+e^4)$, and let $v_0$ be one eighth of the
volume of a radius-$1/4$ ball in the normalized round cylinder at $-1$.
Generalized reduced-volume comparison with these fixed
three-dimensional inputs chooses $\kappa_0$ before the solution.

Fix a nonround solution, a point $o$, and a test radius $r$ at zero
satisfying the backward-cylinder curvature bound.
Choose $\tau_i=i+1$ and minimizing rescalings based at $o$.
A compact asymptotic soliton would be round, and
Lemma~\ref{lem:ancient-compact-round-limit} would make the
original solution round.
The asymptotic model is therefore a cylinder or an involution quotient.

Every such model has scalar curvature one at $-1$.  Its
radius-$1/4$ ball has volume at least $8v_0$ at every center.
For the cylinder this follows by rotations and line translations.
For the quotient, use the two involution formulas proved in
Theorem~\ref{thm:ancient-model-identifications}.  Either formula
makes the sphere coordinate antipodal.  Projection from the
time-$-1$ product metric $2g_{S^2}+dz^2$ to the unit sphere in
$\mathbb R^3$ contracts lengths.  Two points in a radius-$1/4$
product ball have sphere chord distance less than $1/2$, whereas
an antipodal pair has chord distance two.  Hence the quotient map
is injective on that ball.  It preserves local volume and
contracts distances, so its image lies in the quotient ball
with the same radius and retains the full product-ball volume.

Take $U_i=B_{g_i(-1)}(p_i,1/2)$ in a sufficiently late rescaling.
The pointed ball-volume comparison with distortion factor two
gives
\[
 \operatorname{Vol}_{h(-1)}B(p_\infty,1/4)
       \leq8\operatorname{Vol}_{g_i(-1)}U_i,
\]
so $\operatorname{Vol}_{g_i(-1)}U_i\geq v_0$.
Volume on this fixed measurable spatial set is nonincreasing
in forward time, since $\partial_t d\mu=-R\,d\mu$ and $R\geq0$.
Thus $\operatorname{Vol}_{g_i(-2)}U_i\geq v_0$.
The convergence maps also cover the source unit ball about $p_i$
by a buffered compact limit ball; scalar convergence and the
model value one give $R_i(-1,\cdot)<2$ on that whole source ball.
Coverage is essential here: estimates only on a chosen image
would not control every point of $U_i$.

The minimizing seed satisfies $\ell_i(p_i,1)\leq3/2$;
we use the weaker bound three to fix the constants.
For $q\in U_i$, choose a constant-speed minimizing spatial
geodesic from $p_i$ to $q$ in $g_i(-1)$, parametrized on
$1\leq s\leq2$ and smoothly extended slightly past its endpoints.
It lies in the buffered unit ball and has squared speed less
than one.  Scalar monotonicity gives $R_i(-s,q(s))\leq2$.
At each of its spatial points, $0\leq\operatorname{Ric}\leq2g$
on $[-2,-1]$; integration of the metric equation gives
$g_i(-s)\leq e^4g_i(-1)$.  Consequently its squared speed
in $g_i(-s)$ is at most $e^4$.

Join this path to a minimizing reduced-length path ending at
$(p_i,1)$.  In square-root time $a=\sqrt s$, the action is
\[
 \int\left(2a^2R_i(-a^2,\beta(a))
                    +\tfrac12|\beta'(a)|_{g_i(-a^2)}^2\right)da.
\]
The first piece has action at most six.  On the second piece
$\beta(a)=q(a^2)$, for $1\leq a\leq\sqrt2$, so its integrand
is at most $8+4e^4$ and its interval has length less than one.
The corner can be smoothed without increasing the limiting
action.  Subdivide into finitely many coordinate intervals
$[a,b]$, where the coordinate curve is a primitive of
$w\in L^2([a,b])$.  Choose smooth $w_j$ supported in $(a,b)$
and tending strongly to $w$.  For a fixed smooth interior bump
$\chi$ with integral one, replace $w_j$ by
\[
 \widetilde w_j=w_j+\chi\int_a^b(w-w_j).
\]
The correction tends to zero in $L^2$ and makes its integral
exactly the prescribed endpoint displacement.  Its primitive
is constant near both endpoints, so neighboring pieces glue
smoothly with unchanged endpoints.  Cauchy--Schwarz bounds
the uniform primitive error by
$\sqrt{b-a}\,\|\widetilde w_j-w\|_2$; hence all sufficiently
late curves stay in a fixed compact subset of the same chart.
The metric coefficients converge uniformly there.  The kinetic
integral error is bounded by a coefficient error times
$\|\widetilde w_j\|_2^2$, plus a fixed constant times
$\|\widetilde w_j-w\|_2(\|\widetilde w_j\|_2+\|w\|_2)$,
and therefore tends to zero.  The continuous scalar potential
converges uniformly as well.  Summing finitely many pieces
gives smooth admissible comparison paths with the joined
action as their limit.  Thus
\[
 2\sqrt2\,\ell_i(q,2)\leq14+4e^4
                         \leq2L_0\leq2\sqrt2 L_0.
\]

For large $i$, $\rho_i=r/\sqrt{\tau_i}$ has $\rho_i^2\leq2$.
The test cylinder rescales to radius $\rho_i$ and bound $\rho_i^{-2}$.
To use Theorem~\ref{thm:m15-noncollapse}, view this ordinary flow
as its product spacetime.  Complete bounded geometry on $[-3,0]$
supplies the ordinary capture and one stable exponential family
at backward time two, whose endpoint image has full measure.
If $\mathcal S$ is its open stable set and $E_2$ its endpoint map,
put $W=\mathcal S\cap E_2^{-1}(U_i)$.  This is open, its
branch action equals the actual reduced length, and
\[
 E_2(W)=U_i\cap E_2(\mathcal S),\qquad
            \operatorname{Vol}_{-2}E_2(W)
                      =\operatorname{Vol}_{-2}U_i\geq v_0.
\]
The product cylinder is exactly the tested terminal ball times
its time interval, and completeness gives compact ball closure.
These verify the stable-set and cylinder hypotheses of that
theorem with the fixed parameters $(3,2,L_0,v_0)$, giving
$\mu_{g_i(0)}(B(o,\rho_i))\geq\kappa_0\rho_i^3$.
Undo scaling to obtain $\kappa_0r^3$.
For the open-past convention, apply this estimate at radii
$0<r'<r$.  Their closed time intervals lie strictly inside
the given open-past interval, and $r^{-2}\leq(r')^{-2}$.
Ball inclusion gives $\operatorname{Vol}B(o,r)\geq\kappa_0(r')^3$;
let $r'\uparrow r$.

Finally translate a terminal $t\leq0$ to zero.  The translated
solution remains nonround.  Otherwise a complete positive round
slice makes the carrier compact by Bonnet--Myers, and its round
past supplies every pinching factor at arbitrarily early times.
The preservation argument above forces roundness of the original
flow through its own time zero, a contradiction.  Thus the same
$\kappa_0$ works at every terminal time.

For zero volume, a compact slice has finite total volume and $n>0$.
For a noncompact slice, nonflatness excludes dimensions zero and one.
The structural theorem bounds its entire past, and ancient
noncollapsing implies the closed-cylinder version in
Proposition~\ref{prop:bounded-ancient-zero-volume}.  Applying that
proposition with the calibrated measure proves the assertion.
\end{proof}

\section{Compactness at normalized scalar curvature}

Fix $\kappa>0$.  Let $(N_i,g_i(t),p_i)$ be any sequence of
three-dimensional ancient $\kappa$-solutions with $R_i(0,p_i)=1$.
The carriers may differ.

\begin{theorem}[Normalized compactness]
\label{thm:ancient-compactness}
The sequence has a subsequence converging smoothly on compact spatial
sets and compact time intervals in $(-\infty,0]$ to an ancient
$\kappa$-solution $(N_\infty,h(t),p_\infty)$ with
$R_h(0,p_\infty)=1$.  One carrier, subsequence, and family of
convergence maps are retained through time zero.
For each $r>0$ there is $C(r)>0$, independent of $i$, such that
\[
 \begin{split}
 R_i(0,x)&\leq C(r) &&(x\in B_{g_i(0)}(p_i,r)),\\
 A_i(t,x)&\leq C(r) &&(t\leq0,\ x\in B_{g_i(0)}(p_i,r)).
 \end{split}
\]
The limit has globally bounded curvature on each slice.
\lean{PoincareMT.m23NormalizedKappaCompactness}
\uses{thm:ancient-universal,thm:surface-ancient}
\source{MT2007}
\dcref{Theorem 9.64, pp. 225 and 229; Lemma 9.65 and Claim 9.66,
pp. 225--227; Kleiner--Lott, corrected 2013, Theorem 46.1,
pp. 2687--2688, and Appendix E, pp. 2848--2849}
\end{theorem}
\begin{proof}[Proof of local curvature control]
First, for each $\nu>0$ there is $D=D(\kappa,\nu)$ such that
$\mu(B(p,r))\geq\nu r^3$ implies $r^2R(x)\leq D$ on that ball.
Here all balls and scalar values in the hypothesis are at time zero,
and the constant works for all carriers, centers and radii.
Suppose instead that $r_i^2R_i(0,x_i)\to\infty$ with
$x_i\in B_0(p_i,r_i)$ and the indicated volume lower bound.
Put $f=\sqrt{R_i(0,\cdot)}$ and $L=r_if(x_i)/2$ for one term.
Starting at $x_i$, replace a point $y$ by $z$ whenever
\[
 d_0(x_i,z)\leq d_0(x_i,y)+L/f(y),\qquad f(z)>2f(y).
\]
The quantity $d_0(x_i,y)+2L/f(y)$ does not increase at a
replacement, since $2L/f(z)<L/f(y)$.  This procedure terminates:
otherwise the values of $f$ would double indefinitely, contradicting
the bounded curvature of this individual terminal slice.
At its final point $q_i$, set $Q_i=R_i(0,q_i)$ and
$s_i=L/\sqrt{Q_i}$.  The displacement budget gives
\[
 d_0(x_i,q_i)+2s_i\leq r_i,\qquad
 0<s_i\leq r_i,\qquad
 s_i\sqrt{Q_i}=\tfrac12r_i\sqrt{R_i(0,x_i)}.
\]
In particular $q_i\in B_0(p_i,2r_i)$ and
$B_0(q_i,s_i)\subset B_0(p_i,2r_i)$.
For $y$ in the former ball, the stopping rule gives
$R_i(0,y)\leq4Q_i$.  Ancient scalar monotonicity and
nonnegative curvature operator therefore give
$|\operatorname{Rm}_i|(t,y)\leq4Q_i$ for every $t\leq0$
on this same terminal ball.

Recentering loses only the factor $27$.  Indeed
$B_0(p_i,r_i)\subset B_0(q_i,3r_i)$; hence that larger ball
has volume ratio at least $\nu/27$.  Bishop--Gromov gives
\[
 \operatorname{Vol}_0 B(q_i,a)\geq(\nu/27)a^3
                     \qquad(0<a\leq3r_i).
\]
Define $\widehat g_i(t)=Q_i g_i(t/Q_i)$ and
$L_i=s_i\sqrt{Q_i}\to\infty$.
These rescalings retain $\kappa$, have base scalar one,
and satisfy $|\operatorname{Rm}|\leq4$ throughout the past
on their terminal $L_i$-balls.  Every such ball of radius
$0<a\leq L_i$ has volume at least $(\nu/27)a^3$.
For each fixed radius the bounds thus hold eventually.
Any finite exceptional prefix can be absorbed into a larger
local constant, since each individual ancient solution has
a whole-past curvature bound.

Interior compactness supplies one complete limit $H(t)$ for
$t<0$.  The eventual bound $B=4$ passes at every point and
time, so this limit is globally bounded and has nonnegative
curvature operator.  To retain positive volume growth at
an interior reference time, set $c=e^{27B}$.
Metric comparison on a terminal ball gives both tangent-norm
comparisons with factor $c$ between times $-1$ and $0$.
Consequently the terminal radius-$a/c$ ball lies in the
earlier radius-$a$ ball, and its terminal volume is at most
$c^3$ times the earlier volume there.  The terminal lower
bound at radius $a/c$ therefore yields
\[
 \operatorname{Vol}_{\widehat g_i(-1)}B(q_i,a)
       \geq\frac{\nu}{27c^6}a^3
\]
for every fixed $a>0$ and all sufficiently large $i$.
The two factors $c^3$ account separately for the radius
change and the volume comparison.  Complete-ball volume
convergence passes this bound to $H(-1)$ at every radius.

There is also a fixed interior time at which the same limit
is nonflat.  Local second-derivative estimates give
$|\nabla^2\operatorname{Rm}_i|(t,q_i)\leq C$ uniformly.
At these basepoints $0\leq R_i(t,q_i)\leq1$, so contraction
and the scalar evolution equation give
\[
 \partial_tR_i(t,q_i)
 =\Delta R_i+2|\operatorname{Ric}_i|^2
 \leq27C+2=:D_0.
\]
Integration up to the included endpoint zero gives
$R_i(-\delta,q_i)\geq1/2$ for
$\delta=\min\{1/2,1/(2D_0)\}$.
Interior convergence at this fixed time proves
$R_H(-\delta,q_\infty)>0$.

For completeness, positive curvature somewhere at this
time forces nonflatness on every earlier slice.  On a
closed intervening slab let $0\leq R\leq B_1$.
Nonnegative Ricci curvature implies
\[
       (\partial_t-\Delta)R=2|\operatorname{Ric}|^2
                         \leq2R^2\leq2B_1R.
\]
If an earlier slice were flat, the initial scalar would
vanish.  On a compact carrier scalar comparison immediately
keeps it zero.  On a complete noncompact carrier, smooth
distance approximation and the bounded curvature derivatives
give a fixed proper function $\psi\geq1$ with
$\Delta_t\psi\leq C_1$ throughout the slab.
This uses one initial distance smoothing: metric comparison
controls its gradient at later times, and the connection
variation, bounded by $\nabla\operatorname{Ric}$, controls
its Hessian.  Subtract
$\varepsilon e^{A(t-a)}\psi$ from $R$, where $a$ is the
initial time and $A>2B_1+C_1$.
The barrier satisfies
$(\partial_t-\Delta)\bigl(\varepsilon e^{A(t-a)}\psi\bigr)
>2B_1\varepsilon e^{A(t-a)}\psi$.
Its properness and the upper scalar bound confine the
positive part of the difference to a fixed compact sublevel.
The scalar maximum principle there, followed by
$\varepsilon\downarrow0$, gives $R\leq0$ on the slab.
This contradicts the positive value at time $-\delta$.

The limit has a positive noncollapsing constant as well.
On a source slice a bound $|\operatorname{Rm}|\leq a^{-2}$
on an $a$-ball bounds scalar there by $3a^{-2}$.
The past comparison then bounds full curvature by the same
quantity on its terminal $a/3$-ball throughout the past.
In particular the required past cylinder satisfies the
weaker threshold $(a/3)^{-2}$.  Source noncollapsing gives
volume at least $(\kappa/27)a^3$ for the original $a$-ball.
Lemma~\ref{lem:rf-limit-static-noncollapse}, with its strict
smaller-ball curvature margin, passes this static estimate
to every interior limit slice.
A parabolic curvature bound includes its terminal static
bound, so the limit is also parabolically
$(\kappa/27)$-noncollapsed.

Translate $H(-\delta)$ to time zero.  Completeness, bounded
curvature, curvature sign and the preceding backward
nonflatness make this an ancient $(\kappa/27)$-solution.
The slice $H(-1)$ lies in its past, where zero asymptotic
volume is already known.  But its all-radius lower bound
makes that asymptotic volume at least $\nu/(27c^6)>0$.
This contradiction proves the finite-radius estimate without
using compactness at the terminal time.

If normalized local control failed at a fixed radius $r$, choose
$x_i\in B_0(p_i,r)$ with $r^2R_i(0,x_i)\to\infty$.
The preceding estimate implies
$\mu_i(B_0(p_i,r))\to0$.
The small-radius density at $p_i$ is $\omega_3$ and ball volume is
continuous in positive radius.  The intermediate value theorem
therefore supplies $0<\rho_i<r$ with
\[
             \mu_i(B_0(p_i,\rho_i))
                        =\tfrac12\omega_3\rho_i^3 .
\]
These radii tend to zero: otherwise their displayed volumes would
stay bounded below while lying in the collapsing radius-$r$ balls.
Rescale by $\rho_i^{-2}$, retaining $p_i$ and the carrier.
The new unit-ball volume is exactly $\omega_3/2$, whereas the new
base scalar curvature is $\rho_i^2\to0$.
The positive unit volume and the finite-radius estimate give uniform
whole-past curvature bounds on each terminal ball.
For radii at most one, Bishop--Gromov supplies the same lower
volume ratio as the unit ball.  For any fixed larger radius $r$,
ball inclusion supplies the positive ratio $\omega_3/(2r^3)$.
The bound on the radius-$r$ ball may depend on $r$; no global
curvature bound on the resulting limit is assumed.

Complete interior compactness gives one limit for all times
$t<0$, with the same basepoint and nonnegative curvature
operator.  At every fixed $t<0$ the source estimates give
\[
 0\leq|\operatorname{Rm}_i|(t,p_i)
               \leq R_i(0,p_i)=\rho_i^2\longrightarrow0.
\]
Hence the limit scalar vanishes at its basepoint for every
negative time.  On a compactly contained time slab ending at
any such time, the scalar equation
$(\partial_t-\Delta)R=2|\operatorname{Ric}|^2\geq0$ and
the strong maximum principle make scalar curvature vanish
throughout that terminal slice.  Nonnegative curvature
operator then makes the entire curvature tensor zero.
Repeating for each negative terminal time shows that the
whole interior limit is flat.  This argument uses local
parabolic regularity and connectedness, not a global limit
curvature bound or an extension to time zero.

The limit inherits the original parabolic constant $\kappa$.
For a tested radius, use smaller radii whose closed cylinders
lie inside the required past interval; their curvature
thresholds have a strict margin.  Complete-ball coverage
places the corresponding source balls in the convergence
images, where curvature and volume converge.  Apply source
noncollapsing and let the smaller radii increase to the
tested radius.  Since the limit is flat at every interior
time, every radius satisfies the parabolic antecedent.
Thus every interior limit ball has volume at least
$\kappa r^3$, without a loss in the constant.

In particular let $B$ bound full curvature on the terminal unit balls.
Metric comparison and ball monotonicity give, for $0<\delta<1$,
\[
 \mu_{-\delta}(B_{-\delta}(p,1))
       \leq e^{81B\delta}\mu_0(B_0(p,1)).
\]
Choose
$\delta=\min\{1/2,\log(3/2)/(81(B+1))\}$.
Then the earlier unit volume is at most $3\omega_3/4$, uniformly in $i$.
Ball-volume convergence at the fixed interior time $-\delta$
therefore gives a complete flat slice with
\[
 \mu(B(p,r))\geq\kappa r^3\quad(r>0),\qquad
 \mu(B(p,1))\leq\tfrac34\omega_3 .
\]
The deficit is measured at an interior time, so terminal
convergence has not been assumed.

Here is the flat-volume contradiction without a classification
of complete flat quotients.  The local injectivity estimate
with zero curvature and unit-ball volume at least $\kappa$
gives a fixed $\sigma=\sigma(\kappa)>0$ on which the
exponential is injective.  Completeness and minimizing
geodesics identify the image of its tangent $\sigma$-ball
with the intrinsic $\sigma$-ball.  In a parallel frame along
each radial geodesic, the Jacobi equation is $J''=0$;
$J(0)=0$ and $J'(0)=v$ give $J(t)=tv$.  Thus the exponential
pullback metric is Euclidean and its volume density is one.
The intrinsic $\sigma$-ball has volume $\omega_3\sigma^3$.
For any prescribed $r>0$, scale the flat metric by
$(\sigma/r)^2$.  Its unit-ball volume is still at least
$\kappa$ by the all-radius lower bound.  The same argument
gives Euclidean volume at radius $\sigma$ in this scaled
metric, which is radius $r$ in the original metric.
Canceling the volume scale proves
$\operatorname{Vol}B(p,r)=\omega_3r^3$ for every $r>0$.
At $r=1$ this contradicts the strict deficit.
The resulting terminal scalar bound gives the whole-past curvature
bound by Theorem~\ref{thm:ancient-structure}.
\end{proof}
\begin{proof}[Proof of convergence through time zero]
Return to the original normalized sequence.  Its local curvature
bounds give a complete interior limit $h(t)$, $t<0$, with
time-independent convergence maps $e_i$ on an exhaustion of
one carrier.  We retain these maps and this sequence.
For $x\in B_{g_i(0)}(p_i,r)$, the earlier unit ball
$B_{g_i(t-1)}(x,1)$ lies in $B_{g_i(0)}(p_i,r+1)$.
Its curvature is therefore bounded on $[t-1,t]$ by a
constant depending only on $r$.  Completeness makes its
closure compact.  The local derivative estimate on that
unit time interval gives, for every $m$,
\[
 |\nabla^m\operatorname{Rm}_i|(t,x)\leq D_m(r)
                         \qquad(t\leq0),
\]
uniformly in $i$, including $t=0$.

Fix a compact coordinate ball in the limit carrier.
Completeness at time $-1$ and interior convergence place its
image under $e_i$ in one bounded source ball at time $-1$,
and hence in a terminal ball of the same radius.
All preceding derivative estimates apply on this image.
Write $B_i(t,x)$ for the bilinear coefficients of
$e_i^*g_i(t)$ in this fixed chart.  At $t=-1$ their spatial
derivatives converge, and compactness gives constants
$\alpha,\beta>0$ with
$\alpha|v|^2\leq B_i(-1,x)(v,v)\leq\beta|v|^2$
for all sufficiently large $i$.
If full curvature is bounded by $C$ on the image and
$d\geq1$, metric comparison on $[-d,0]$ gives
\[
 e^{-54Cd}\alpha|v|^2\leq B_i(t,x)(v,v)
                         \leq e^{54Cd}\beta|v|^2.
\]
Thus both the metric and its inverse are uniformly bounded
on every such closed slab.

The fixed-coordinate derivative argument in the proof of
Lemma~\ref{lem:ancient-positive-ratio-obstruction} applies
with reference time $-1$.  Once spatial derivatives through
order $m-1$ are bounded, the highest derivative satisfies
\[
 \|\partial_tD_x^mB_i\|
                   \leq C_m(1+\|D_x^mB_i\|).
\]
Gronwall propagates its reference bound throughout $[-d,0]$.
The Ricci-flow coordinate equation then bounds all mixed
derivatives using finitely many spatial derivatives and the
uniform inverse bound.  Continuity from negative times
extends these bounds to the closed time boundary.  The
index after which a bound holds may depend on the compact
set and the derivative order; each argument uses only
finitely many such bounds at once.

On the closed half-cylinder, the next derivative bound
makes each derivative equicontinuous.  Arzela--Ascoli and
diagonal extraction give a smooth candidate coefficient
field $B$ up to zero.  On every negative-time compact set
all candidate derivatives agree with those of the already
retained interior metric, by uniqueness of limits.
This auxiliary extraction does not change the convergence
sequence in the theorem.  To see this, let $J_i$ be any
fixed mixed derivative and $J$ its candidate limit on a
compact spatial set.  A next-derivative bound supplies a
common time-Lipschitz constant $L$ on $[-1,0]$ for $J_i$;
interior convergence and continuity give the same bound
for $J$.  For a compact spacetime set $K$, replace
$(t,x)$ by $(\min\{t,-\eta\},x)$, where $0<\eta<1$.
The image $K_\eta$ is compact and has strictly negative
times.  The triangle inequality gives
\[
       \sup_K\|J_i-J\|
       \leq 2L\eta+\sup_{K_\eta}\|J_i-J\|.
\]
Choose $\eta$ first and then use the known interior
convergence on $K_\eta$.  Every derivative therefore
converges through zero along the original retained sequence.

Symmetry and the positive lower quadratic bound pass to
$B$.  It defines a positive smooth metric in the open
coordinate ball.  Spatial differentiation and one-sided
time differentiation of the convergent coefficients pass
the equation $\partial_tB=-2\operatorname{Ric}(B)$ to the
limit: Ricci depends continuously on the positive metric,
its inverse and its first two spatial derivatives.
On overlapping charts the tensors agree for $t<0$ with
the same interior metric.  Continuity of their values on
fixed tangent vectors gives agreement also at $t=0$.
They consequently define a smooth closed ancient flow
on the original carrier, with the original chart maps
and the one-sided Ricci-flow equation at zero.

We next prove terminal completeness using only local
curvature control.  For $x\in B_{h(0)}(p_\infty,A)$,
choose a smooth path from $p_\infty$ to $x$ of length
less than $A$.  Its image is compact even though terminal
completeness has not yet been proved.  Terminal coefficient
convergence bounds the pulled-back source tangent norm
by twice the limit tangent norm on this path for all
sufficiently large $i$.  Thus
$e_i(x)\in B_{g_i(0)}(p_i,2A)$.
Passing source curvature bounds to the limit at each
fixed $t\leq0$ and $x$ gives one bound on
$B_{h(0)}(p_\infty,A)$ throughout the past.
The bound is independent of $x$ and $t$, although the
index used to pass it to the limit need not be.

Let $(x_j)$ be a terminal Cauchy sequence.  It lies in
some terminal $r$-ball.  Almost minimizing terminal paths
between two points of that ball can be chosen inside the
terminal $3r$-ball.  If $K$ is the curvature bound there
on $[-1,0]$, metric comparison along these paths gives
\[
              d_{h(-1)}(x_j,x_k)
                        \leq e^{27K}d_{h(0)}(x_j,x_k).
\]
Hence the sequence is Cauchy in the complete metric
$h(-1)$ and converges to a point of the same carrier.
Both smooth Riemannian metrics induce its manifold
topology, so this is also convergence for $h(0)$.
This proves terminal completeness without using compactness
of terminal balls in the argument that establishes it.

Nonnegative curvature operator passes to zero by continuity
of its quadratic form on each fixed tangent frame.
The scalar time estimate used above gives, at the basepoint,
$1-D_0(-t)\leq R_h(t,p_\infty)\leq1$ for $t<0$.
Continuity at zero yields $R_h(0,p_\infty)=1$.

Finally the original parabolic constant $\kappa$ is retained
through zero.  Interior convergence retains that constant
by the smaller-cylinder argument already used for the flat
limit.  Suppose a terminal $r$-ball satisfies
$|\operatorname{Rm}|\leq r^{-2}$ on its cylinder
$(-r^2,0]\times B_{h(0)}(p,r)$.
Fix $0<\rho<r$ and choose
$\rho^2-r^2<s<0$.  Nonnegative Ricci curvature gives
$B_{h(s)}(p,\rho)\subset B_{h(0)}(p,r)$.
The earlier cylinder ending at $s$ has time interval
$(s-\rho^2,s]\subset(-r^2,0]$, and its curvature threshold
$\rho^{-2}$ is larger.  Interior noncollapsing and local
volume comparison on the fixed terminal ball give
\[
 \kappa\rho^3
 \leq\operatorname{Vol}_{h(s)}B(p,\rho)
 \leq e^{81r^{-2}(-s)}\operatorname{Vol}_{h(0)}B(p,r).
\]
First let $s\uparrow0$, with $\rho$ fixed, then let
$\rho\uparrow r$.  The result is the required
$\kappa r^3$ lower bound.  There is no loss in $\kappa$
at the included endpoint.
\end{proof}
\begin{proof}[Proof of the global bound]
\emph{The auxiliary bounded product.}
At this stage the limit has only local curvature bounds.
Write $g=h(0)$ and $R=R_g$.  Suppose $R$ is unbounded.
Fix $p$.  Completeness and continuity bound $R$ on each
closed ball, so we can choose $x_i$ outside
$\overline B(p,2(i+1)^2)$ with $R(x_i)>\max\{i,16\}$.
Apply the doubling selection of
Lemma~\ref{lem:ancient-selected-line}, with budget
$d(p,x_i)/2$.  Here termination uses compactness of the
closed budget ball: infinitely many doublings of
$\sqrt R$ inside that ball contradict continuity.
We obtain $q_i,\rho_i$ with
\[
 Q_i:=R(q_i)\geq R(x_i),\quad
 d_i:=d(p,q_i)\geq\tfrac12d(p,x_i),\quad
 \rho_i\sqrt{Q_i}=\tfrac14d(p,x_i)\sqrt{R(x_i)},
\]
and $R\leq4Q_i$ on $B(q_i,\rho_i)$.
Put $r_i=\rho_i/(i+1)$, $L_i=r_i\sqrt{Q_i}$ and
$D_i=d_i\sqrt{Q_i}$.  Then
\[
 d_i,Q_i,L_i,D_i\longrightarrow\infty,
       \qquad 0<r_i/d_i\leq1/(i+1)\longrightarrow0.
\]
The unscaled radii $r_i$ need not diverge.

Scalar monotonicity passes from the original sequence
to $h$ at fixed points and times, by smooth convergence
and then continuity at time zero.  Thus the actual
rescaled flows
\[
 H_i(t)=Q_i h(t/Q_i),\qquad t\leq0,
\]
satisfy $R_{H_i}(0,q_i)=1$ and
$|\mathrm{Rm}_{H_i(t)}|\leq4$ throughout
$B_{H_i(0)}(q_i,L_i)$ for all $t\leq0$.
Here the sharp nonnegative-operator inequality is
$|\mathrm{Rm}|\leq R$.
The original completeness and parabolic noncollapse
are preserved by these homotheties.
The curvature bound verifies noncollapse on the
terminal half-unit ball, giving a positive volume bound
at the selected basepoint.  Expanding-cylinder
compactness therefore applies using only the indicated
controlled balls.

Fix for the moment any positive cap $\eta$ on a
backward time shift.  The basepoint Shi estimate and
scalar evolution give a constant $D>0$, independent
of the eventual shift, such that
$1-R_{H_i}(-\delta,q_i)\leq D\delta$ for all large $i$.
Choose $0<\delta<1$ below $\eta$ and $1/(2D)$.
The shifted sequence $H_i(t-\delta)$ has a complete
pointed limit $G(t)$, $t<\delta$, with
\[
 |\mathrm{Rm}_G|\leq4,\qquad
 \operatorname{Rm}_G\geq0,\qquad R_G(0,o)\geq\tfrac12.
\]
The global bound follows because every fixed radius
eventually lies inside the controlled radius $L_i$.
The same radius-three argument used in the finite-radius
estimate gives static noncollapse constant $\kappa/27$:
a terminal curvature bound on an $r$-ball controls its
past on the concentric $r/3$-ball by scalar monotonicity.
This lower volume bound passes to the complete limit.
It implies parabolic noncollapse with that constant,
since the terminal slice belongs to each tested
parabolic cylinder.  The bounded scalar comparison
already proved prevents an earlier flat slice.
Thus $G|_{(-\infty,0]}$ is an actual bounded ancient
$\kappa/27$-solution.

There is a line through $o$ in this same limit.  Use
the farther-vertex and opposite-segment construction
in Lemma~\ref{lem:ancient-geometric-selected-line} after choosing
the convergence subsequence.  That construction uses
only nonnegative sectional curvature of the fixed
metric $g$ and the displayed escape and scale limits.
Here take its margin constant to be $4e^{12\delta}$.
On the controlled balls the curvature bound gives
$0\leq\operatorname{Ric}_{H_i}\leq12H_i$.
For endpoints in the trimmed radius-$s_i$ ball, the
closed-time distance estimate becomes
\[
 d_{H_i(0)}(x,y)\leq d_{H_i(-\delta)}(x,y)
       \leq d_{H_i(0)}(x,y)+108\delta.
\]
This is the earlier additive estimate with dimension
three, cutoff scale one and Ricci bound twelve.
It extends to time zero because metric comparison
on a surrounding controlled ball gives
$d_{H_i(t)}(x,y)\to d_{H_i(0)}(x,y)$ as $t\uparrow0$.
The error is independent of $s_i\to\infty$, so the
earlier comparison cosines still tend to $-1$.
Corresponding-side comparison and pointed convergence
then give the line exactly as in that lemma.

Apply Lemmas~\ref{lem:ancient-line-busemann} and
\ref{lem:ancient-parallel-persistence} to $G$, then use the
fixed-level product and curvature construction in the proof of
Lemma~\ref{lem:ancient-fixed-factor}.  That construction uses
the parallel unit gradient; its later positive-volume conclusion
is not needed here.  It gives one fixed product diffeomorphism
$N\times\mathbb R\to G$ whose pullback metric is
$k(t)+dz^2$ for every $t\leq0$.
Completeness, bounded curvature and nonflatness pass
to $k$.  Its parabolic noncollapse constant can be
taken to be $\kappa/54$: an ambient ball of radius $r$
lies in $B_k(x,r)\times(-r,r)$, whose product volume is
$2r\operatorname{Vol}_k B_k(x,r)$; the factor's tested
curvature bound verifies the ambient cylinder condition.
Theorem~\ref{thm:surface-ancient} now identifies $N$
as a round sphere or a round projective plane.

\emph{Necks in the original terminal slice.}
It remains to obtain necks in $h(0)$, although the
product was obtained at the earlier rescaled time
$-\delta$.  Fix the desired $\varepsilon>0$, and put
$m=\lfloor\varepsilon^{-1}\rfloor$ and
$J=[-\varepsilon^{-1},\varepsilon^{-1}]$.
For either round model let $\Phi:S^2\times\mathbb R\to G$
be its normalized parametrization or cover, centered at
$o$.  With $\theta=R_G(0,o)$ it satisfies
\[
 \theta\Phi^*G(0)=2g_{S^2}+dz^2,
       \qquad \tfrac12\leq\theta\leq1,
       \qquad1-\theta\leq D\delta.
\]
The parametrization includes the constant axial dilation
needed for this scalar normalization.

The finite-jet time estimate supplies a constant $B$
before $\delta$, $G$ and $\Phi$ are selected.  Through
the retained convergence maps $e_i$, all coefficient
jets through order $m$ of $(e_i\Phi)^*H_i$ change by at
most $B\delta$ from $-\delta$ to zero on this slab.
Here is why its uniformity is available.  Unit past
cylinders give curvature derivative bounds of every
required order.  In centered sphere-product charts,
the round reference tensors with
$1/2\leq\theta\leq1$ give fixed initial jet bounds and
ellipticity between $1/2$ and $7$.  The metric-jet
time argument used above, with its highest derivative
entering linearly, then bounds the time derivatives by
a constant depending only on those initial bounds and
the finite order.  The compact image of the model slab
has bounded source distance; its radius affects the
index from which the estimate holds, but not $B$.
This also applies in locally invertible projective
charts.

Let $Z$ bound the model coefficient jets through order
$m$, and let $C_0$ convert coefficient-jet errors to
the sum of squared model-covariant errors.
Both depend only on the fixed model slab and order.
If the earlier coefficient error from the limit is
at most $\tau$, the terminal coefficient error from
the scalar-one cylinder is at most
\[
 \tau+B\delta+|\theta^{-1}-1|Z
       \leq\tau+(B+2DZ)\delta.
\]
Choose the time cap and coefficient tolerance in advance
as
\[
 \eta=\frac{\varepsilon}
              {2(C_0+1)(B+2DZ+1)}.
\]
Construct the buffered limit with $\delta\leq\eta$;
then convergence gives $\tau\leq\eta$ for large $i$.
Since $C_0\leq(C_0+1)^2$, the terminal sum of squared
covariant errors is at most $\varepsilon^2/4$.
If the two model branches give different constants,
choose the smaller of their time caps before extracting
the common limit.  The order of choices therefore
allows either round factor without changing the selected
source sequence.

The maps $e_i\Phi$ now give spherical necks or projective
neck covers in $H_i(0)$.  Compactness of the closed
slab and eventual injectivity of $e_i$ preserve exactly
the original model fibers; the projective fibers are
$(q,z)\sim(-q,z)$.
Undoing the homothety retains these very maps in
$g=h(0)$, with centers $q_i$ and exact scales
$s_i=Q_i^{-1/2}\to0$.
For fixed $\varepsilon$, every half-slab has diameter
from its center at most $C(\varepsilon)s_i$ by the
spherical-path and axial-path comparison.  Since
$d(p,q_i)\to\infty$, it eventually excludes $p$.
These are the shrinking, escaping necks required below,
all in one complete terminal metric.

\emph{The convex exhaustion and its smooth approximations.}
The obstruction now uses nonnegative curvature and a convex exhaustion,
and differs from the strict-positive-curvature flux argument above.
Work entirely in the fixed complete terminal metric $g=h(0)$.
Fix $p$ and let
\[
 f(x)=\sup\bigl(\{0\}\cup\{-b_\gamma(x):\gamma(0)=p\}\bigr),
 \qquad b_\gamma(x)=\lim_{T\to\infty}(d(x,\gamma(T))-T),
\]
where $\gamma$ ranges over minimizing unit-speed rays.
The construction in Section~\ref{sec:rf-harnack} gives a
nonnegative, convex, $1$-Lipschitz exhaustion with $f(p)=0$
and compact totally convex sublevels $C_b=\{f\leq b\}$.
It also gives one constant $R_*>0$ such that
\[
                  d(p,x)\leq R_*f(x)\qquad(f(x)\geq1).
\]
Indeed choose $R_*$ larger than every distance from $p$
in the compact unit sublevel.  If a minimizing segment
from $p$ to $x$ has length $L>R_*$, its point at distance
$R_*$ has exhaustion value greater than one.  Convexity
on the segment bounds that value above by $R_*f(x)/L$.
Thus $L<R_*f(x)$; the case $L\leq R_*$ is immediate.

We first construct the required smooth approximations.
For $c>0$ let $d_c(x)=d(x,\partial C_c)$ on $C_c$.
Then
\[
 d_c(x)\geq c-f(x),\qquad
 d_c(x)=c-f(x)\quad\hbox{if }f(x)>0.
\]
For the lower bound, every point outside $C_c$ violates
one ray inequality $b_\gamma\geq-c$; its distance from
$x$ is at least $c-f(x)$ by the Lipschitz bound.
Distance to the complement equals distance to the
frontier: a minimizing segment from an interior point
to an exterior point crosses the frontier, while each
frontier point is a limit of exterior points.

For the reverse bound, fix a ray $\gamma$ from $p$.
Limits of minimizing segments from $x$ to $\gamma(T)$
give a ray $\alpha$ from $x$ on which
$b_\gamma(\alpha(t))=b_\gamma(x)-t$.
Indeed the minimizing-segment identity, first at a
fixed parameter and then as $T\to\infty$, gives this
equality; properness provides the convergent segments.
At time $c+b_\gamma(x)+e$, for $e>0$, that ray is outside
$C_c$, and hence
$d(x,C_c^c)\leq c+b_\gamma(x)$ on letting $e\downarrow0$.
If $f(x)>0$, values $-b_\gamma(x)$ approach $f(x)$
by the definition of the supremum.  This proves the
equality.  At the zero-level core only the lower bound
is needed.

We need a quantitative distance smoothing on a compact
set $T$ away from a closed nonempty set $S$.
Suppose $d(x,S)\geq r>0$ there.  For any $e,\zeta>0$
there is a globally smooth $\rho$ such that on $T$
\[
 |\rho-d(\cdot,S)|\leq e,\qquad
 |\nabla\rho|\leq1+\zeta,\qquad
 \operatorname{Hess}\rho\leq
           \left(\frac4{3r}+\zeta\right)g.
\]
Here is the construction in the nonnegative-curvature
case used now.  A nearest point of $S$ exists by
properness.  Join it to $x$ by a minimizing segment of
length $d=d(x,S)$ and discard its initial quarter.
The remaining segment has a smooth local inverse
exponential branch at $x$.  Its radial norm plus $d/4$
is a smooth upper support for distance to $S$, equal
to that distance at $x$, with unit gradient there.
The index form tested on the linearly growing parallel
field on the remaining segment bounds its Hessian by
$1/(3d/4)=4/(3d)$ times the metric.  The radial and
mixed Hessian terms vanish.  This also handles cut
points of the original distance.

Use a normalized exponential chart, shrunk so that its
metric is as close to the Euclidean metric as desired
and its Christoffel symbols are uniformly small.
The upper supports just constructed and the formula
$D^2v=\operatorname{Hess}v+dv(\Gamma)$ imply that the
coordinate distance function minus $C|y|^2/2$ is
concave, with $C$ arbitrarily close to the required
Hessian bound.  This implication follows by restriction
to each affine segment and the one-dimensional
touching-upper-support criterion for concavity.
The coordinate distance is Lipschitz, with constant
arbitrarily close to one.  Extend it as a Lipschitz
function to Euclidean space and convolve with a
nonnegative normalized smooth kernel of sufficiently
small support.  On a smaller coordinate ball all
translated arguments remain in the original ball.
Convolution preserves the Lipschitz constant and the
quadratic concavity bound, since both are inequalities
preserved under positive averaging.  Its value error
is at most the Lipschitz constant times the kernel
radius.  Converting its derivatives back to the metric
gives the desired local bounds with arbitrarily small
loss.

For clarity, the gluing preserves this one-sided Hessian
bound as follows.  Choose finitely many nested smooth
cutoffs $0\leq\psi_j,\chi_j\leq1$, with $\chi_j=1$
near the support of $\psi_j$, with $\chi_j$ supported
in the $j$th local domain, and with some $\psi_j=1$
near every point of $T$.  Let $v_j$ be local
approximations with common value error $\xi$ and the
common gradient and Hessian bounds.  Choose a constant
$Q\geq\sup_T d(\cdot,S)$ and form the globally smooth
branches
\[
 b_j=\chi_jv_j+(1-\chi_j)Q+a(1-\psi_j),\qquad a>0.
\]
On $T$ all branches are at least $d(\cdot,S)-\xi$;
some branch is at most $d(\cdot,S)+\xi$.
Where $\psi_j=0$, the $j$th branch is at least
$d(\cdot,S)-\xi+a$.  Where $\psi_j\ne0$, it equals
$v_j+a(1-\psi_j)$ on a neighborhood, so its gradient
and Hessian losses are bounded by $a$ times fixed
derivative bounds of the cutoff.

Let $A_\tau$ be the convolution of absolute value with
an even probability kernel of radius $\tau$.
It is smooth and convex, $|A_\tau'|\leq1$, and
$|z|\leq A_\tau(z)\leq|z|+\tau$; it equals $|z|$
when $|z|\geq\tau$.  The smooth minimum
\[
 m_\tau(x,y)=\tfrac12(x+y-A_\tau(x-y))
\]
is concave, nondecreasing in each argument, and lies
between $\min(x,y)-\tau/2$ and $\min(x,y)$.
Iterate it over the $k+1$ branches.  The resulting
function of their values has nonnegative differential
weights summing to one, negative semidefinite Hessian,
and value error at most $k\tau/2$ from their minimum.
A branch more than $2k\tau$ above the minimum has
zero weight; this follows inductively from exact
equality with the smaller argument for a separated pair.
Choose $a$ first to make all cutoff derivative losses
small, then $\xi,\tau$ so that
$2\xi+2k\tau<a$ and $\xi+k\tau/2\leq e$.
Every branch with $\psi_j=0$ is now inactive on $T$.
The chain rule gives a convex combination of the
remaining branch gradients, and a convex combination
of their Hessians plus a negative semidefinite term.
This proves the global smoothing bounds, after splitting
the allowed loss between local smoothing and gluing.

Apply this result with $S=\partial C_c$ and with the
prescribed compact set enlarged to contain $p$.
Given the desired Hessian loss $H>0$, take
$r=8/(3H)$ and choose
$c>\sup_T f+r$.  The depth lower bound puts $T$ at
distance at least $r$ from $S$.
Smooth with value error $e$ and derivative loss $H/2$,
and put $u=c-\rho$.
Then $\operatorname{Hess}u\geq-Hg$,
$|\nabla u|\leq1+H/2$, $u(p)\leq e$ and
$u\leq f+e$ on $T$.
Where $f>0$, the exact depth identity also gives
$|u-f|\leq e$.  Thus value error and Hessian loss
can be prescribed independently.

\emph{Uniform regular bands and level transport.}
Here is the uniformity needed on regular bands.
Fix $0<a\leq b$ and the open band
$U=\{a/2<f<b+1/2\}$ inside $C_{b+1}$.
Let $R>0$ bound $d(p,x)$ on this compact convex set.
Choose a smooth approximation $u$ with value error
$e\leq\min\{a/8,1/8\}$, with $u(p)\leq e$,
$u\leq f+e$ on $C_{b+1}$ and $|u-f|\leq e$ where
$f>0$, and with $\operatorname{Hess}u\geq-Hg$.
For $x\in U$, these estimates imply $u(x)-u(p)\geq a/4$.
Along a unit-speed minimizing segment from $p$ to $x$
of length $L\leq R$, total convexity retains the Hessian
bound.  Twice integrating the one-dimensional inequality
$(u\circ\gamma)''\geq-H$ gives
\[
 u(x)-u(p)\leq L|\nabla u(x)|+\tfrac12HL^2.
\]
If $HR^2\leq a/4$, it follows that
$|\nabla u|\geq a/(8R)=:l_{a,b}>0$ on $U$.
This lower bound is chosen before either approximation
error is prescribed; a smaller Hessian error preserves it.
Moreover
$\{x\in U:a\leq u(x)\leq b\}
=C_{b+1}\cap\{a\leq u\leq b\}$:
the value-error margins force every point of the latter
set into $U$.  This proves compactness of the band used
for continuation of the gradient flow.

A second constant is uniform even as the upper level
increases.  Set $l=1/(8R_*)$.  On $C_{b+1}$ choose
$e\leq1/8$ and a Hessian error satisfying
$H\leq1$ and $H R_*(b+1)\leq l$.
For $f(x)\geq1$ the linear bound and approximation give
\[
 H d(p,x)\leq2l,\qquad
             2l\,d(p,x)\leq u(x)-u(p),\qquad
             |\nabla u(x)|\leq2.
\]
The gradient upper bound is included in the smooth
approximation, and follows from its Lipschitz bound with
arbitrarily small smoothing error.
For a remote neck center $x_0$ with $f(x_0)\geq2$, take
$b=f(x_0)+1$, $a=1$, and retain the band lower bound
$l_b=l_{1,b}$.  Impose also
$H\leq l_b^2\log(2)/b$.
All these requirements can be met by decreasing the
Hessian error after the constants have been chosen.

The vector field $X=-\nabla u/|\nabla u|^2$ decreases
$u$ at unit speed.  A tangent variation $v$ to a transported
level remains tangent, since $u(\Phi_s(x))=u(x)-s$.
Differentiating its squared norm, the derivative of the
denominator has zero pairing with $v$, and hence
\[
 \frac{d}{ds}|v|^2
 =-\frac{2\operatorname{Hess}u(v,v)}{|\nabla u|^2}
 \leq\frac{2H}{l_b^2}|v|^2.
\]
The compact regular band supplies the flow for the entire
time $u(x_0)-1$.  Its endpoint map $Q$ transports the
upper level to level one with tangential expansion at most
$e^{Hb/l_b^2}\leq2$.  The opposite normalized gradient
flow supplies an inverse.  Thus $Q$ is a diffeomorphism
of levels and carries each connected component onto a
connected component, with no connectedness assumption
on the whole level.

\emph{Axial control on a remote neck.}
We next explain why the relevant upper component is a
small graph.  Write $\Phi:S^2\times(-\varepsilon^{-1},
\varepsilon^{-1})\to M$ for the neck parametrization at
scale $s$.  In the projective case it is a local
diffeomorphism with precisely the fibers
$(q,z)\sim(-q,z)$.  Let $W>0$ be fixed, and choose the
accuracy threshold so that
$\varepsilon\leq\varepsilon_0\leq1/[4(W+1)]$.
Then $W<1/(4\varepsilon)$, so the central $W$-slab
lies strictly inside the half-neck.
For the projective cover, a spherical path followed by an
axial path gives
\[
 d(\Phi(q,z),x_0)\leq
          \bigl(2W+4(\pi+1)\bigr)s\quad (|z|\leq W).
\]
The same coarse bound applies to an ordinary neck.  Require
also $(2W+4(\pi+1))s\leq1/2$.
Since $f$ is one-Lipschitz and $f(x_0)\geq2$, every point
of this closed slab has $f\geq3/2$ and lies in
$U=\{1/2<f<f(x_0)+3/2\}$.  In particular the same
radial and Hessian inequalities hold throughout the slab.
The constants $l,W$ and the accuracy threshold are fixed
before imposing this smallness condition on $s$.
The fixed point $p$
lies outside the half-neck for the selected escaping
sequence.  A minimizing unit-speed segment from a point
$x$ in the central slab to $p$ must exit that half-neck.
Its first exit has length at least $s/(8\varepsilon)$,
since its axial displacement is at least
$1/(4\varepsilon)$ and the metric comparison bounds
axial speed by $2/s$.
Here is the long-neck comparison needed at this step.
It is the scale-dependent form of
\cite[Lemma A.4, pp. 498--499]{MT2007}.
The model metric is $g_{\mathrm{cyl}}=2g_{S^2}+dz^2$.
In either case let $a$ be the axial coordinate on the
image of $\Phi$.  In the projective case it is well
defined because both points of each fiber have the same
$z$, and it is smooth by the local inverse theorem.
Thus $a\circ\Phi(q,z)=z$.

First the metric jets give, with a universal $K$,
\[
 |\operatorname{Hess}a(v,v)|\leq K\varepsilon/s^2
                                      \qquad(|v|_g=1).
\]
To see the powers of $s$ in this estimate, center a
spherical coordinate chart at the chosen preimage and
use its product with the axial coordinate.  In these
stereographic coordinates $y\in\mathbb R^2$ the model is
$32(|y|^2+4)^{-2}|dy|^2+dz^2$.
Its first derivatives, and hence its Christoffel symbols,
vanish at $y=0$.  The first model-covariant derivative
of the metric error is therefore its ordinary derivative
there.  Each centered coordinate basis vector has model
norm at most two, so the tensor bound gives at most
$8\varepsilon$ for each of its three-slot components.
Expanding each slot in the three coordinate basis vectors
then bounds the full derivative norm by
$3^3\cdot8\varepsilon=216\varepsilon$.
The normalized first-jet error therefore bounds the
derivative of the actual metric coefficients there by
$d_0s^2\varepsilon$, whereas the zeroth-order comparison
gives $B\geq a_0s^2 I$ for fixed $a_0>0$.
The constants remain uniform under the fixed linear
identification of the product coordinate space with
$\mathbb R^3$: the derivative has three covariant slots,
so the change multiplies its norm by at most the cube
of that linear map's norm.  At the chart center the
coordinate expression of $a$ is a fixed linear
functional $\ell$.  Koszul's formula gives
\[
 \|\Gamma\|\leq\frac{3}{2a_0s^2}\|DB\|,
 \qquad |w|^2\leq\frac1{a_0s^2}
                         \quad\hbox{if }B(w,w)=1.
\]
Since $\operatorname{Hess}a(v,v)=-\ell(\Gamma(w,w))$,
one may take $K=3\|\ell\|d_0/(2a_0^2)$.
Only a local chart is used here, so the same estimate
holds at every preimage in a projective cover.

Let $\gamma:[0,L_0]\to M$ be the unit-speed segment
up to the first half-neck exit, and put $F=a\circ\gamma$.
The geodesic equation gives
$|F''|\leq\delta:=K\varepsilon/s^2$.
For every $0\leq b\leq c\leq L_0$, minimality and a
competitor inside the full neck give
\[
 c-b\leq s\sqrt{1+\varepsilon}
          \bigl(|F(c)-F(b)|+C\bigr),
       \qquad C=\sqrt2(\pi+1).
\]
Indeed choose preimages of the two endpoints, join
their spherical coordinates by a path of unit-sphere
length less than $\pi+1$, and interpolate their axial
coordinates linearly.  The path stays in the slab and
has model length at most $\sqrt2(\pi+1)+|F(c)-F(b)|$.
The metric comparison gives the displayed upper bound
for its image length.  Global minimality of the segment
is a lower bound for the length of every such competitor.
This construction also works for the projective cover;
it requires no choice of a global inverse.

Set
\[
 w=\frac{s}{200\sqrt\varepsilon},\qquad
 A=s\sqrt{1+\varepsilon}.
\]
Since $L_0\geq s/(8\varepsilon)>s/(100\varepsilon)$
and $\varepsilon<1/2$, we have $2w\leq L_0$.
Every $t\in[0,L_0]$ consequently belongs to a subinterval
$[b,c]$ of length $w$ lying in $[0,L_0]$: use the
forward interval if it fits, and the backward one otherwise.
On that interval the acceleration estimate implies
$|F'(r)|\leq|F'(t)|+\delta w$.
Integration and the competitor inequality yield
\[
 w\leq A\bigl((|F'(t)|+\delta w)w+C\bigr),
 \qquad
 |F'(t)|\geq A^{-1}-C/w-\delta w.
\]
After multiplication by $s$, the lower bound becomes
\[
 (1+\varepsilon)^{-1/2}
       -(200C+K/200)\sqrt\varepsilon,
\]
which tends to one as $\varepsilon\downarrow0$,
independently of the scale and segment length.
In particular $F'$ never vanishes for small enough
$\varepsilon$.  Its sign is constant by continuity,
and the mean value theorem identifies that sign with
$F(L_0)-F(0)$.  Choose $\sigma\in\{1,-1\}$ so that
$\sigma F(L_0)>\sigma F(0)$.  Given $\beta>0$, decreasing
the accuracy threshold gives
$s\sigma F'(t)\geq1-\beta$ throughout the closed
interval, including both endpoints.

For a preimage of $\gamma(t)$ write
$E=s^{-1}d\Phi(0,1)$ and $v=\gamma'(t)$.
The zeroth-order tensor error and the lower metric
bound give
\[
 |E|_g^2\leq1+\varepsilon,\qquad
 |\langle v,E\rangle_g-s\,da(v)|\leq2\varepsilon.
\]
For the second estimate, lift $v$ through the local
differential of $\Phi$.  Its model norm is at most
$2/s$, and the mixed tensor error against $(0,1)$
is at most $\varepsilon$ times the two model norms.
It follows that
\[
 |\sigma v-E|_g^2
 \leq1+(1+\varepsilon)-2(1-\beta-2\varepsilon)
 =2\beta+5\varepsilon.
\]
Taking $\beta$ and then $\varepsilon$ small proves
$|\sigma v-E|_g<\alpha$ for any prescribed $\alpha>0$.
This is a uniform estimate for the actual axial tangent
of the parametrization, including the projective case.

The segment stays in the totally convex region.
If its length is $L=d(x,p)$, the radial gap and Hessian
bound imply
\[
 L\,du_x(v)\leq u(p)-u(x)+\tfrac12HL^2\leq-lL.
\]
Taking $\alpha=l/4$ and using $|\nabla u|\leq2$
therefore gives
\[
       |\partial_z(u\circ\Phi)(q,z)|\geq(l/2)s
                                      \qquad(|z|\leq W).
\]
The neck-accuracy threshold here depends on $l$ and $W$,
which have been fixed independently of the remote level.

\emph{The graph and its diameter.}
Choose $W=(4\pi+6)/l$ in the spherical case and
$W=(8(\pi+1)+6)/l$ in the projective case.
Take value error at most $\min\{s,1/8\}$.
The central sphere diameter bounds and the Lipschitz
bound on $f$ give, respectively,
\[
 |u(\Phi(q,0))-u(x_0)|\leq(2\pi+2)s
       \quad\hbox{or}\quad(4(\pi+1)+2)s.
\]
Both are strictly less than $(l/2)sW$.
For each $q$ the axial derivative has constant sign on
$[-W,W]$, by continuity and its nonvanishing.
The endpoint values therefore straddle $u(x_0)$.
The intermediate value theorem and strict monotonicity
give a unique height $z=h(q)\in(-W,W)$ with
$u(\Phi(q,h(q)))=u(x_0)$.
The inverse function theorem makes $h$ smooth; uniqueness
makes the local height functions agree.  On a projective
cover, $\Phi(-q,z)=\Phi(q,z)$ implies $h(-q)=h(q)$,
so the graph descends to a projective cross-section.
The graph is the entire level inside the open central
slab.  The image $V$ of that slab is open: every point
has a local inverse for $\Phi$, including in the projective
case.  If $A=\{x\in U:u(x)=u(x_0)\}$, the graph image
is exactly $A\cap V$.  It is therefore open in $A$.
It is also compact, as the continuous image of $S^2$,
and hence closed in the Hausdorff space $A$.
Its connectedness places it inside one component of $A$;
being both open and closed prevents that component from
containing anything else.  This proves capture of a whole
component without assuming that the sphere parametrization
is injective.

Put $P(q)=\Phi(q,h(q))$ and $m=l/2$.
Differentiating its constant level value yields
\[
 dh(v)=-\frac{du(d\Phi(v,0))}{du(d\Phi(0,1))}.
\]
Metric comparison gives $|d\Phi(v,0)|\leq2s|v|_{S^2}$,
so $|dh(v)|\leq(4/m)|v|_{S^2}$ and
\[
              |dP(v)|\leq2s(1+8/l)|v|_{S^2}.
\]
After applying $Q$ the bound doubles.  Any two points
of the corresponding lower-level component have preimages
on $S^2$.  Indeed the whole upper level $A$ is compact,
being a closed subset of the compact regular band.
The restriction of $Q$ is a continuous bijection from $A$
onto the Hausdorff lower level, so its inverse is continuous.
This homeomorphism takes the captured upper component onto
the entire lower component.  Choose spherical preimages of
the two upper points; they can be
joined by a unit-sphere geodesic of length at most $\pi$.
Its image under $Q\circ P$ therefore has length at most
$4\pi s(1+8/l)$.  This argument needs only surjectivity
of the sphere parametrization, so it applies unchanged
to the projective component.

\emph{The fixed lower-level obstruction.}
There is a positive lower bound $d$ for the diameters
of all these level-one components, independent of $u$
and of the remote upper level.  Indeed approximation
error at most $1/8$ places their entire level in the
fixed compact set
$K=C_2\cap\{3/4\leq f\leq5/4\}$, contained in
the fixed open regular band $O=\{1/2<f<3/2\}$.
The regular-level atlas makes each component a compact
connected surface.  To construct its collar, integrate
$\nabla u/|\nabla u|^2$ on a neighborhood of the compact
level.  A finite cover supplies a common positive flow
time, and compactness allows that time to be decreased
so that the orbits stay in the regular neighborhood.
Along each orbit $u(\Phi_t(x))=u(x)+t$.
The map $(x,t)\mapsto\Phi_t(x)$ from the level times
a small interval consequently has inverse
\[
 y\longmapsto\bigl(\Phi_{1-u(y)}(y),u(y)-1\bigr).
\]
The inverse identities follow from uniqueness for the
same smooth vector field.  Both maps are continuous,
and their common image is an open scalar band in a
neighborhood of the level.  Restricting to a connected
component gives the required two-sided collar in $O$.

Cover $K$ by finitely many coordinate balls whose larger
closed coordinate balls still lie in $O$, and take a
Lebesgue number $d>0$.  If a component had diameter less
than $d$, it would fit in one of these coordinate balls.
A compact connected collared surface $Y$ in
$\mathbb R^3$ separates it into two regions with $Y$
as common boundary.  The following construction gives
the precise separation fact used here.
Write the collar as $E:Y\times(-r,r)\to U$.
Clamp $t/r$ to $[-1/2,1/2]$ and take its class in
$\mathbb R/\mathbb Z$.  Near either end of the collar
this is the same class $1/2$, so it extends continuously
by that constant outside $U$.  Compactness of $Y$
makes the part where the map is nonconstant compactly
contained in $U$.
Since $\mathbb R^3$ is simply connected, the map has
a continuous real lift $F$, normalized to vanish at
one point of $Y$.  On the connected collar the lift
must equal the clamped real coordinate, by uniqueness
of lifts.  Outside the collar its values belong to
$1/2+\mathbb Z$, so cannot be zero.  Consequently
$F^{-1}(0)=Y$.

The sets $A=\{F<0\}$ and $B=\{F>0\}$ are disjoint
nonempty open sets covering the complement, and their
intersections with the collar are its connected
negative and positive halves.  Continuity gives
$\partial A,\partial B\subset Y$; approaching from
either half gives the reverse inclusions.
Every component of $A$ meets the collar.  Otherwise
its closure avoids the collar, hence avoids
$\partial A$; it is then both open and closed in
$\mathbb R^3$, an impossibility.  Since $A\cap U$
is connected, it lies in one component and meets all
of them, so $A$ is connected.  The same holds for $B$.
Finally the exterior of a large ball containing $Y$
is connected and must lie entirely in one of $A,B$.
The other side is bounded, and both cannot be bounded
because their union with the compact $Y$ is all of
$\mathbb R^3$.

The bounded region
lies in the coordinate ball: the exterior of that ball
is connected and unbounded, so belongs to the other region.
The coordinate expression of $u$ has constant value on
the boundary of the bounded region.  Its compact closure
has a maximum and a minimum.  Either one is attained
in the interior, or both have the boundary value and
the function is constant throughout.  In either case
there is an interior critical point, contradicting
$du\ne0$ on $O$.
The finite cover is independent of the smoothing, so
the same $d$ works for every component under consideration.
Thus
\[
                       d\leq4\pi s(1+8/l),
\]
which contradicts $s\to0$ for the selected neck sequence.
This proves the obstruction component by component;
neither a one-endedness conclusion nor a globally
connected exhaustion boundary is needed.

The terminal scalar is consequently globally bounded.
Harnack bounds full curvature on the entire past.
Terminal nonflatness follows from $R_h(0,p_\infty)=1$.
If an earlier slice were flat, bounded-curvature scalar evolution
rigidity would force the terminal slice to be flat as well.
Thus the limit satisfies every ancient-solution hypothesis.
Scalar monotonicity alone would not justify this last backward
nonflatness assertion.
\end{proof}

\begin{remark}[Compactness inputs]
\label{rem:ancient-compactness-gaps}
The point-selection estimate, terminal convergence, selected blowup,
convex smoothing and componentwise obstruction have been developed
above.  Their shared inputs include complete pointed compactness,
corresponding-side comparison, radial inverse and index-form estimates,
and the coordinate metric-jet evolution, all used before the
terminal global curvature bound is established.
\end{remark}

\section{The actual quotient models}

The normalized cylinder has time-dependent metric
\[
                  b_t=(-2t)g_{S^2}+dz^2,\qquad t<0,
\]
where $g_{S^2}$ has Gaussian curvature one.  The line metric is fixed.
A model identification means one fixed diffeomorphism transporting
the metrics at every negative time.

\begin{theorem}[Refined model identifications]
\label{thm:ancient-model-identifications}
Every supplied model in Theorem~\ref{thm:three-shrinkers} has the
following refinement on its actual carrier and with its actual maps.
In the compact case a finite group acts freely through $SO(4)$ on
the unit three-sphere, and a fixed quotient map $\pi:S^3\to N$ satisfies
\[
                      \pi^*G(t)=(-4t)g_{S^3}\quad(t<0).
\]
In the product case the supplied sphere-line model is retained.
In the involution case, in the supplied sphere and line coordinates,
the involution is exactly one of
\[
                 J(x,z)=(-x,z),\qquad J(x,z)=(-x,2c-z)
\]
for one fixed real $c$.
The first quotient has a homeomorphism to
$\mathbb{RP}^2\times\mathbb R$ commuting with the original quotient
projection.  The second has a homeomorphism to punctured
$\mathbb{RP}^3$ and the compatible standard smooth double cover.
The quotient projection is a local diffeomorphism transporting
the time-dependent metric and sectional curvature.
No orientability hypothesis is imposed.
\lean{PoincareMT.m24ModelCertificates}
\uses{thm:three-shrinkers}
\source{MT2007}
\dcref{Theorem 1.11 and its space-form consequence, pp. 8--9;
Theorem 9.42, pp. 206--208; Corollary 9.54, p. 218;
Proposition 9.58, pp. 220--221}
\end{theorem}
\begin{proof}
On a compact round slice at $-1$, the soliton equation says
$\operatorname{Hess}f=(1/2-2K)h$.
At the maximum and minimum of $f$ the Hessian has opposite signs,
so $K=1/4$ and $\operatorname{Ric}_h=h/2$.
For each $t<0$, naturality under the supplied slice homothety
gives $\operatorname{Ric}_{G(t)}=G(t)/(2(-t))$.
For fixed tangent vectors the Ricci-flow equation is consequently
$\partial_tG(t)(u,v)=G(t)(u,v)/t$.
The derivative of $G(t)(u,v)/t$ vanishes; its value at $-1$
gives $G(t)=(-t)h$ on the original carrier.  This argument does
not differentiate a choice of slice homotheties.

The metric $h/4$ has sectional curvature one.  The round covering
construction used above, now on the simply connected three-sphere,
gives a surjective smooth local isometry
$\pi:(S^3,g_{S^3})\to(N,h/4)$.  Hence
$\pi^*h=4g_{S^3}$ and $\pi^*G(t)=(-4t)g_{S^3}$ with this one
fixed map.  Define its deck group as the ambient orthogonal
motions $L$ for which $\pi\circ L=\pi$.
For two points in a fiber, the tangent isometry
$(d\pi_y)^{-1}d\pi_x$, together with $x\mapsto y$, extends to
an orthogonal motion of $\mathbb R^4$.  First-order uniqueness
of local isometries shows that this motion preserves $\pi$ on
the connected sphere.  Thus the deck orbits are exactly its fibers.
Uniqueness of lifts shows that a deck motion fixing one point
is the identity.  Evaluation at a point therefore injects the
deck group into a fiber.  A fiber is closed in the compact sphere
and discrete because $\pi$ is a covering; it is finite, as is
the group.

Its matrices have determinant one.  Indeed, for a nonidentity
deck matrix $L$, freeness gives $\det(L-I)\ne0$.  Orthogonality
gives $L(I-L^T)=L-I$, while in even dimension
$\det(I-L^T)=\det(L-I)$.  Taking determinants and cancelling
this nonzero number gives $\det L=1$.  The identity has the same
determinant.  This constructs the claimed finite free $SO(4)$
action without an orientability assumption on $N$.

For an involution quotient write $b_t=(-t)a+b$, where
$a=2g_{S^2}$ and $b=dz^2$.  Isometry at times $-1$ and $-2$
implies separately $J^*a=a$ and $J^*b=b$.
Their kernels are the factor distributions.  Positive definiteness
on each factor shows that the derivative of the sphere component
in the line direction is zero, and conversely for the line
component in every sphere direction.  A smooth map with zero
derivative is constant on a connected slice.  Since the sphere
and line are connected, this gives
$J(x,z)=(A(x),j(z))$ with no dependence on the other factor.
The equality $J^2=\mathrm{id}$ makes both factors smooth
involutions.

The separated metric identities give $j'(z)^2=1$ and make $A$
preserve the round metric.  The derivative bound makes $j$
distance nonincreasing; applying the same inequality to $j(x),j(y)$
and using involutivity gives equality of distances.
For $A$, lengths of curves are preserved in both directions,
so intrinsic distances on the unit sphere are preserved.
The formula
\[
 d_{S^2}(x,y)=\arccos\langle x,y\rangle,
 \qquad |x-y|^2=2-2\langle x,y\rangle
\]
then gives preservation of chord distances and inner products.

For completeness the line classification uses only distances to
$0$ and $1$.  Put $d=j(0)$.  Since $|j(1)-d|=1$, its sign is
either $+1$ or $-1$.  Subtracting the squared-distance equations
to these two points gives respectively $j(z)=d+z$ or $j(z)=d-z$.
In the first case involutivity forces $d=0$; in the second put
$c=d/2$.  Thus $j$ is the identity or $z\mapsto2c-z$.
Each case has a fixed line coordinate, so freeness of $J$ makes
$A$ fixed-point free.

We now prove that $A$ is antipodal using its actual distance
properties.  If $A(x)\ne-x$, put
\[
             m=\frac{x+A(x)}{|x+A(x)|}\in S^2.
\]
Inner-product preservation and $A^2=\mathrm{id}$ give
$\langle A(m),x\rangle=\langle m,A(x)\rangle$ and
$\langle A(m),A(x)\rangle=\langle m,x\rangle$.
Adding and dividing by $|x+A(x)|$ yields
$\langle A(m),m\rangle=\langle m,m\rangle=1$.
Both are unit vectors, so $|A(m)-m|^2=0$, contradicting freeness.
Therefore $A(x)=-x$ for every $x$, proving both literal formulas
in the supplied coordinates.

The topology of the given quotient must also be retained.
Write $p:S^2\times\mathbb R\to Q$ for its supplied projection.
The metric pullback identity at time $-1$ makes $dp$ injective:
$dp(v)=0$ implies $b_{-1}(v,v)=0$, hence $v=0$.
Equal dimensions make $dp$ invertible, so the inverse function
theorem makes $p$ a local diffeomorphism.  It is surjective and
open, and hence is a quotient map.
For the first formula $k(x,z)=([x],z)$ is also an open quotient
map.  The antipodal sphere projection is open because the
saturation of an open set $U$ is $U\cup(-U)$; its product with
the identity remains open and surjective.  The fibers of $k$
are exactly the fibers of $p$.  Define $E(p(x,z))=k(x,z)$.
Equality of fibers makes $E$ a well-defined bijection, and the
quotient properties of both maps make $E$ and $E^{-1}$ continuous.
This supplies the asserted homeomorphism on $Q$.

For the second formula put $w=z-c$ and use the equivariant
diffeomorphism
\[
 H(x,w)=\frac{(w,x)}{\sqrt{1+w^2}}
              \in S^3\setminus\{(1,0),(-1,0)\}.
\]
Its inverse at $(a,u)$, where $u\ne0$, is
$(u/|u|,a/|u|)$.  Both maps are smooth since their denominators
are nonzero.  Simultaneous negation of $(x,w)$ becomes the
antipodal map on $S^3$, and the excluded poles represent one
point of $\mathbb{RP}^3$.  Composing $H$ with the restricted
antipodal projection gives an open quotient map with precisely
the prescribed involution fibers.  The same equality-of-fibers
argument yields the homeomorphism from $Q$ onto the complement
of that projective point.

The smooth covering statement concerns this original $Q$.
Compose $H^{-1}$, the inverse center shift, and $p$ to obtain
\[
 S^3\setminus\{(1,0),(-1,0)\}\longrightarrow Q.
\]
It is surjective, has exactly antipodal pairs as fibers, and is
a local diffeomorphism as a composition of local diffeomorphisms.
The homeomorphism just constructed intertwines it with the
standard punctured projective projection.  These statements
give the compatible smooth double cover without assuming a
smooth structure on a topological quotient in advance.

Finally the metric identities make the fixed projection a local
isometry at every negative time.  Uniqueness of the metric-compatible
torsion-free connection transports the Levi-Civita connection,
and then its curvature tensor.  The numerator and Gram denominator
in sectional curvature are both preserved, giving sectional-curvature
transport at each time.
Composition with the fixed identification of the original flow
gives the simultaneous time-dependent transports claimed.
\end{proof}

Recognizing a high-curvature neighborhood in a surgery flow requires
the additional construction of its blowup sequence and transfer
through the retained convergence maps.  That is the singular-limit
argument.  Ancient classification alone does not identify such a
point or supply the quantitative neck and cap parameters for surgery.
